\section{Barendregt--Geuvers: proof checking and dependent types}
\label{sec:p12-har-types}

This presentation follows the author-hosted 2001 chapter, whose published
extent is pp.~1149--1238. Its \(\lambda\mathrm{HOL}\), pure type systems,
extensions, and historical assistants have their own standing; they are not
identified with the current Rocq or Lean kernels. The following proof-checking
development follows the chapter's proof-checking systems, historical
theory development and dated assistant comparison, retaining their
source-specific notation and assumptions
\cite[pp.~1151--1230]{BarendregtGeuvers2001}.

\subsection{What is checked}

\begin{hbdefinition}{def:p12-har-checking}{Formal theory, proof checker, and proof assistant}
A formalized theory supplies primitive notions, well-formed definitions,
axioms, and proofs in a declared logic. Proof checking verifies the
well-formedness of its definitions and the correctness of its supplied
proofs. In the chapter's type-theoretic approach, an interactive
proof-development system constructs a term, and a separate checker verifies
its type. Their combination is a proof assistant. Proof generation and proof
verification remain different tasks even when automation is available in the
development system.

The formal assertion is \(\Gamma\vdash_T p:A\): the context, proof term,
and proposition represented as a type are explicit linguistic objects.
If a logic \(L\) and a type theory \(T\) have a justified
propositions-as-types translation \([\ ]\), its intended correspondence is
\[
 \Gamma\vdash_L A\text{ using }p
 \quad\longleftrightarrow\quad
 [\Gamma]\vdash_T[p]:[A].
\]
The biconditional is a calculus-relative adequacy obligation, not a
consequence of writing brackets. Definitions and the implementation of the
checker are part of the trust boundary. A set-theoretic encoding of
computation may be possible without being feasible as a checking mechanism;
type-theoretic reduction is intended to make this mechanism explicit
\cite[pp.~1151--1154]{BarendregtGeuvers2001}.
\end{hbdefinition}

\begin{hbdefinition}{def:p12-har-three-problems}{TCP, TSP, and TIP}
For a fixed type theory and well-formed context, the chapter distinguishes
\[
\begin{array}{lll}
 \Gamma\vdash_T M:A?&\mathrm{TCP}&\text{type checking},\\
 \Gamma\vdash_T M:?&\mathrm{TSP}&\text{type synthesis},\\
 \Gamma\vdash_T ?:A&\mathrm{TIP}&\text{type inhabitation}.
\end{array}
\]
TCP asks whether the alleged proof proves the stated proposition; TSP asks
whether the given term has a type; TIP asks for a proof of a given
proposition. Decidability of TCP and TSP depends on the exact system,
conversion relation, and annotations carried by terms. It does not follow
that TIP is decidable. A faithful effective embedding of a logic with
undecidable theoremhood makes inhabitation undecidable for its image.

The chapter's de Bruijn criterion asks for explicit proof objects checkable
by an easy algorithm. A small inspectable checker can provide a strong
reliability argument, but not absolute assurance against implementation or
hardware failure. Large proof objects and a small checker are compatible;
the checker follows local steps rather than discovering the proof
\cite[pp.~1154--1156]{BarendregtGeuvers2001}.
\end{hbdefinition}

\begin{hbdefinition}{def:p12-har-propositions-types}{The constructive reading of the source type formers}
For the direct interpretation, the chapter reads an implication proof as a
method taking a proof of its antecedent to a proof of its consequent, a
conjunction proof as a pair, and a disjunction proof as a tagged alternative.
There is no proof of falsehood. A universal proof is a method with a proof
for every argument; an existential proof supplies a witness and its proof.
Using the chapter's logical symbols and type-former symbols,
\[
\begin{array}{rcl@{\qquad}rcl}
 [A\supset B]&=&[A]\to[B]&[A\mathbin{\&}B]&=&[A]\times[B],\\
 {}[A\lor B]&=&[A]\cup[B]&[\bot]&=&\varnothing,\\
 {}[\forall x\in X.A_x]&=&\Pi x:X.[A_x]&
 [\exists x\in X.A_x]&=&\Sigma x:X.[A_x].
\end{array}
\]
Here \(\cup\) is the source's disjoint-union type, not untagged set union;
its inhabitants are \(\langle0,p\rangle\) with \(p:P\), or
\(\langle1,q\rangle\) with \(q:Q\). The empty-type symbol denotes
the absence of constructors. Negation is \(A\supset\bot\).
The dependent product has functions whose value at \(a:A\) has type
\(P a\), while the dependent sum has pairs \(\langle a,p\rangle\)
with \(p:P a\). These are the constructive intuitions; whether existential
propositions actually permit extraction into data depends on the chosen
sorts and elimination rules \cite[pp.~1156--1157]{BarendregtGeuvers2001}.
\end{hbdefinition}

\subsection{A proof term and its entire application chain}

\begin{hbdefinition}{def:p12-har-relation-terms}{The source relation operators}
The chapter uses \(\mathrm{Set}\) for domains and
\(\mathrm{Prop}\) for propositions. With these local sort roles, define
\[
\begin{aligned}
 \mathrm{Rel}&:=\lambda X:\mathrm{Set}.X\to X\to\mathrm{Prop},\\
 \mathrm{AntiSym}&:=\lambda X:\mathrm{Set}.\lambda R:\mathrm{Rel}\,X.
       \Pi x,y:X.(Rxy)\to(Ryx)\to\bot,\\
 \mathrm{Irrefl}&:=\lambda X:\mathrm{Set}.\lambda R:\mathrm{Rel}\,X.
       \Pi x:X.(Rxx)\to\bot.
\end{aligned}
\]
The identifier \(\mathrm{AntiSym}\) is retained from the source, but its
formula is \emph{asymmetry}, not ordinary antisymmetry. Ordinary
antisymmetry allows \(Rxx\), for example when \(R\) is equality, and
does not imply irreflexivity. Unfolding a definition is \(\delta\)-reduction;
applying an abstraction is \(\beta\)-reduction. Thus
\(\mathrm{Rel}\,X\to_\delta
(\lambda X:\mathrm{Set}.X\to X\to\mathrm{Prop})X
\to_\beta X\to X\to\mathrm{Prop}\).
The product in the type of \(\mathrm{AntiSym}\) depends on its domain:
\[
 \mathrm{AntiSym}:\Pi X:\mathrm{Set}.
       (X\to X\to\mathrm{Prop})\to\mathrm{Prop}.
\]
The abstract product/application surface used in this development is
\hyperref[sheet:p12-har-types]{the source-rule sheet}
\cite[pp.~1157--1158]{BarendregtGeuvers2001}.
\end{hbdefinition}

\begin{hbtheorem}{thm:p12-har-asymmetry}{proved=annex:p12.har.asymmetry}{Checking the source asymmetry proof}
In the source's direct dependent-product setting, the term
\[
 \lambda X:\mathrm{Set}.\lambda R:\mathrm{Rel}\,X.
 \lambda h:\mathrm{AntiSym}\,X R.\lambda x:X.\lambda q:Rxx.hxxqq
\]
has type
\[
 \Pi X:\mathrm{Set}.\Pi R:\mathrm{Rel}\,X.
       (\mathrm{AntiSym}\,X R)\to(\mathrm{Irrefl}\,X R).
\]
Every application uses the same \(x\) and the same proof \(q\); no
equality premise or excluded-middle principle is used.
\end{hbtheorem}

\begin{hbmethod}{def:p12-har-synthesis}{Synthesis followed by comparison}
For an annotated abstraction \(\lambda x:X.M\), synthesize the body type
\(B\) under the added declaration \(x:X\), then return \(\Pi x:X.B\).
For \(F N\), synthesize \(A\) for \(N\) and \(C\) for \(F\);
expose a product \(\Pi x:D.B\) by permitted conversion of \(C\),
compare \(D\) with \(A\), and return \(B[N/x]\) only if they are
convertible. The product-formation and context checks must also succeed.
An application-domain failure is not repaired by comparing \(D\) with
the whole function type \(C\).

To solve TCP by this TSP procedure and a final comparison requires
uniqueness of types up to the same conversion relation, sound and complete
synthesis, and a decidable conversion test. Computing a normal form is a
possible conversion method under the hypotheses below, not always a good
implementation strategy. The source emphasizes reduction only as needed,
with definitions unfolded sparingly: definitions can conceal most relevant
\(\beta\)-redexes even when proof terms appear \(\beta\)-normal
\cite[pp.~1158--1160]{BarendregtGeuvers2001}.
\end{hbmethod}

\begin{hbtheorem}{thm:p12-har-normal-form-test}{proved=annex:p12.har.normal-form-test}{The conditional normal-form conversion test}
Let a finitely branching effective reduction on well-typed terms have
decidable normal-form recognition, preserve well-typedness, be confluent,
and give every well-typed term a normal form. Suppose alpha-equivalence of
normal forms is decidable. Breadth-first reduction search computes a normal
form, and two well-typed terms are convertible exactly when their computed
normal forms are alpha-equivalent. Hence conversion is decidable in this
regime. This statement proves neither normalization nor confluence of an
arbitrary dependent type theory or PTS specification.
\end{hbtheorem}

\subsection{Three framework signatures, not three direct interpretations}

\begin{hbdefinition}{def:p12-har-framework}{The chapter's logical-framework encoding}
In a direct interpretation, logical connectives and their rules are mirrored
by type formers and type-theoretic rules. In the chapter's other approach,
called its \emph{logical-framework} or \emph{shallow} encoding, a signature
declares the object logic inside a host type theory. This published use of
``shallow'' is retained locally; it must not be silently identified with an
alternative modern deep/shallow taxonomy.

Object proposition names have type \(\mathrm{prop}\); they are not
themselves host types to be inhabited. A family
\(T:\mathrm{prop}\to\mathrm{type}\) assigns each name its proof type.
Binding and substitution use the host's binding discipline, while logical
rules become declared constants. Adequacy and faithfulness must be checked
for each signature; a framework's flexibility is not a universal encoding
theorem \cite[pp.~1160--1163]{BarendregtGeuvers2001}.
\end{hbdefinition}

\begin{hbdefinition}{def:p12-har-prop-signature}{The minimal implication signature}
For the \(\{\supset\}\)-fragment of minimal propositional logic, the
source signature \(\Sigma_{\mathrm{PROP}}\) is
\[
\begin{aligned}
 \mathrm{prop}&:\mathrm{type},&
 \mathrm{imp}&:\mathrm{prop}\to\mathrm{prop}\to\mathrm{prop},\\
 T&:\mathrm{prop}\to\mathrm{type},\\
 \mathrm{imp\_intr}&:\Pi p,q:\mathrm{prop}.(Tp\to Tq)\to T(\mathrm{imp}\,p q),\\
 \mathrm{imp\_el}&:\Pi p,q:\mathrm{prop}.
             T(\mathrm{imp}\,p q)\to Tp\to Tq.
\end{aligned}
\]
Declare each atomic proposition symbol as a variable of type
\(\mathrm{prop}\), and set
\([A\supset B]=\mathrm{imp}[A][B]\). For instance,
\[
 \mathrm{imp\_intr}[A][A](\lambda x:T[A].x)
       :T(\mathrm{imp}[A][A]).
\]
The type of an object proposition's proof is \(T[A]\), not
\([A]\) itself \cite[p.~1161]{BarendregtGeuvers2001}.
\end{hbdefinition}

\begin{hbdefinition}{def:p12-har-pred-signature}{The minimal predicate signature}
For the one-sorted \(\{\supset,\forall\}\)-fragment, extend the
propositional declarations (including \(T\)) with
\[
\begin{aligned}
 A&:\mathrm{type},&c&:A,&f&:A\to A,&R&:A\to A\to\mathrm{prop},\\
 \mathrm{forall}&:(A\to\mathrm{prop})\to\mathrm{prop},\\
 \mathrm{forall\_intr}&:\Pi P:A\to\mathrm{prop}.
             (\Pi x:A.T(Px))\to T(\mathrm{forall}\,P),\\
 \mathrm{forall\_elim}&:\Pi P:A\to\mathrm{prop}.
             T(\mathrm{forall}\,P)\to\Pi x:A.T(Px).
\end{aligned}
\]
The source's constant \(c\), unary function \(f\), and binary
predicate \(R\) constitute the chosen object signature. Translation
preserves them and sets
\([\forall x:A.Px]=\mathrm{forall}(\lambda x:A.[Px])\).
Object free variables are declared with type \(A\).

For the complete instantiation of the source example, let
\(U=\mathrm{forall}(\lambda x:A.
\mathrm{forall}(\lambda y:A.Rxy))\). Then
\[
\begin{aligned}
 &\mathrm{forall\_intr}\,(\lambda z:A.\mathrm{imp}\,U(Rzz))\\
 &\quad(\lambda z:A.\mathrm{imp\_intr}\,U(Rzz)
   (\lambda h:T U.
     \mathrm{forall\_elim}\,(\lambda y:A.Rzy)
      (\mathrm{forall\_elim}\,(\lambda x:A.
          \mathrm{forall}(\lambda y:A.Rxy))\,h\,z)\,z))
\end{aligned}
\]
has type \(T(\mathrm{forall}(\lambda z:A.
\mathrm{imp}\,U(Rzz)))\). Thus it represents
\(\forall z:A.((\forall x,y:A.Rxy)\supset Rzz)\), with every
implicit product argument in the chapter's abbreviated display now given
\cite[p.~1162]{BarendregtGeuvers2001}.
\end{hbdefinition}

\begin{hbdefinition}{def:p12-har-lambda-signature}{The framework signature for untyped beta-equality}
The source's third signature \(\Sigma_{\mathrm{lambda}}\) declares
\[
\begin{aligned}
 D&:\mathrm{type},&\mathrm{app}&:D\to D\to D,&
 \mathrm{abs}&:(D\to D)\to D,\\
 \mathrm{eq}&:D\to D\to\mathrm{type},\\
 \mathrm{refl}&:\Pi x:D.\mathrm{eq}\,x x,\\
 \mathrm{sym}&:\Pi x,y:D.\mathrm{eq}\,x y\to\mathrm{eq}\,y x,\\
 \mathrm{trans}&:\Pi x,y,z:D.
       \mathrm{eq}\,x y\to\mathrm{eq}\,y z\to\mathrm{eq}\,x z,\\
 \mathrm{mon}&:\Pi x,x',z,z':D.
       \mathrm{eq}\,x x'\to\mathrm{eq}\,z z'
           \to\mathrm{eq}(\mathrm{app}\,z x)(\mathrm{app}\,z'x'),\\
 \mathrm{xi}&:\Pi F,G:D\to D.
       (\Pi x:D.\mathrm{eq}(Fx)(Gx))
           \to\mathrm{eq}(\mathrm{abs}\,F)(\mathrm{abs}\,G),\\
 \mathrm{beta}&:\Pi F:D\to D.\Pi x:D.
       \mathrm{eq}(\mathrm{app}(\mathrm{abs}\,F)x)(Fx).
\end{aligned}
\]
Its translation is \([x]=x\),
\([P Q]=\mathrm{app}[P][Q]\), and
\([\lambda x.P]=\mathrm{abs}(\lambda x:D.[P])\). Free variables are
declared of type \(D\). Object \(\beta\)-equality is represented by
\(\mathrm{eq}[P][Q]\); it is not host definitional equality.
For example, \(\mathrm{beta}(\lambda x:D.x)y\) checks as
\(\mathrm{eq}(\mathrm{app}(\mathrm{abs}(\lambda x:D.x))y)y\).
No object \(\eta\)-axiom is included
\cite[p.~1163]{BarendregtGeuvers2001}.
\end{hbdefinition}

\subsection{The three framework converses with their exact host}

\begin{hbdefinition}{def:p12-har-fw-host}{The beta-only LF presentation used by the three signatures}
For the framework signatures in
\cref{def:p12-har-prop-signature,def:p12-har-pred-signature,def:p12-har-lambda-signature},
use weak dependent LF with beta conversion, opaque signature constants,
and object-variable contexts. Its split syntax has kinds
\(K::=\mathrm{type}\mid\Pi x:B.K\), families
\(B::=a\mid\Pi x:B.C\mid\lambda x:B.C\mid BM\), and
objects \(M::=c\mid x\mid\lambda x:B.M\mid MN\).
Here \(a\) is a family constant and \(c\) an object constant.
In the PTS view the distinguished kind \(\mathrm{type}\) is
\(\ast\), the top classifier is \(\Box\), the axiom is
\(\ast:\Box\), and the product triples are exactly
\((\ast,\ast,\ast)\), \((\ast,\Box,\Box)\).
The seven underlying PTS rules are the already registered
\hyperref[sheet:p12-har-pts]{Figure~3 surface}, at this specification.
Treat signature declarations as fixed opaque entries when using that view;
do not bind or substitute them as additional object variables.

Thus \(\mathrm{prop}\), \(A\), \(D\) are fixed family constants,
and \(T:\mathrm{prop}\to\mathrm{type}\) and
\(\mathrm{eq}:D\to D\to\mathrm{type}\) are families, not
proposition-valued predicates inside the direct lambdaHOL system.
There is no kind-domain product pair, equality reflection, inductive
eliminator or additional proof-producing axiom. This is the beta-only
LF interface discussed by the primary framework account, not its
separately discussed eta extension and not the three-sort lambdaHOL
\cite[pp.~1160--1163]{BarendregtGeuvers2001}
\cite[sections~2.1 and 2.4]{HarperHonsellPlotkin1993LF}.
\end{hbdefinition}




\begin{hbtheorem}{thm:p12-har-framework-forward}{proved=annex:p12.har.framework-forward}{The three forward encoding constructions}
At the beta-only LF host of \cref{def:p12-har-fw-host},
every natural-deduction derivation in minimal implication logic translates
to a term of \(T[A]\) over \(\Sigma_{\mathrm{PROP}}\), with
its hypotheses translated to variables of their proof types. The same holds
for the displayed one-sorted minimal predicate logic over
\(\Sigma_{\mathrm{PRED}}\). Every derivation of untyped
\(\beta\)-equality translates to a term of
\(\mathrm{eq}[P][Q]\) over \(\Sigma_{\mathrm{lambda}}\).
The translation respects capture-avoiding object substitution by host
substitution.
\end{hbtheorem}

\begin{proofboundary}{def:p12-har-framework-inverse}{Faithfulness is a separate direction}
The converse for each of these exact signatures and admissible open contexts
is proved at \cref{thm:p12-har-fw-faithfulness}, using typed normalization,
normal-head inversion and source reification. Induction on an object
derivation alone proves only the forward direction. The converse does not
extend merely by writing the same constants in a host with additional
axioms, equality reflection or eliminators
\cite[pp.~1161--1163]{BarendregtGeuvers2001}.
\end{proofboundary}

\begin{hblemma}{thm:p12-har-fw-normalization}{proved=annex:p12.har.fw-normalization}{The exact LF host has typed beta-normal representatives}
Every derivable judgment over a well-formed signature/context in this
host has a beta-normal subject of the same type. The identity sort/rule
inclusion in exact CC gives strong normalization; the source rule
construction gives LF subject reduction. Normal-form generation exposes
the head constant/variable and its complete typed application spine.
Opaque atomic families \(Tq\), \(\mathrm{eq}\,uv\),
\(\mathrm{prop}\), \(A\), and \(D\) cannot become dependent
products by beta conversion. This is not an assumption that every
normal term is beta-convertible to an eta-long canonical term.
\end{hblemma}

\begin{hbdefinition}{def:p12-har-fw-contexts}{Admissible contexts and pointwise reification}
For PROP, the context after \(\Sigma_{\mathrm{PROP}}\) consists
of finitely many distinct atom names \(a_i:\mathrm{prop}\) and
proof variables \(h_j:T(q_j)\). For PRED, after the full signature
(including inherited \(T\)), it consists of object variables
\(x_i:A\) and proof variables \(h_j:T(q_j)\).
The \(q_j\) are well-typed formula names, normalized when decoding.
For the untyped-lambda converse, the context after
\(\Sigma_{\mathrm{lambda}}\) consists only of free object
variables \(x_i:D\); an extra equality-proof assumption would assert
an extra equation, not pure beta equality.

Normalize before reification. PROP formula names are atoms or
\(\mathrm{imp}\,p q\). PRED object names are variables, \(c\),
or \(ft\); formula names are \(Rst\), \(\mathrm{imp}\,p q\),
or \(\mathrm{forall}\,F\). Decode the first two formula cases
structurally. For the third, decode
\(\mathrm{forall}\,F\) as
\(\forall z:A.\mathcal D(\mathrm{nf}_\beta(Fz))\), with \(z\)
fresh. A normal \(F:A\to\mathrm{prop}\) is either
\(\lambda z:A.q\) or \(Rs\); the latter case directly yields
\(\forall z:A.R(\mathcal D(s),z)\).

For normal untyped-lambda names, decode \(x:D\) to the source
variable, \(\mathrm{app}\,uv\) to the source application, and
\(\mathrm{abs}\,F\) to
\(\lambda z.\mathcal D(\mathrm{nf}_\beta(Fz))\).
A normal \(F:D\to D\) is either \(\lambda z:D.u\) or
\(\mathrm{app}\,u\); the latter directly yields
\(\lambda z.\mathcal D(u)z\). These pointwise clauses handle
eta-short LF functions without adding an object-language eta equation.
For example \(\mathrm{abs}(\mathrm{app}\,u)\) decodes to
\(\lambda z.\mathcal D(u)z\), not to \(\mathcal D(u)\).
Freshness is relative to the entire current context and the arguments.
\end{hbdefinition}

\begin{hblemma}{thm:p12-har-fw-reification}{proved=annex:p12.har.fw-reification}{The reifiers are defined and commute with open substitution}
In the stated contexts, the normal-name grammars above are exhaustive,
independent of proof variables, and their reifiers terminate. Reification
of the source encoding recovers each source expression up to bound-name
renaming. Reification is invariant under LF beta conversion and commutes
with substitutions of object names for object variables, including
capture-avoiding substitution under \(\mathrm{forall}\) and
\(\mathrm{abs}\). In PRED these are arbitrary first-order term
substitutions, not assignments only to closed constants; in the lambda
case they are arbitrary finite open untyped-term substitutions.
\end{hblemma}

\begin{hbtheorem}{thm:p12-har-fw-faithfulness}{proved=annex:p12.har.fw-faithfulness}{Derivability reflection for all three source framework signatures}
At exactly the host/signature/context discipline above:
\begin{enumerate}
 \item If \(\Sigma_{\mathrm{PROP}},\vec a,\vec h
 \vdash p:T([\varphi])\), then the decoded assumptions
 \(\mathcal D(q_j)\) derive \(\varphi\) in minimal implication
 natural deduction. With no \(\vec h\), \(\varphi\) is a theorem.
 \item If \(\Sigma_{\mathrm{PRED}},\vec x,\vec h
 \vdash p:T([\varphi])\), then those decoded assumptions derive
 \(\varphi\) in the specified one-sorted minimal
 \(\{\supset,\forall\}\)-logic with \(c,f,R\).
 The constructed universal-introduction rule obeys its eigenvariable
 condition, and elimination uses the actual decoded open argument.
 \item If \(\Sigma_{\mathrm{lambda}},\vec x:D
 \vdash p:\mathrm{eq}[P][Q]\), then \(P=_\beta Q\)
 in the untyped lambda calculus, with its actual open free variables.
\end{enumerate}
The proof constructs a source derivation from every normal proof head,
including function-valued proof arguments to introduction/xi constants.
Combined with \cref{thm:p12-har-framework-forward}, these give both
existence directions. They do not assert a bijection on all LF terms or
faithfulness after adding arbitrary signature axioms
\cite[pp.~1161--1163]{BarendregtGeuvers2001}.
\end{hbtheorem}

\begin{hbcorollary}{thm:p12-har-fw-no-eta}{proved=annex:p12.har.fw-no-eta}{The lambda signature has not acquired object eta}
In the pure lambda signature with \(u:D\), no term inhabits
\(\mathrm{eq}\,u\,[\lambda x.ux]\), for \(x\) fresh.
Pointwise reification of LF function arguments is not the object
equation \(u=\lambda x.ux\).
\end{hbcorollary}




\subsection{Algorithms, graphs, and computation inside proof checking}

\begin{hbdefinition}{def:p12-har-functions}{Algorithms and functional relations}
A term \(f:A\to B\) is an algorithm in the source's operational reading.
It may be axiomatically declared, constructed by abstraction, or obtained by
primitive recursion. An axiomatically declared \(f\) is an unknown
algorithm, not an executable definition. For \(b:B\) and
\(f:\mathrm{nat}\to B\to B\), the recursor has type
\(\mathrm{Rec}\,b f:\mathrm{nat}\to B\) and the computation
clauses on the source-rule sheet.

A relation \(R:A\to B\to\mathrm{Prop}\) with a proof of
\(\forall x:A.\exists!y:B.Rxy\) is a function \emph{graph}.
This assertion alone does not name the image of \(a\), supply an
algorithm, or make its existential elimination into data permissible.
In Church's selected description-operator regime, for
\(P:A\to\mathrm{Prop}\) the term \(\iota P:A\) has the axiom
\(\Pi P:A\to\mathrm{Prop}.(\exists!x:A.Px)\to P(\iota P)\).
It chooses an element satisfying a uniquely inhabited predicate;
otherwise it denotes an unspecified element. Such a
constant can create closed natural-number terms that are not constructor
numerals. It is not the recursor's \(\iota\)-reduction despite the shared
Greek letter in the two uses.

The constructive existence property for a closed proof of
\(\forall x:A.\exists y:B.Rxy\) is a metatheoretic claim about the
chosen theory. It need not be an internal extraction operation. Interpreting
the existential as a dependent sum internalizes a particular extraction
operation as follows \cite[pp.~1164--1165]{BarendregtGeuvers2001}.
\end{hbdefinition}

\begin{hbtheorem}{thm:p12-har-sigma-extraction}{proved=annex:p12.har.sigma-extraction}{The dependent-sum extraction pair}
Suppose dependent-sum projections into data and proofs are allowed and
\(p:\Pi x:A.\Sigma y:B.Rxy\). Then
\[
 f:=\lambda x:A.\pi_1(px):A\to B,\qquad
 q:=\lambda x:A.\pi_2(px):\Pi x:A.Rx(fx).
\]
If \(p\) reduces to \(\lambda x:A.\langle t,u\rangle\), the
extracted function reduces to \(\lambda x:A.t\). The assertion does not
permit a \(\mathrm{Prop}\)-existential to be eliminated into data in a
system that forbids that elimination.
\end{hbtheorem}

\begin{hbdefinition}{def:p12-har-subject-reduction}{Subject reduction and a proof detour}
For a specified reduction relation, subject reduction says that
\(\Gamma\vdash M:A\) and \(M\twoheadrightarrow N\) imply
\(\Gamma\vdash N:A\). For data this preserves type under evaluation;
for proof terms it preserves the proved proposition under detour removal.
The source's implication example is
\[
 (\lambda h:A\to A\to B.\lambda z:A.hzz)q
       \to_\beta\lambda z:A.qzz
 \qquad(q:A\to A\to B).
\]
Both terms have type \(A\to B\): applying \(q\) to \(z\)
gives \(A\to B\), applying again gives \(B\), and abstracting
gives the displayed type. This local verification is not a general
subject-reduction proof for arbitrary PTS or inductive extensions.

In developments, a proved term may be named by a definition and made opaque
to prevent automatic unfolding. The source's \(\mathrm{lemma1}\) names
\(\lambda h:A\to A\to B.\lambda z:A.hzz\); using the name
does not require repeatedly reducing its proof. A reduction detour need not
literally shorten the written term \cite[p.~1166]{BarendregtGeuvers2001}.
\end{hbdefinition}

\begin{hbexample}{def:p12-har-plus-conversion}{The source addition and conversion calculation}
Write \(\overline n\) for the numeral \((S^+)^n0\), not a
metalevel natural number. The chapter defines
\(\mathrm{plus}:=\lambda x,y:\mathrm{nat}.
\mathrm{Rec}\,x(\lambda z:\mathrm{nat}.S^+)y\). Thus
\[
\begin{aligned}
 \mathrm{plus}\,\overline1\,\overline0
 &\to_\delta
 (\lambda x,y:\mathrm{nat}.
   \mathrm{Rec}\,x(\lambda z:\mathrm{nat}.S^+)y)
       \overline1\,\overline0\\
 &\to_\beta(\lambda y:\mathrm{nat}.
   \mathrm{Rec}\,\overline1(\lambda z:\mathrm{nat}.S^+)y)\overline0\\
 &\to_\beta\mathrm{Rec}\,\overline1
             (\lambda z:\mathrm{nat}.S^+)\overline0
 \to_\iota\overline1.
\end{aligned}
\]
Given
\(p:\Pi x,y,z:\mathrm{nat}.(x>\mathrm{plus}\,y z)\to(x>z)\)
and \(q:\overline2>\overline1\), first obtain
\(p\overline2\overline1\overline0:
(\overline2>\mathrm{plus}\,\overline1\overline0)
\to(\overline2>\overline0)\). Conversion changes this type to
\((\overline2>\overline1)\to(\overline2>\overline0)\), leaving
the term unchanged. Application to \(q\) gives
\(p\overline2\overline1\overline0q:\overline2>\overline0\).
No arithmetic-order theorem is being inferred from reduction alone: the
given \(p\) supplies the implication and \(q\) its premise
\cite[pp.~1166--1167]{BarendregtGeuvers2001}.
\end{hbexample}

\begin{hbdefinition}{def:p12-har-conversion}{Conversion and the source Poincare principle}
The chapter uses the conversion relation \(=_{\beta\iota\delta}\)
generated by beta reduction, admitted constructor/recursor computation, and
definition unfolding. Its conversion rule changes a type to a convertible
type without changing the proof term. Decidability of this relation is an
explicit premise of the implicit comparison mechanism. A richer equality
that is not decidable needs additional evidence or another checking design;
this is not licensed simply by adding a side condition to the same algorithm.

The source's Poincare principle separates proof from computation:
computationally equal expressions require no new user proof of the
computation. The conversion checker still executes the computation; the
principle does not trust an external answer or remove the need for a sound
internal decision procedure \cite[pp.~1167--1168]{BarendregtGeuvers2001}.
\end{hbdefinition}

\begin{hbtheorem}{thm:p12-har-reflection}{proved=annex:p12.har.reflection}{The source two-level reflection term}
Assume a type \(\mathrm{Class\!\mbox{-}\!of\!\mbox{-}\!Form}\),
terms \(\mathrm{Dec},\mathrm{Check}:
\mathrm{Class\!\mbox{-}\!of\!\mbox{-}\!Form}\to\mathrm{Prop}\),
and a checked soundness term
\[
 \mathrm{ok}:\Pi t:\mathrm{Class\!\mbox{-}\!of\!\mbox{-}\!Form}.
                    (\mathrm{Check}\,t)\to(\mathrm{Dec}\,t).
\]
If \(\mathrm{Dec}\,t=_{\beta\iota\delta}T\),
\(\mathrm{Check}\,t=_{\beta\iota\delta}\top\), and
\(\mathrm{true}:\top\), then
\(\mathrm{ok}\,t\,\mathrm{true}:T\).
The construction uses the soundness implication, not a completeness claim
that every true formula makes the checker return \(\top\).
\end{hbtheorem}

\subsection{Definitional equality, book equality, and setoid constructions}

\begin{hbdefinition}{def:p12-har-equalities}{Two standings of equality}
Definitional equality is the underlying equality used by the type system:
it can be a judgment \(\Gamma\vdash M=N:A\), or a conversion side
condition on pseudoterms. Book equality is a proposition/type
\(M=_A N\), with a proof term sought inside the theory. An external
statement that two functions compute the same values does not make them
definitionally equal in an intensional system.

The source's baseline conversion omits \(\eta\). For annotated raw
terms, \(\lambda x:A.(\lambda y:B.y)x\) can beta-reduce to
\(\lambda x:A.x\), or eta-reduce to \(\lambda y:B.y\); their
annotations need not have a common reduct unless the domains do. Typed
beta/eta metatheory therefore requires its own hypotheses, not the untyped
confluence argument. The example is a warning about annotated raw terms,
not a well-typed counterexample to a correctly restricted system.

An extensional reflection rule that promotes a proof of pointwise book
equality to definitional function equality changes the conversion problem:
it can require discovering an inhabitant of
\(\Pi x:A.Mx=_B Nx\). Merely adding function extensionality as a
\emph{book-equality axiom} is a different operation and need not add such
reflection to conversion \cite[pp.~1168--1172]{BarendregtGeuvers2001}.
\end{hbdefinition}

\begin{hbdefinition}{def:p12-har-book-equality}{The three equality presentations}
The chapter distinguishes the following choices, keeping their sort roles.
\begin{enumerate}
 \item In its impredicative higher-order setting, Leibniz equality is
 \(t=_A q:=\Pi P:A\to\mathrm{Prop}.Pt\to Pq\).
 Quantification over all proposition-valued predicates is part of the
 allowed type theory, not a predicative abbreviation.
 \item Inductive equality \(\mathrm{Eq}_A:A\to A\to\mathrm{Prop}\)
 is generated by the sole constructor
 \(\mathrm{Refl}:\Pi x:A.\mathrm{Eq}_A x x\), together with the
 specified induction/elimination principle. A constructor declaration
 alone does not supply that principle.
 \item The Martin--Lof presentation uses
 \(\mathrm{Id}_A a b:\mathrm{Set}\), reflexivity, and identity
 elimination. The chapter displays transport from
 \(P:A\to\mathrm{Set}\), \(q:Pa\), and
 \(e:\mathrm{Id}_A a b\) to a term of \(Pb\).
 It does not give the full dependent identity eliminator there.
\end{enumerate}
Leibniz equality transports proposition-valued predicates. To transport
\(t:\mathrm{Fin}\,n\) along \(e:n=_{\mathrm{nat}}m\) to
\(\mathrm{Fin}\,m\), elimination into the relevant data sort is
needed. That property depends on the exact equality and elimination
restriction; it is not established by the word ``inductive'' or by a
modern assistant's name \cite[pp.~1170--1171]{BarendregtGeuvers2001}.
\end{hbdefinition}

\begin{hbtheorem}{thm:p12-har-leibniz}{proved=annex:p12.har.leibniz}{Leibniz equality and proposition transport}
In the displayed impredicative setting, \(=_A\) is reflexive,
symmetric and transitive. If \(e:a=_A b\), \(P:A\to\mathrm{Prop}\)
and \(u:Pa\), then \(e P u:Pb\). These terms transport only within
the admitted predicate class; they do not provide arbitrary data-valued
dependent transport.
\end{hbtheorem}

\begin{hbdefinition}{def:p12-har-setoids}{Setoids and their function space}
A setoid is a carrier \(A:\mathrm{Set}\) with an equivalence relation
\(=_A:A\to A\to\mathrm{Prop}\), including proofs of its laws.
The source's \([A,=_A]\) is a structured object, not itself the
carrier type \(A\). Given setoids \([A,=_A]\), \([B,=_B]\),
its function-space carrier and equality are
\[
\begin{aligned}
 A\to^s B&:=\Sigma f:A\to B.
      \Pi x,y:A.(x=_A y)\to(fx=_B fy),\\
 f=_{A\to^s B}g&:=\Pi x:A.(\pi_1f)x=_B(\pi_1g)x.
\end{aligned}
\]
The carrier contains a map \emph{and} its preservation proof. The equality
ignores that proof component. In particular
\(\langle f,p\rangle=_{A\to^s B}\langle f,q\rangle\) for any
two preservation proofs. This is a chosen book equality, not a declaration
that all host proof terms are definitionally equal.

For an equivalence relation \(Q\) on \(A\) with \(=_A\subseteq Q\),
the quotient setoid is \([A,Q]\), retaining the same carrier.
For \(P:A\to\mathrm{Prop}\), the subsetoid has carrier
\(\Sigma x:A.Px\) and equality
\(q(=_A\vert P)r:=(\pi_1q)=_A(\pi_1r)\).
The source does \emph{not} require \(P\) to respect \(=_A\).
It may retain only some representatives of a class
\cite[pp.~1172--1173]{BarendregtGeuvers2001}.
\end{hbdefinition}

\begin{hbtheorem}{thm:p12-har-setoid-constructions}{proved=annex:p12.har.setoid-constructions}{The three displayed setoid constructions}
The function-space equality and subsetoid equality are equivalence
relations. For equality-preserving \(f,g\), pointwise equality is
equivalent to
\(\Pi x,y:A.(x=_A y)\to fx=_B gy\).
If a map \(f:A\to B\) preserves \(Q\), it induces a
setoid map \([A,Q]\to[B,=_B]\) with the same underlying function.
These assertions need no definitional equality of preservation proofs.
\end{hbtheorem}

\begin{hbexample}{def:p12-har-rational-setoid}{Rational representatives and the noninvariant predicate}
For \(A=\mathrm{int}\times\mathrm{nat}\), set
\[
 \langle x,p\rangle=_A\langle y,q\rangle
       \quad\text{when}\quad x(q+1)=y(p+1),\qquad
 P\langle x,p\rangle\quad\text{when}\quad
       \gcd(x,p+1)=1.
\]
The encoded denominator is always positive. The subsetoid keeps the reduced
representative of each rational class, although \(P\) is not invariant:
\(\langle1,1\rangle=_A\langle2,3\rangle\), but only the first
pair satisfies \(P\). The reduced representative of zero is
\(\langle0,0\rangle\). This example illustrates the source's
reason for not demanding predicate invariance
\cite[p.~1173]{BarendregtGeuvers2001}.
\end{hbexample}

\subsection{Rational subsetoids and quotient mapping}

\begin{hbtheorem}{thm:p12-har-rational-normalization}{proved=annex:p12.har.rational-normalization}{Reduced representatives and the actual subsetoid isomorphism}
Use the source carrier \(A=\mathbb Z\times\mathbb N\), with
\((a,p)=_A(b,q)\) when \(a(q+1)=b(p+1)\), and
\(P(a,p)\) when \(\gcd(|a|,p+1)=1\). Its subsetoid
\(S=\Sigma z:A.Pz\) compares first components by \(=_A\),
ignoring the proof witnesses. Every class has exactly one coprime
integer-numerator/positive-denominator pair. The construction is
\[
 d=\gcd(|a|,p+1),\qquad
 r(a,p)=\left(a/d,(p+1)/d-1\right),
\]
together with the proved coprimality witness. The first-component
inclusion \(j:S\to A\) and this \(r:A\to S\) respect the
declared equalities, and \(jr\) and \(rj\) are pointwise equal
to their respective identities. This is a setoid isomorphism, not
definitional equality of all proof-witness terms. Zero has reduced pair
\((0,0)\), and negative numerators require no sign change of the
positive denominator \cite[p.~1173]{BarendregtGeuvers2001}.
\end{hbtheorem}

\begin{hbtheorem}{thm:p12-har-setoid-quotient-universal}{proved=annex:p12.har.setoid-quotient-universal}{The full quotient-setoid mapping property}
For a source setoid \([A,=_A]\) and equivalence relation \(Q\)
with \(=_A\subseteq Q\), the identity underlying map
\(q:[A,=_A]\to[A,Q]\) is equality preserving. Maps from
\([A,Q]\) to \([B,=_B]\) correspond precisely to source
setoid maps from \([A,=_A]\) that preserve \(Q\).
The descended map is the same underlying function and is unique up to
the source's pointwise function-setoid equality. No quotient selector,
choice axiom or host proof-irrelevance principle is required.
\end{hbtheorem}

\begin{proofboundary}{def:p12-har-fw-scope-boundary}{The equality and completeness regimes remain separate}
The lambda framework's \(\mathrm{eq}\) is a type representing
object beta equality. It is not host definitional equality, the
source's impredicative Leibniz equality, or a proof-dependent identity
eliminator. The three reflection theorems do not supply the source-cited
global extensional/setoid interpretation for every dependent constructor,
or the separately attributed weak completeness results for Martin--Lof
presentations. Those need their actual type, transport, universe and
elimination hypotheses. The rational and quotient arguments here close
the stated concrete setoid constructions only; no larger type theory
is identified with them.
\end{proofboundary}



\subsection{Coercions and the logical interpretation boundary}

\begin{hbdefinition}{def:p12-har-coercions}{Source coercions are explicit terms inserted by elaboration}
If a setoid is represented as
\(\Sigma A:\mathrm{Set}.\Sigma R:A\to A\to\mathrm{Prop}.
\mathrm{Is\_eq\_rel}\,A R\), it is not a type suitable as an
unexpanded variable annotation. Likewise,
\(\mathrm{Eq\_Rel}_A:=\Sigma R:A\to A\to\mathrm{Prop}.
\mathrm{Is\_eq\_rel}\,A R\) is a pair type, so an element
\(R:\mathrm{Eq\_Rel}_A\) cannot be applied directly as a relation.
A subsetoid element \(\langle a,p\rangle\) is not literally a
carrier element \(a\). The needed projections are real terms.

The chapter's coercion notation \(c:A\subseteq B\) marks a map for
implicit insertion, not a new ordinary subtype judgment. It has three uses:
a map to \(\mathrm{Set}\) or \(\mathrm{Prop}\) expands a
variable's structured annotation; a map to a function type expands a
structured term used in function position; a map to an argument's required
type expands the argument. For maps \(c_1:A\to B\),
\(c_2:B\to C\), composition is the term
\(\lambda x:A.c_2(c_1x)\). The source requires a unique declared
coercion path and no self-coercion, with a declaration-time graph check.
These are elaboration/coherence conditions, not an assertion that an
arbitrary graph with these names is coherent. The resulting expanded term
still must pass the underlying checker
\cite[pp.~1173--1175]{BarendregtGeuvers2001}.
\end{hbdefinition}

\begin{hbdefinition}{def:p12-har-interpretation}{Soundness, completeness, and strong completeness}
The chapter separates meaning in a chosen model from correspondence with
derivability in a chosen logic. For an appropriate source logic \(L\)
and its selected type theory \(S(L)\), soundness sends a
natural-deduction derivation \(\Delta\) of \(\varphi\) to a
typed proof term \([\Delta]\). Completeness requires that inhabitation
of the translated \(\varphi\) imply a derivation in \(L\).
Strong completeness additionally requires a translation of the given proof
term back into a derivation. These are three different obligations.

The distinction between domains, terms, formulas and derivations in
ordinary staged logic makes proof-dependent objects problematic for a
literal inverse translation. In a presentation with only one universe
\(\mathrm{Set}\), the context
\[
 A:\mathrm{Set},\ a:A,\ P:A\to\mathrm{Set},\ h:Pa,
 \ Q:(Pa)\to\mathrm{Set},\ f:(Pa)\to A
\]
can treat \(Pa\) as either a domain or a proposition; \(Q\) and
\(f\) then depend on proofs if \(h\) is read as a logical assumption.
A first-order logic does not automatically interpret such dependencies.
The chapter attributes weaker completeness and higher-order failure results
to specific cited presentations, not to all type theories with a shared
family name. Its discussion of higher-order sum interpretations and
impredicative existential encodings does not supply their full metatheoretic
proofs here \cite[pp.~1175--1178]{BarendregtGeuvers2001}.
\end{hbdefinition}

\begin{hbexample}{def:p12-har-proof-dependent-field}{The reciprocal and absolute-value examples}
An everywhere typed \(\mathrm{recip}:A\to A\) leaves
\(\mathrm{recip}\,0\) unspecified even if its reciprocal law is
restricted to nonzero arguments. Instead the source considers
\(\mathrm{recip}:(\Sigma x:A.x\ne0)\to A\).
To make this independent of the nonzero proof, choose a field setoid and
require \(\mathrm{recip}\) to preserve the subsetoid equality. Then
\[
 a=_A a',\quad p:a\ne_A0,\quad q:a'\ne_A0
 \quad\Longrightarrow\quad
 \mathrm{recip}\langle a,p\rangle=_A
 \mathrm{recip}\langle a',q\rangle.
\]
Taking \(a=a'\) demonstrates proof independence in the declared setoid
equality. It does not force equality of the proof components themselves.

In an ordered field, for
\(p:\Pi x:F.(x\ge0\lor x\le0)\), define
\(\mathrm{abs}_p x\) by the \(\mathrm{inl}\) branch \(x\)
and the \(\mathrm{inr}\) branch \(-x\). A second proof \(q\)
can choose a different branch at zero. The substantive independence claim
is \(\Pi x,x':F.(x=_F x')\to
\mathrm{abs}_p x=_F\mathrm{abs}_q x'\), not definitional equality
of the two functions \cite[pp.~1178--1180]{BarendregtGeuvers2001}.
\end{hbexample}

\begin{hbtheorem}{thm:p12-har-field-independence}{proved=annex:p12.har.field-independence}{The proof-independent field constructions}
For a reciprocal that preserves the source subsetoid equality, changing
the nonzero proof does not change its value up to \(=_A\).
In a linearly ordered field setoid whose order and operations respect
\(=_F\), with the displayed disjunction-to-data case elimination
admitted, the two absolute-value definitions above agree on equal
arguments for every pair of sign proofs \(p,q\).
\end{hbtheorem}

\begin{proofboundary}{def:p12-har-premise-bridge}{Interfaces and nonidentifications}
The abstract reduction and normalization terminology is specified at
\cref{def:p12-klop1987-ars,def:p12-klop1987-properties}; the arrow-only
example has the restricted interface of
\cref{def:p07-stlc-meta-setting,def:p07-curry-howard}. Neither interface
supplies the source's higher-order dependent \(\lambda\mathrm{HOL}\)
rules or its metatheory. The family specification in
\cref{def:p12-pts-lambda-cube} and the trust distinctions in
\cref{def:p12-computational-reflection} are compatible navigation points,
not proofs of source-specific adequacy, reduction, normalization, or
elimination claims. The source's logical-framework terminology and historical
assistant comparison must retain their own context and date.
\end{proofboundary}

\subsection{The chapter's constructive higher-order predicate logic}

\begin{hbdefinition}{def:p12-har-hol-language}{HOL domains, terms, and formulas (source Definition 3.1)}
The chapter's \(\mathrm{HOL}\) has domains
\(D::=B\mid\Omega\mid D\to D\), where the basic domains \(B\)
form a countable collection and \(\Omega\) is the domain of propositions.
For each \(\sigma\in D\), there are constants \(c_i^\sigma\)
and variables \(x_i^\sigma\). The typed term sets have the following
formation clauses: constants and variables of domain \(\sigma\) are
in \(\mathrm{Term}_\sigma\); for \(\varphi,\psi:\Omega\),
\(\forall x^\sigma.\varphi:\Omega\) and
\((\varphi\Rightarrow\psi):\Omega\); application takes
\(M:\sigma\to\tau\), \(N:\sigma\) to \(MN:\tau\);
abstraction takes \(M:\tau\) to
\(\lambda x^\sigma.M:\sigma\to\tau\).
Set \(\mathrm{Term}=\bigcup_{\sigma\in D}\mathrm{Term}_\sigma\)
and \(\mathrm{form}=\mathrm{Term}_\Omega\).
Substitution is capture avoiding, and beta conversion includes the
typed abstractions in this language.

Domains associate to the right, application to the left; application binds
more tightly than the logical implication \(\Rightarrow\).
There is no primitive product domain. A binary relation on \(A\) has
domain \(A\to A\to\Omega\), by currying. The source natural
deduction rules, including the eigenvariable condition and beta-conversion
rule, are on the registered sheet. They define provability from a finite
set of formulas, not classical higher-order logic with excluded middle
\cite[pp.~1181--1183, Fig.~1]{BarendregtGeuvers2001}.
\end{hbdefinition}

\begin{hbexample}{def:p12-har-hol-expressions}{All five source HOL language examples}
Take domains \(\mathbb N,A\), constants
\(>:\mathbb N\to\mathbb N\to\Omega\),
\(0:\mathbb N\), \(S:\mathbb N\to\mathbb N\), and
variables \(R,Q:A\to A\to\Omega\). The source expressions are:
\[
\begin{aligned}
 &\lambda x^{\mathbb N}.x>0:\mathbb N\to\Omega,\\
 \mathsf{ind}&:=\forall P^{\mathbb N\to\Omega}.
   P0\Rightarrow(\forall x^{\mathbb N}.Px\Rightarrow P(Sx))
       \Rightarrow\forall x^{\mathbb N}.Px,\\
 \mathsf{trans}(R)&:=\forall x^Ay^Az^A.
       Rxy\Rightarrow Ryz\Rightarrow Rxz,\\
 R\subseteq Q&:=\forall x^Ay^A.Rxy\Rightarrow Qxy,\\
 &\lambda x^Ay^A.\forall Q^{A\to A\to\Omega}.
      \mathsf{trans}(Q)\Rightarrow(R\subseteq Q)\Rightarrow Qxy
       :A\to A\to\Omega.
\end{aligned}
\]
The last expression is the source transitive-closure predicate: it
intersects all transitive relations containing \(R\). It does not add
reflexive closure. \(\mathsf{ind}\) is a formula expressing induction,
not a theorem inferred merely from the domain named \(\mathbb N\).
The sans-serif names are abbreviations, whereas \(Rab\) is always
successive application \((Ra)b\)
\cite[p.~1182, Example~3.3]{BarendregtGeuvers2001}.
\end{hbexample}

\begin{hbdefinition}{def:p12-har-hol-connectives}{The source impredicative logical encodings}
For \(\varphi,\psi:\Omega\), with quantified variables fresh,
\[
\begin{aligned}
 \varphi\mathbin{\&}\psi&:=\forall z^\Omega.
       (\varphi\Rightarrow\psi\Rightarrow z)\Rightarrow z,\\
 \varphi\lor\psi&:=\forall z^\Omega.
       (\varphi\Rightarrow z)\Rightarrow(\psi\Rightarrow z)\Rightarrow z,\\
 \bot&:=\forall z^\Omega.z,\qquad
 \neg\varphi:=\varphi\Rightarrow\bot,\\
 \exists x^\sigma.\varphi&:=\forall z^\Omega.
       (\forall x^\sigma.\varphi\Rightarrow z)\Rightarrow z,\\
 t=_A t'&:=\forall P^{A\to\Omega}.Pt\Rightarrow Pt'.
\end{aligned}
\]
These quantify over propositions or predicates at the same order as the
defined expression and are impredicative. The source calls its logic
constructive because it lacks double-negation elimination and has
disjunction/existence properties at its stated regime; it does not claim
that these encodings are predicative constructions. Their introduction and
elimination terms below give the local rule justification. The global
disjunction/existence-property proofs require the relevant normalization
and canonical-form analysis, not the formulas alone
\cite[pp.~1183--1185, Example~3.6]{BarendregtGeuvers2001}.
\end{hbdefinition}

\subsection{The actual three-sort system \(\lambda\mathrm{HOL}\)}

\begin{hbdefinition}{def:p12-har-lambda-hol}{Pseudo terms and the source Figure 2 calculus}
The source's \(\lambda\mathrm{HOL}\) has three sorts
\(\mathcal S=\{\mathrm{Prop},\mathrm{Type},\mathrm{Type}'\}\)
and pseudo terms
\[
 \mathcal T::=\mathrm{Prop}\mid\mathrm{Type}\mid\mathrm{Type}'
    \mid V\mid\Pi V:\mathcal T.\mathcal T
    \mid\lambda V:\mathcal T.\mathcal T\mid\mathcal T\mathcal T.
\]
Its axioms are \(\mathrm{Prop}:\mathrm{Type}\) and
\(\mathrm{Type}:\mathrm{Type}'\), with no type for
\(\mathrm{Type}'\). Product-sort pairs are exactly
\[
 (\mathrm{Type},\mathrm{Type}),\quad
 (\mathrm{Prop},\mathrm{Prop}),\quad
 (\mathrm{Type},\mathrm{Prop}).
\]
The complete seven-rule Figure~2 surface is on the source sheet. Contexts
are finite declarations of distinct variables; a new variable is fresh.
Conversion is \emph{beta equality of pseudo terms}, not the
beta/iota/delta combination of the earlier informal discussion. Definitions
and their delta rule are expressly omitted from this formal calculus.
A pseudo term is well typed if it occurs either to the left or to the
right of the colon in a derivable judgment. A context is well formed if
it occurs in some derivable judgment.

The three product pairs respectively form data-function domains,
implications, and universal propositions. In this specific system the
first two are nondependent; the last can depend on its bound data variable.
The absence of \((\mathrm{Prop},\mathrm{Type})\) prevents
data/functions or formulas from depending on proof arguments. The absence
of a product with domain sort \(\mathrm{Type}'\) prevents
abstraction over arbitrary data-type variables. These are system-specific
facts, not properties of every dependent type theory
\cite[pp.~1185--1187, Definition~3.7; Fig.~2]{BarendregtGeuvers2001}.
\end{hbdefinition}

\begin{hbdefinition}{def:p12-har-hol-translation}{An explicit direct HOL translation}
Following the source's direct interpretation, send a basic domain \(B\)
to a fresh declared \(B:\mathrm{Type}\), \(\Omega\) to
\(\mathrm{Prop}\), and domain arrows to type arrows. Translate
object constants and variables to declarations of their translated domain,
application homomorphically, and
\(\lambda x^\sigma.M\) to \(\lambda x:[\sigma].[M]\).
Translate a logical implication to \([\varphi]\to[\psi]\)
and a quantifier to \(\Pi x:[\sigma].[\varphi]\).
For a derivation \(\Delta\), put all domains, constants and free
object variables into \(\Gamma_\Delta\); an undischarged assumption
\(\varphi_i\) becomes a fresh \(z_i:[\varphi_i]\).
Logical proof abstraction/application, distinct from object
abstraction/application, encode hypothesis discharge and logical
elimination. The source pp.~1187--1190 motivates rather than defines these
maps; the explicit clauses here make the reconstruction unambiguous.
\end{hbdefinition}

\begin{hbexample}{def:p12-har-lambda-hol-judgments}{The source Example 3.10 judgments}
The following are the chapter's concrete formal-system instances, with
the indicated declarations supplying their contexts:
\begin{enumerate}
 \item \(\mathbb N:\mathrm{Type},0:\mathbb N,
 >:\mathbb N\to\mathbb N\to\mathrm{Prop}\) gives
 \(\lambda x:\mathbb N.x>0:\mathbb N\to\mathrm{Prop}\).
 \item \(\mathbb N:\mathrm{Type},0:\mathbb N,
 S:\mathbb N\to\mathbb N\) gives the formula
 \[
 \Pi P:\mathbb N\to\mathrm{Prop}.
 P0\to(\Pi x:\mathbb N.Px\to P(Sx))
       \to\Pi x:\mathbb N.Px:\mathrm{Prop}.
 \]
 \item \(A:\mathrm{Type},R:A\to A\to\mathrm{Prop}\) gives
 \(\Pi x,y,z:A.Rxy\to Ryz\to Rxz:\mathrm{Prop}\).
 \item \(A:\mathrm{Type}\) gives
 \[
 \lambda R,Q:A\to A\to\mathrm{Prop}.
       \Pi x,y:A.Rxy\to Qxy
       :(A\to A\to\mathrm{Prop})\to
        (A\to A\to\mathrm{Prop})\to\mathrm{Prop}.
 \]
 \item \(A:\mathrm{Type}\) gives
 \(\lambda x,y:A.\Pi P:A\to\mathrm{Prop}.Px\to Py
       :A\to A\to\mathrm{Prop}\), the equality relation.
 \item \(A:\mathrm{Type},x:A,y:A\) gives the equality-symmetry
 term in \cref{def:p12-har-encoded-proof-terms}, at type
 \((x=_A y)\to(y=_A x)\).
\end{enumerate}
For the source derivation of irreflexivity, the formal context is
\(A:\mathrm{Type},R:A\to A\to\mathrm{Prop},
z:\Pi x,y:A.Rxy\to Ryx\to\bot\), and the checked term is
\(\lambda x:A.\lambda q:Rxx.zxxqq:
\Pi x:A.Rxx\to\bot\). It specializes the application chain of
\cref{thm:p12-har-asymmetry}; it does \emph{not} abstract over
\(A:\mathrm{Type}\) inside \(\lambda\mathrm{HOL}\), which
has no permitted product with domain sort \(\mathrm{Type}'\).
Thus the earlier five-binder informal Set-based term and this concrete
formal judgment keep different quantification regimes
\cite[pp.~1187--1189, Example~3.10]{BarendregtGeuvers2001}.
\end{hbexample}

\begin{hbtheorem}{thm:p12-har-hol-forward}{proved=annex:p12.har.hol-forward}{The direct HOL embedding, with conversion}
Every HOL object term \(t:\sigma\) translates to a term of
\([\sigma]\); every HOL formula translates to a term of
\(\mathrm{Prop}\). Every derivation of \(\varphi\) from
\(\varphi_1,\ldots,\varphi_n\) translates to a derivation
\[
 \Gamma_\Delta,z_1:[\varphi_1],\ldots,z_n:[\varphi_n]
       \vdash_{\lambda\mathrm{HOL}}[\Delta]:[\varphi].
\]
The translation commutes with capture-avoiding substitution and preserves
beta conversion. No classical axiom is added.
\end{hbtheorem}

\begin{hbexample}{def:p12-har-encoded-proof-terms}{Connective terms and the source equality-symmetry term}
In \(\lambda\mathrm{HOL}\), replace source HOL quantification by
\(\Pi\) and implication by \(\to\) in the impredicative encodings.
Thus \(\varphi\mathbin{\&}\psi=
\Pi\alpha:\mathrm{Prop}.(\varphi\to\psi\to\alpha)\to\alpha\).
The source's complete conjunction terms are
\[
\begin{aligned}
 \pi_1&:=\lambda p:(\varphi\mathbin{\&}\psi).
       p\varphi(\lambda x:\varphi.\lambda y:\psi.x),\\
 \pi_2&:=\lambda p:(\varphi\mathbin{\&}\psi).
       p\psi(\lambda x:\varphi.\lambda y:\psi.y),\\
 \langle-,-\rangle&:=\lambda x:\varphi.\lambda y:\psi.
       \lambda\alpha:\mathrm{Prop}.
       \lambda h:\varphi\to\psi\to\alpha.hxy.
\end{aligned}
\]
These are proposition-level encodings, not primitive data products.
For disjunction the two introductions choose the corresponding supplied
branch; its eliminator at \(\chi:\mathrm{Prop}\) is
\(\lambda p.\lambda f:\varphi\to\chi.
\lambda g:\psi\to\chi.p\chi f g\), with \(p\) annotated
by the encoded disjunction. Falsehood elimination at \(\chi\) is
\(\lambda p:\bot.p\chi\). Existential introduction from
\(t:A\), \(u:\varphi[t/x]\) is
\(\lambda\alpha:\mathrm{Prop}.
\lambda h:(\Pi x:A.\varphi\to\alpha).htu\);
elimination at \(\chi\), with \(x\) absent from \(\chi\), is
\(\lambda p:(\exists x:A.\varphi).
\lambda h:(\Pi x:A.\varphi\to\chi).p\chi h\).

For source Leibniz equality, its second, explicitly displayed symmetry
term is
\[
 \lambda r:(x=_A y).\lambda P:A\to\mathrm{Prop}.
   r(\lambda z:A.Pz\to Px)(\lambda q:Px.q)
       :(x=_A y)\to(y=_A x).
\]
It instantiates the predicate with \(Pz\to Px\), different from
the also valid equality-valued predicate used in
\cref{thm:p12-har-leibniz}; the two constructions are not conflated
\cite[pp.~1187--1189, Examples~3.10 and connective discussion]{BarendregtGeuvers2001}.
\end{hbexample}

\begin{hbtheorem}{thm:p12-har-connective-checks}{proved=annex:p12.har.connective-checks}{Typing and beta computations for the impredicative encodings}
The conjunction terms have types
\((\varphi\mathbin{\&}\psi)\to\varphi\),
\((\varphi\mathbin{\&}\psi)\to\psi\), and
\(\varphi\to\psi\to(\varphi\mathbin{\&}\psi)\).
Their projections of an encoded pair beta-reduce to its corresponding
component. The displayed disjunction, falsehood and existential terms
give the introduction/elimination rules at proposition-valued conclusions,
with the usual existential eigenvariable restriction. The displayed
source symmetry term has exactly its stated equality type.
\end{hbtheorem}

\begin{hblemma}{thm:p12-har-hol-contexts}{proved=annex:p12.har.hol-contexts}{Valid contexts, projection, and thinning for the exact source system}
A context is valid when it is empty, or extends a valid prefix by a
fresh \(x:A\) with an annotation judgment \(A:s\). This is
equivalent to occurring in a source \(\lambda\mathrm{HOL}\)
judgment. Every declaration has its projection judgment in the complete
context. A judgment remains derivable in any valid context containing
all its original literal declarations. Bound variables may be
consistently alpha-renamed fresh for the enlarged context.
\end{hblemma}

\begin{hblemma}{thm:p12-har-hol-generation}{proved=annex:p12.har.hol-generation}{Generation and beta-product injectivity}
Removing terminal conversion and weakening from a source typing judgment
recovers the premises of its subject's constructor rule, at a generated
type beta-equal to its given type. This covers variables, both typed
sorts, products, abstractions and applications; \(\mathrm{Type}'\)
has no subject typing judgment. Beta-equal products have beta-equal
domains and bodies under the same fresh binder. Distinct sorts, and a
sort versus a product, are not beta-equal.
\end{hblemma}

\begin{hblemma}{thm:p12-har-hol-type-correctness}{proved=annex:p12.har.hol-type-correctness}{Type correctness}
If \(\Gamma\vdash M:A\) in the source system, either \(A\)
is a sort or \(\Gamma\vdash A:s\) for a source sort \(s\).
The application case uses product generation and the actual-argument
substitution theorem, not subject reduction or normalization.
\end{hblemma}

\begin{hblemma}{thm:p12-har-hol-context-conversion}{proved=annex:p12.har.hol-context-conversion}{Replacing a typed annotation by a convertible typed annotation}
If \(\Gamma,x:B,\Theta\) is valid,
\(\Gamma\vdash B':s'\), and \(B=_\beta B'\), then every
judgment in that context is derivable in
\(\Gamma,x:B',\Theta\), with its subject, type and later raw
annotations unchanged. The new context is valid. Both annotation
judgments are required; an untyped raw beta-equivalent expression is
not sufficient.
\end{hblemma}

\begin{hblemma}{thm:p12-har-disambiguation}{proved=annex:p12.har.disambiguation}{Source Lemma 3.11: the three context strata}
A well-formed \(\lambda\mathrm{HOL}\) context can be permuted to
\(\Gamma_D,\Gamma_L,\Gamma_P\), preserving judgments, where
\(\Gamma_D\) declares domain variables \(A:\mathrm{Type}\),
\(\Gamma_L\) declares object variables of types formed in
\(\Gamma_D\), and \(\Gamma_P\) declares proofs of formulas
formed in \(\Gamma_D,\Gamma_L\). Data types are simple arrows
ending in a declared domain variable or \(\mathrm{Prop}\); data
terms require no proof declarations. If a product of sort
\(\mathrm{Prop}\) has a proposition-sorted domain, its codomain
does not depend on that proof variable. The annex supplies the full
domain/object grammar classification, beta-stable separation and
smaller-context reconstruction; it does not assume a general
strengthening principle to obtain this result.
\end{hblemma}

\begin{hblemma}{thm:p12-har-hol-uniqueness}{proved=annex:p12.har.hol-uniqueness}{Uniqueness of source lambdaHOL types}
For the exact source specification, if
\(\Gamma\vdash M:A\) and \(\Gamma\vdash M:B\), then
\(A=_\beta B\). The proof uses generation, product injectivity
and the functional sort rules, independently of checker completeness.
\end{hblemma}

\begin{hbtheorem}{thm:p12-har-hol-decoder}{proved=annex:p12.har.hol-decoder}{The inverse direct HOL proof construction}
Using the proved disambiguation, generation and type correctness,
over a translated HOL domain/object signature,
every derivation \(\Gamma\vdash M:\varphi\), with
\(\Gamma\vdash\varphi:\mathrm{Prop}\), can be decoded to a
HOL natural deduction of the decoded formula from the formulas in
\(\Gamma_P\). The decoding handles each typing rule and conversion;
it is not an assertion of faithfulness for the separate framework
signatures of \cref{def:p12-har-framework}.
\end{hbtheorem}

\subsection{The checker and its actual proof dependencies}

\begin{hbtheorem}{thm:p12-har-pseudo-confluence}{proved=annex:p12.har.pseudo-confluence}{Beta confluence on the source pseudo terms}
For the source grammar over any fixed alphabet of sort atoms, with
variables, annotated abstraction,
dependent product and application, compatible beta reduction is confluent
on all pseudo terms: if \(M\twoheadrightarrow_\beta N_1\) and
\(M\twoheadrightarrow_\beta N_2\), there is \(P\) with
\(N_1\twoheadrightarrow_\beta P\) and
\(N_2\twoheadrightarrow_\beta P\). No normalization, typing,
eta rule, inductive computation, or definition-unfolding rule is used.
\end{hbtheorem}

\begin{hbtheorem}{thm:p12-har-hol-subject-reduction}{proved=annex:p12.har.hol-subject-reduction}{Beta subject reduction for exact source lambdaHOL}
If \(\Gamma\vdash M:T\) and
\(M\twoheadrightarrow_\beta M'\), then
\(\Gamma\vdash M':T\). Reduction in annotations and dependent
product bodies is included. Combined with raw confluence, any two
reduction paths from a typed term have a common reduct of the same type.
No normalization or checker premise is used.
\end{hbtheorem}

\begin{hbdefinition}{def:p12-har-metatheory-premises}{The separate source metatheory results}
The source asserts beta confluence on \emph{all pseudo terms}
(Proposition~3.12), subject reduction for typed terms (3.13), confluence
with a typed common reduct (3.14), and strong normalization of typed terms
(3.15). Its discussion points outside the chapter for the involved
normalization proof. These assertions have distinct uses: confluence
compares conversion paths; subject reduction keeps reductions in the
typed regime; normalization terminates the checker comparisons. They do
not follow from soundness of a checker that has assumed those comparisons.
Raw confluence, typed preservation and the structural classification
are now proved in
\cref{thm:p12-har-pseudo-confluence,thm:p12-har-hol-subject-reduction,thm:p12-har-disambiguation}.
Strong normalization is supplied by the complete exact-CC argument
and separately proved transport in
\cref{thm:p12-har-cc-normalization,thm:p12-har-hol-cc-transport,thm:p12-har-hol-normalization}.
The checker uses these results, rather than deriving them from its own
claimed termination
\cite[pp.~1191--1192, Propositions~3.12--3.15]{BarendregtGeuvers2001}.
\end{hbdefinition}

\begin{hbdefinition}{def:p12-har-type-algorithm}{Source Definition 3.16 and the corrected terminating guard}
Let \(\mathrm{false}\) mean failure, not a term. The source mutual
functions are \(\mathrm{Ok}(\Gamma)\) and
\(\mathrm{Type}_\Gamma(M)\), with these clauses. All context
extensions require freshness, and unsuccessful tests return failure.
\begin{enumerate}
 \item \(\mathrm{Ok}(\langle\rangle)=\mathrm{true}\);
 \(\mathrm{Ok}(\Gamma,x:A)\) tests
 \(\mathrm{Type}_\Gamma(A)\in\mathcal S\).
 \item A variable returns its declared type if
 \(\mathrm{Ok}(\Gamma)\) and its declaration is present.
 \(\mathrm{Prop}\) returns \(\mathrm{Type}\), and
 \(\mathrm{Type}\) returns \(\mathrm{Type}'\), each only in an
 accepted context. \(\mathrm{Type}'\) always fails.
 \item For \(MN\), compute \(C=\mathrm{Type}_\Gamma(M)\)
 and \(D=\mathrm{Type}_\Gamma(N)\). Expose
 \(\Pi x:A.B\) by weak-head normalization of \(C\), test
 \(A=_\beta D\), and return \(B[N/x]\).
 \item For \(\Pi x:A.B\), compute
 \(s_1=\mathrm{Type}_\Gamma(A)\),
 \(s_2=\mathrm{Type}_{\Gamma,x:A}(B)\). Return \(s_2\)
 exactly when both are sorts and their pair is one of the three allowed
 pairs.
 \item For \(\lambda x:A.M\), first compute
 \(B=\mathrm{Type}_{\Gamma,x:A}(M)\). The first source version
 then recursively checks \(\mathrm{Type}_\Gamma(\Pi x:A.B)\)
 and returns \(\Pi x:A.B\) if its result is a sort.
\end{enumerate}
The last recursive input can exceed the original term's size; its
termination is not established by the source's proposed size measure.
The source Definition~3.20 replaces that input by a syntactic guard.
For the corrected version, call \(B\) a data-type shape if it is a
simple arrow type ending in \(\mathrm{Prop}\) or a variable declared
\(z:\mathrm{Type}\) in \(\Gamma\). On a successfully typed
body, use the following guard instead:
\begin{enumerate}
 \item If \(B\) has that data-type shape, require
 \(\mathrm{Type}_\Gamma(A)=\mathrm{Type}\).
 \item Otherwise, if \(B\) is neither \(\mathrm{Type}\) nor
 \(\mathrm{Type}'\), require
 \(\mathrm{Type}_\Gamma(A)\in\{\mathrm{Prop},\mathrm{Type}\}\).
 \item Otherwise fail.
\end{enumerate}
The second branch keeps \emph{both} allowed product pairs with
proposition-valued codomain. The published guard retains only
\(\mathrm{Prop}\) there, which would reject universal-proof
abstraction; it is not the active algorithm used here. The classification
requires the disambiguation and type-correctness premises, not a guess
about an arbitrary pseudo term's sort
\cite[pp.~1192--1195, Definitions~3.16, 3.20; Lemma~3.19]{BarendregtGeuvers2001}.
\end{hbdefinition}

\begin{hbtheorem}{thm:p12-har-guard-counterexample}{proved=annex:p12.har.guard-counterexample}{A valid universal abstraction rejected by the printed guard}
In \(\Gamma=A:\mathrm{Type},P:A\to\mathrm{Prop},
h:\Pi x:A.Px\), the term \(\lambda x:A.hx\) has type
\(\Pi x:A.Px\). Its body type \(Px\) is not a data-type
shape, and the printed second guard would demand
\(\mathrm{Type}_\Gamma(A)=\mathrm{Prop}\), though the result
is \(\mathrm{Type}\). Hence the source's printed modification
is incomplete. The corrected second branch accepts this case.
\end{hbtheorem}

\begin{hbtheorem}{thm:p12-har-checker-conditional}{proved=annex:p12.har.checker-conditional}{Corrected type synthesis and checking for exact source lambdaHOL}
With the proved strong normalization, confluence, subject
reduction, type correctness, generation, substitution, context conversion,
uniqueness and disambiguation results, the corrected mutual algorithm
terminates on all pseudo-term/context inputs, accepts exactly the
well-formed contexts, synthesizes a type exactly for typable terms, and
returns a type beta-equal to every derivable type of that term. Therefore
comparison with that synthesized type decides TCP. The theorem is
specific to the declared beta-only system and corrected guard. It does
not establish checking for the warned or inductive extensions.
\end{hbtheorem}

\begin{hbdefinition}{def:p12-har-syntax-directed}{The source modified judgment surface}
Source Definition~3.22 introduces \(\Gamma\vdash^{\mathrm{tc}}M:A\)
and \(\Gamma\vdash^{\mathrm{tc}}\mathrm{ok}\). It removes
the ordinary axiom, weakening, variable and conversion rules; adds empty
context, context extension, projection and sort rules; retains product
and abstraction; and replaces application by a weak-head product exposure
and beta comparison. Its complete rules are on the sheet. For completeness
of these modified rules the meaningful statement is
\[
 \Gamma\vdash M:A\ \Longrightarrow\
 \exists A'.\bigl(A'=_\beta A\text{ and }
       \Gamma\vdash^{\mathrm{tc}}M:A'\bigr),
\]
with the modified turnstile in the consequent. This supplies a
syntax-directed intermediate obligation rather than assuming that
reversing any derivation's last rule automatically reverses conversion
and weakening \cite[pp.~1195--1196, Definition~3.22; Proposition~3.23]{BarendregtGeuvers2001}.
\end{hbdefinition}

\subsection{The source PTS specifications and genuine maps}

\begin{hbdefinition}{def:p12-har-pts-specifications}{Source Definition 3.24 and its four distinct specifications}
A source PTS \(\lambda(\mathcal S,\mathcal A,\mathcal R)\)
has sort axioms \(\mathcal A\subseteq\mathcal S^2\) and product
triples \(\mathcal R\subseteq\mathcal S^3\), with the full
Figure~3 rules on the sheet. The pair \((s_1,s_2)\) abbreviates
\((s_1,s_2,s_2)\), not an arbitrary result-sort choice.
Besides the already specified \(\lambda\mathrm{HOL}\), the chapter
gives these distinct systems:
\begin{enumerate}
 \item \(\lambda\mathrm{PRED}\):
 \(\mathcal S=\{\mathrm{Set},\mathrm{Type}^s,
 \mathrm{Prop},\mathrm{Type}\}\),
 \(\mathcal A=\{\mathrm{Set}:\mathrm{Type}^s,
 \mathrm{Prop}:\mathrm{Type}\}\), and
 \[
 \mathcal R=\{(\mathrm{Set},\mathrm{Set}),
 (\mathrm{Set},\mathrm{Type}),(\mathrm{Prop},\mathrm{Prop}),
 (\mathrm{Set},\mathrm{Prop})\}.
 \]
 \item \(\lambda\mathrm{PRED}\omega\) has the same sorts and
 axioms, and extends those pairs by
 \((\mathrm{Type},\mathrm{Type})\),
 \((\mathrm{Type},\mathrm{Prop})\).
 \item \(\mathrm{CC}\): \(\mathcal S=\{\ast,\Box\}\),
 \(\mathcal A=\{\ast:\Box\}\), and
 \(\mathcal R=\{(\ast,\ast),(\ast,\Box),
 (\Box,\ast),(\Box,\Box)\}\).
 \item \(\lambda\ast\): \(\mathcal S=\{\ast\}\),
 \(\mathcal A=\{\ast:\ast\}\),
 \(\mathcal R=\{(\ast,\ast)\}\).
\end{enumerate}
\(\lambda\mathrm{PRED}\) is minimal first-order predicate logic
\emph{with higher-order function domains}; it forbids quantification
over predicate domains, not every higher-order function. The ordinary
connectives are not recovered there by quantifying over propositions.
\(\lambda\mathrm{PRED}\omega\) adds the stated predicate-domain
formation/quantification. The two-sort \(\mathrm{CC}\) is the source
Calculus of Constructions, not an inductive extension. \(\lambda\ast\)
is inconsistent in the source's stated logical sense and is not even
weakly normalizing: the actual typed witness, transport and normal-head
obstruction are supplied in
\cref{thm:p12-har-hurkens-bottom,thm:p12-har-star-inconsistency,thm:p12-har-star-no-normal-form}
\cite[pp.~1196--1202]{BarendregtGeuvers2001}.
\end{hbdefinition}

\begin{hbdefinition}{def:p12-har-pts-premises}{General and restricted source PTS properties}
The source separates general PTS thinning, substitution, strengthening
and subject reduction (Propositions~3.25--3.28) from uniqueness of types
(3.30), which it states for functional specifications. Functionality
requires each sort-axiom source to have at most one target and each fixed
pair of product input sorts to have at most one result sort. General
strengthening removes \(x:B\) only when \(x\) is absent from
the later context, term and type; it is not proved by substituting an
inhabitant of \(B\), which need not exist.

The chapter points to separate metatheoretic sources for these proofs.
It also states strong normalization of \(\mathrm{CC}\), then transfers
it to specified mapped systems. None of these assertions establishes
normalization for arbitrary PTSs; the displayed \(\lambda\ast\)
is its explicit contrary case \cite[pp.~1199--1201]{BarendregtGeuvers2001}.
\end{hbdefinition}

\begin{hbtheorem}{thm:p12-har-pts-substitution}{proved=annex:p12.har.pts-substitution}{Source Proposition 3.26: substitution with the entire later context}
In the exact Figure~3 PTS calculus, if
\(\Gamma_1,x:B,\Gamma_2\vdash M:A\) and
\(\Gamma_1\vdash N:B\), then
\[
 \Gamma_1,\Gamma_2[N/x]\vdash M[N/x]:A[N/x].
\]
Substitution acts on every later declaration's type as well as the
term and its type. Bound variables are freshened before substitution.
This theorem requires no inhabitant of a type unrelated to the given
\(B\), no strengthening premise, and no normalization theorem.
\end{hbtheorem}

\begin{hbdefinition}{def:p12-har-pts-morphism}{PTS morphisms and the actual comparison map}
A PTS morphism \(f\) maps sorts to sorts and preserves every sort
axiom and product triple. Extend it to terms and contexts homomorphically,
fixing variables and preserving abstraction, product and application.
The source map from \(\lambda\mathrm{PRED}\omega\) to
\(\lambda\mathrm{HOL}\) is
\[
 f(\mathrm{Prop})=\mathrm{Prop},\quad
 f(\mathrm{Set})=f(\mathrm{Type})=\mathrm{Type},\quad
 f(\mathrm{Type}^s)=\mathrm{Type}'.
\]
It identifies two source sorts but does not identify the systems'
published primitive syntax. A morphism to \(\mathrm{CC}\) sends
\(\mathrm{Set},\mathrm{Prop}\) to \(\ast\), and
\(\mathrm{Type}^s,\mathrm{Type}\) to \(\Box\).
Each contracted beta-redex maps to a contracted beta-redex at the same
position. Noninjective sort maps need not reflect the relation between
\emph{arbitrary chosen preimages}; that stronger claim is unnecessary
for transfer of an infinite reduction path
\cite[pp.~1200--1202, Definition~3.31; Proposition~3.32]{BarendregtGeuvers2001}.
\end{hbdefinition}

\begin{hbtheorem}{thm:p12-har-pts-morphism}{proved=annex:p12.har.pts-morphism}{Judgment and reduction transport along a PTS morphism}
For any sort-axiom/triple-preserving PTS morphism,
\(\Gamma\vdash_T M:A\) implies
\(f(\Gamma)\vdash_{T'}f(M):f(A)\). Every specified beta
contraction is preserved. Consequently, if the target PTS is strongly
normalizing, every well-typed source term is strongly normalizing.
This is a transfer theorem with a target-normalization premise.
\end{hbtheorem}

\begin{hbtheorem}{thm:p12-har-inverse-counterexample}{proved=annex:p12.har.inverse-counterexample}{Failure of the source's printed reverse comparison}
In source \(\lambda\mathrm{HOL}\),
\[
 A:\mathrm{Type}\vdash(A\to\mathrm{Prop})\to A:\mathrm{Type}.
\]
The printed reverse clauses send the context to
\(A:\mathrm{Set}\), the domain \(A\to\mathrm{Prop}\)
has sort \(\mathrm{Type}\) in source
\(\lambda\mathrm{PRED}\omega\), and the codomain \(A\)
has sort \(\mathrm{Set}\). The required outer product pair
\((\mathrm{Type},\mathrm{Set})\) is absent. Thus this proposed
reverse translation does not preserve that judgment; it cannot justify
normalization transfer by the claimed equivalence.
\end{hbtheorem}

\begin{proofboundary}{def:p12-har-pts-inverse}{The printed reverse comparison is not a completed equivalence}
The source's reverse map uses
\(g(\mathrm{Prop})=\mathrm{Prop}\),
\(g(\mathrm{Type})=\mathrm{Set}\),
\(g(\mathrm{Type}')=\mathrm{Type}^s\), but overrides the
result sort of \(M:\mathrm{Type}\) according to its data-type
shape: a chain ending in a domain variable receives \(\mathrm{Set}\),
while one ending in \(\mathrm{Prop}\) receives \(\mathrm{Type}\).
It is not an ordinary inverse sort-map morphism, and the stated
judgment-dependent clauses also fail on
\cref{thm:p12-har-inverse-counterexample}. The valid forward morphism
does not supply an inverse. The source attributes HOL normalization
transfer to this comparison. The valid replacement is the separately
proved judgment-dependent direct CC bridge and its exact normalization
argument in
\cref{thm:p12-har-hol-cc-transport,thm:p12-har-cc-normalization,thm:p12-har-hol-normalization}.
It is not the unmodified source \(\lambda\mathrm{PRED}\omega\)
\cite[pp.~1201--1202]{BarendregtGeuvers2001}.
\end{proofboundary}

\begin{hbtheorem}{thm:p12-har-hol-cc-transport}{proved=annex:p12.har.hol-cc-transport}{A valid judgment-dependent lambdaHOL-to-CC bridge}
Define \(h(\mathrm{Prop})=h(\mathrm{Type})=\ast\),
\(h(\mathrm{Type}')=\Box\), and extend homomorphically on
raw syntax, fixing variables. For a domain expression \(D\), let
\(\sigma(D)=\ast\) when its rightmost atom is a declared domain
variable, and \(\sigma(D)=\Box\) when that atom is
\(\mathrm{Prop}\). Valid contexts map annotation-wise, and
source judgments map by these classified clauses:
\[
\begin{array}{rcl}
 \mathrm{Type}:\mathrm{Type}'&\mapsto&\ast:\Box,\\
 D:\mathrm{Type}&\mapsto&hD:\sigma(D),\\
 O:D&\mapsto&hO:hD,\\
 p:F&\mapsto&hp:hF.
\end{array}
\]
Here \(O\) is an object expression, \(F\) a proposition, and
\(p\) a proof term in the proved stratification.
Each mapped judgment is derivable in exact source CC, and each designated
beta contraction is preserved. Thus a separately proved strong-normalization
theorem for that exact CC entails source lambdaHOL strong normalization.
This is an editorial replacement for the failed printed inverse, not a
PTS sort morphism and not a proof of target CC normalization.
\end{hbtheorem}

\begin{hbdefinition}{def:p12-har-cc-normalization-method}{The separate 1994 primary proof and its candidate adaptation}
For exact two-sort CC, Geuvers's December~1994 Computing Science
Report~94/50 gives a direct two-interpretation normalization proof.
The held report is distinct from the 1995 LNCS publication cited by
the chapter. The selected base argument is sections~2--3, printed
pp.~2--10; later report extensions are not imported here
\cite[pp.~2--10]{Geuvers1994Report9450}.

The report's \(\mathsf{SAT}\) consists of sets of untyped lambda
terms contained in \(\mathsf{SN}\), containing its base terms
(variables or inert \(d\), closed under application to SN arguments),
and closed under its specified key-redex expansion when the expanded
term is SN. The key redex is a redex itself, or the inherited key redex
of an application's left child. The proof below explicitly uses the
alternative reduction-closed neutral-expansion candidates:
\(R\subseteq\mathsf{SN}\), closure under contraction, and
membership of a neutral term when all its immediate reducts belong
to \(R\). Variables, \(d\), and applications are neutral;
lambdas are not. This is a declared proof-method adaptation, not an
identification of the definitions. Every closure, controlled expansion,
evaluation/substitution and valuation property used is proved in the annex.

Kinds \(A:\Box\) have finite shape \(\kappa(A)\), with
\(\kappa(\ast)=\mathbf b\), object-domain products taking the
codomain shape, and constructor-domain products taking the arrow of
their domain and codomain shapes. If \(\mathcal C\) is the set of
these candidates, define
\(V(\mathbf b)=\mathcal C\) and
\(V(k\to\ell)=V(\ell)^{V(k)}\). These are ordinary finite
iterated function sets, not an impredicative universe assumed to contain
itself. Constructor valuation \(\xi\) and raw-term valuation
\(\rho\) are distinct. The latter interpretation retains annotations:
\[
\begin{aligned}
 E(\ast)&=d,&E(v)&=v,&E(MN)&=E(M)E(N),\\
 E(\lambda v:A.M)&=\mathrm K(\lambda v.E(M))E(A),\\
 E(\Pi v:A.B)&=dE(A)(\lambda v.E(B)),&
 \mathrm K&=\lambda p.\lambda q.p.
\end{aligned}
\]
The report writes its raw interpretation as \((M)_\rho\),
with equivalent syntactic code in Fact~3.12; here that code is
\(E(M)\) and \(E_\rho(M)=E(M)\rho\). The constructor/kind
interpretation \([\cdot]_\xi\) follows the same two-level
architecture over the declared candidate variant. Each source beta
contraction has a positive target path, including annotation steps
\cite[Definitions~3.1--3.12; Theorems~3.14--3.15]{Geuvers1994Report9450}.
\end{hbdefinition}

\begin{hbtheorem}{thm:p12-har-cc-normalization}{proved=annex:p12.har.cc-normalization}{Strong normalization for exact source CC}
For \(\mathcal S=\{\ast,\Box\}\),
\(\mathcal A=\{\ast:\Box\}\), and all four product triples
\((\ast,\ast,\ast),(\ast,\Box,\Box),
(\Box,\ast,\ast),(\Box,\Box,\Box)\), with Figure~3 typing
and full compatible annotated beta reduction, every pseudo term on
either side of a derivable judgment has no infinite beta-reduction path.
Open valid contexts, reductions beneath binders, and reductions in
types/annotations are included. No CC normalization, checker correctness
or inhabited-context premise is assumed in the proof.
\end{hbtheorem}

\begin{hbtheorem}{thm:p12-har-hol-normalization}{proved=annex:p12.har.hol-normalization}{Strong normalization for the exact beta-only lambdaHOL}
Every pseudo term well typed in the exact three-sort lambdaHOL of
\cref{def:p12-har-lambda-hol} is strongly normalizing for full compatible
annotated beta reduction. The proof uses the proved judgment-dependent
CC transport and exact CC normalization, not the failed printed PRED-omega
inverse. Definitions/delta, primitive recursor/iota, strong-sum and
inductive extensions are outside this conclusion.
\end{hbtheorem}

\subsection{The exact one-sort paradox and its typed witness}

\begin{hbdefinition}{def:p12-har-star-specification}{The source lambda-star system}
The chapter's \(\lambda\ast\) has
\(\mathcal S=\{\ast\}\), \(\mathcal A=\{(\ast,\ast)\}\),
and \(\mathcal R=\{(\ast,\ast,\ast)\}\). Its pseudo terms are
\[
 \mathcal T::=\ast\mid v\mid\Pi v:\mathcal T.\mathcal T
       \mid\lambda v:\mathcal T.\mathcal T\mid\mathcal T\mathcal T.
\]
Use exactly the \hyperref[sheet:p12-har-pts]{seven registered Figure~3 PTS rules}, at that specification,
fresh context variables and capture-avoiding substitution, with compatible
beta reduction also in annotations. There is no added equality reflection,
inductive type, fixed point, choice or classical axiom. The false type is
\(\bot_\ast:=\Pi p:\ast.p\). This sort self-typing regime is
not the two-sort CC or a stratified universe hierarchy
\cite[p.~1201]{BarendregtGeuvers2001}.
\end{hbdefinition}

\begin{hbdefinition}{def:p12-har-hurkens-source}{The actual Hurkens source hierarchy and notation}
Hurkens's selected construction uses \(\lambda U^-\) with sorts
\(\ast,\Box,\triangle\), axioms
\(\ast:\Box\), \(\Box:\triangle\), and product pairs
\[
 (\ast,\ast),\quad(\Box,\Box),\quad(\Box,\ast),
 \quad(\triangle,\Box),
\]
where a pair \((s,t)\) means \((s,t,t)\). The paper's
\(\lambda U\) additionally permits \((\triangle,\ast)\);
that additional pair is not used here. Preserve the separate sort levels.
For \(X:\Box\), its predicate-space notation is
\(\wp X:=X\to\ast\), not a primitive set-theoretic powerset.
\(\Lambda X:\Box.M\) is ordinary annotated abstraction, and
\(\{bX\}\) is ordinary application at the displayed universe argument.
For small formulas \(\varphi,\chi:\ast\),
\(\varphi\Rightarrow\chi:=\Pi 0:\varphi.\chi\), with an
unused proof binder, \(\bot:=\Pi p:\ast.p\), and
\(\neg\varphi:=\varphi\Rightarrow\bot\). Numerals used as
proof-binder names are the source's assumption labels, not arithmetic
constants. All helper definitions below are finite syntactic abbreviations;
their expansion introduces no new primitive reduction rule
\cite[sections~2--3, pp.~266--269]{Hurkens1995LambdaStar}.
\end{hbdefinition}

\begin{hbdefinition}{def:p12-har-hurkens-data}{The source short universe and its explicit two recipes}
The selected short construction is
\[
 \mathcal U:=\Pi X:\Box.
       ((\wp\wp X\to X)\to\wp\wp X).
\]
For \(t:\wp\wp\mathcal U\) and \(s:\mathcal U\), use the
source substitution recipes
\[
\begin{aligned}
 \tau t&:=\Lambda X:\Box.\lambda f:\wp\wp X\to X.
      \lambda p:\wp X.
      t(\lambda x:\mathcal U.p(f(\{xX\}f))),\\
 \sigma s&:=\{s\mathcal U\}
      (\lambda t:\wp\wp\mathcal U.\tau t).
\end{aligned}
\]
The paper does not postulate sigma/tau constants or their desired equality:
each displayed use expands this actual expression. This \(\mathcal U\)
is the double-predicate-space universe of section~3, distinct from the
single-predicate-space variant and its longer proof in section~5.4.
\end{hbdefinition}

\begin{hblemma}{thm:p12-har-hurkens-beta}{proved=annex:p12.har.hurkens-beta}{Every data type and the computed identity}
In the exact source hierarchy,
\(\mathcal U:\Box\), \(\tau t:\mathcal U\), and
\(\sigma s:\wp\wp\mathcal U\). For \(p:\wp\mathcal U\),
define \(q_p:=\lambda y:\mathcal U.p(\tau\sigma y)\).
Then \(q_p:\wp\mathcal U\), and actual beta contractions give
\[
 (\sigma(\tau t))p=_\beta t q_p.
\]
No extensionality or an axiom asserting this identity is used.
\end{hblemma}

\begin{hbdefinition}{def:p12-har-hurkens-proof}{The full R0, M0, L0 witness}
Define
\[
\begin{aligned}
 I(p)&:=\Pi x:\mathcal U.(\sigma x\,p\Rightarrow px),\\
 K(x)&:=\Pi p:\wp\mathcal U.
               (\sigma x\,p\Rightarrow p(\tau\sigma x)),\\
 \Delta&:=\lambda x:\mathcal U.\neg K(x),\\
 \Omega&:=\mathrm{nf}_\beta
                (\tau(\lambda p:\wp\mathcal U.I(p))),\\
 W&:=\Pi p:\wp\mathcal U.(I(p)\Rightarrow p\Omega).
\end{aligned}
\]
The indicated finite normal form is explicit, not an appeal to global
normalization of this inconsistent system. If
\(q_{X,f,p}:=\lambda y:\mathcal U.p(f(\{yX\}f))\), it is
\[
 \Omega=\Lambda X:\Box.\lambda f:\wp\wp X\to X.
 \lambda p:\wp X.
 \Pi x:\mathcal U.
    (\sigma x\,q_{X,f,p}\Rightarrow p(f(\{xX\}f))).
\]
The three source proof terms, written with ordinary applications and
the original numeric assumption labels, are
\[
\begin{aligned}
 R_0&:=\lambda p:\wp\mathcal U.\lambda1:I(p).
            (1\Omega)(\lambda x:\mathcal U.1(\tau\sigma x)),\\
 M_0&:=\lambda x:\mathcal U.\lambda2:\sigma x\Delta.
       \lambda3:K(x).
       ((3\Delta)2)(\lambda p:\wp\mathcal U.3q_p),\\
 L_0&:=\lambda0:W.
       ((0\Delta)M_0)(\lambda p:\wp\mathcal U.0q_p),\\
 D_0&:=L_0R_0.
\end{aligned}
\]
These are the section~3 witness, with its section~7 R0/M0/L0 factoring.
The annex checks every application and all conversion steps; naming
well-foundedness or diagonalization is not substituted for these terms.
\end{hbdefinition}

\begin{hbtheorem}{thm:p12-har-hurkens-bottom}{proved=annex:p12.har.hurkens-bottom}{The closed source term of false type}
In exact \(\lambda U^-\),
\(R_0:W\), \(M_0:I(\Delta)\), \(L_0:W\Rightarrow\bot\),
and \(\vdash D_0:\bot\). Every displayed domain, predicate and
conversion follows from the stated axioms/product pairs and the computed
sigma/tau identity. The proof contains no assumption already asserting
falsehood or inhabitation of every type
\cite[sections~3 and 7, pp.~269,276--277]{Hurkens1995LambdaStar}.
\end{hbtheorem}

\begin{hbtheorem}{thm:p12-har-star-inconsistency}{proved=annex:p12.har.star-inconsistency}{Actual transport to lambda-star and all well-formed types}
Map each of \(\ast,\Box,\triangle\) to the single \(\ast\),
fix variables, and preserve product, abstraction and application.
This is a PTS morphism into exact lambda-star; let \(D_\ast\)
be the resulting fully expanded image of \(D_0\). Then
\(\vdash D_\ast:\Pi p:\ast.p\).
For every well-formed context \(\Gamma\) and actual type
\(\Gamma\vdash A:\ast\), the explicit term
\(D_\ast A\) inhabits \(A\) in \(\Gamma\).
This concerns derivable types in this exact self-sort calculus, not every
raw expression or every universe discipline.
\end{hbtheorem}

\begin{hbtheorem}{thm:p12-har-star-no-normal-form}{proved=annex:p12.har.star-no-normal-form}{The explicit witness has no beta-normal form}
The closed typed \(D_\ast\) above has no beta-normal form under
full compatible annotated beta reduction. Therefore exact lambda-star
is not even weakly normalizing on all typed terms, and in particular
is not strongly normalizing. The argument supplies the needed typed
preservation and normal-head obstruction; it does not infer this merely
from the word ``inconsistent.''
\end{hbtheorem}


\subsection{Products, sums, and the source strong-sum variants}

\begin{hbdefinition}{def:p12-har-product-sum-extensions}{Source Figures 4--5 and their polymorphic realizations}
The source desired product and sum rules are on the registered sheet.
Primitive \(A_1\times A_2\) has pairing/projections; primitive
\(A_1+A_2\) has two injections and a case operator. In a polymorphic
sort the chapter realizes these operations by encodings. For source
lambdaHOL this construction is available at \(\mathrm{Prop}\), not
at arbitrary data sort \(\mathrm{Type}\):
\[
\begin{aligned}
 A_1\times A_2&:=\Pi\alpha:\mathrm{Prop}.
                  (A_1\to A_2\to\alpha)\to\alpha,\\
 A_1+A_2&:=\Pi\alpha:\mathrm{Prop}.
                  (A_1\to\alpha)\to(A_2\to\alpha)\to\alpha,\\
 \mathrm{in}_1&:=\lambda p:A_1.\lambda\alpha:\mathrm{Prop}.
     \lambda h_1:A_1\to\alpha.\lambda h_2:A_2\to\alpha.h_1p,\\
 \mathrm{in}_2&:=\lambda p:A_2.\lambda\alpha:\mathrm{Prop}.
     \lambda h_1:A_1\to\alpha.\lambda h_2:A_2\to\alpha.h_2p,\\
 \mathrm{case}(f_1,f_2)&:=\lambda x:(A_1+A_2).x C f_1f_2.
\end{aligned}
\]
Here \(A_1,A_2,C:\mathrm{Prop}\), \(f_i:A_i\to C\).
The product projections and pairing are the full terms of
\cref{def:p12-har-encoded-proof-terms}, with
\(\varphi=A_1,\psi=A_2\). Expanding those abbreviations gives
exactly the source Figure~4 operations, not a new primitive data product.
The plus sign is a sum-type former here, not arithmetic addition.
Plain PTS encodings do not supply an induction principle for their
encoded data merely because a recursion term can be written
\cite[pp.~1202--1204, Figs.~4--5]{BarendregtGeuvers2001}.
\end{hbdefinition}

\begin{hbtheorem}{thm:p12-har-sum-checks}{proved=annex:p12.har.sum-checks}{The source encoded sum computations}
At the stated proposition sort,
\(\mathrm{in}_i:A_i\to(A_1+A_2)\) and
\(\mathrm{case}(f_1,f_2):(A_1+A_2)\to C\).
For \(p:A_i\), expanding abbreviations and beta-reducing gives
\(\mathrm{case}(f_1,f_2)(\mathrm{in}_i p)
\twoheadrightarrow_\beta f_i p\). The corresponding product
projection computations are supplied in
\cref{thm:p12-har-connective-checks}; neither claim adds data-level
polymorphism to the source lambdaHOL specification.
\end{hbtheorem}

\begin{hbdefinition}{def:p12-har-strong-sums}{The actual Figure 6 strong sum and its warned variants}
The displayed source strong sum has a data domain and propositional
fiber: from \(A:\mathrm{Type}\) and
\(x:A\vdash\varphi:\mathrm{Prop}\), form
\(\Sigma x:A.\varphi:\mathrm{Type}\), with pairing and
both projections. The second projection has the dependent type
\(\varphi[\pi_1t/x]\). Its two primitive projection computations
are included on the sheet; they are not beta contractions of an encoded
existential. The extended conversion relation must admit these declared
projection computations to check dependent projections.

The chapter considers two different changes. A sum over
\(A:\mathrm{Type}'\) with Type-valued fiber is displayed with
a Type-valued result, then explicitly warned against: in particular
\(\Sigma x:\mathrm{Type}.B\) must have the higher result sort
\(\mathrm{Type}'\), not \(\mathrm{Type}\), in its proposed
universe discipline. Separately, giving the strong sum with Type domain
and Prop fiber a Prop result is proposed as an existential interpretation,
but the source states this addition is inconsistent for its lambdaHOL.
Both warned rule surfaces are recorded as separate variants on the sheet,
not as baseline rules or as proved consistency claims. Their paradox and
consistency qualifications require their own complete metatheoretic
arguments; they are not supplied by the ordinary beta-only theorem
\cite[pp.~1204--1206, Fig.~6]{BarendregtGeuvers2001}.
\end{hbdefinition}

\begin{hbexample}{def:p12-har-group-record}{The full source group-record construction and its formation boundary}
For \(A:\mathrm{Type}\), write the operation \(\circ:A\to A\to A\),
unit \(e:A\), and inverse \(\mathrm{inv}:A\to A\). The three
listed source axioms are
\[
\begin{aligned}
 \mathrm{Assoc}&:=\Pi x,y,z:A.(x\circ y)\circ z=_A x\circ(y\circ z),\\
 \mathrm{Unit}&:=\Pi x:A.e\circ x=_A x,\\
 \mathrm{Inverse}&:=\Pi x:A.(\mathrm{inv}\,x)\circ x=_A e.
\end{aligned}
\]
Its group-structure type is
\(\mathrm{Group\!\mbox{-}\!Str}(A)=
(A\to A\to A)\times(A\times(A\to A))\).
Writing \(\Phi=\mathrm{Assoc}\land
(\mathrm{Unit}\land\mathrm{Inverse})\), its intended full record is
\[
 \mathrm{Group}(A):=
 \Sigma\circ:A\to A\to A.\Sigma e:A.
       \Sigma\mathrm{inv}:A\to A.\Phi.
\]
The outer two sums have data-valued fibers. Thus this nested expression
needs data/data strong-sum formation in addition to the data/Prop rule
actually displayed in Figure~6; that figure alone does not type it.
When these stated record formations are admitted, witnesses
\(f,a,h,p_1,p_2,p_3\) give the record
\[
 \langle f,\langle a,\langle h,
       \langle p_1,\langle p_2,p_3\rangle\rangle
       \rangle\rangle\rangle.
\]
There are three axiom proofs, matching the three listed axioms, not a
fourth unspecified proof field. For \(t:\mathrm{Group}(A)\), its
operation, unit and inverse are \(\pi_1t\),
\(\pi_1(\pi_2t)\), and \(\pi_1(\pi_2(\pi_2t))\).
The axiom package is \(\pi_2(\pi_2(\pi_2t))\); its first
projection gives associativity for that extracted operation.
Abstracting also over the carrier gives the intended
\(\mathrm{Group}:=\Sigma A:\mathrm{Type}.\mathrm{Group}(A)\),
at the higher result sort specified by the source's universe warning.
This is a record schema with explicit formation requirements, not a
derivable baseline lambdaHOL type
\cite[pp.~1204--1205]{BarendregtGeuvers2001}.
\end{hbexample}

\subsection{The source inductive mechanisms and their worked computations}

\begin{hbdefinition}{def:p12-har-inductive-scheme}{The chapter's strictly positive constructor scheme}
The source inductive declaration is
\[
\begin{aligned}
 &\mathrm{Inductive}\ \mu:\mathrm{Type}:=\\
 &\quad\mathrm{constr}_i:
    \sigma_1^i(\mu)\to\cdots\to\sigma_{m_i}^i(\mu)\to\mu
          \quad(1\le i\le n).
\end{aligned}
\]
Each argument scheme is \(A_1\to\cdots\to A_k\to X\),
with \(\mu\) absent from every \(A_j\), and either
\(X=\mu\) or \(\mu\) absent from \(X\). This is the
chapter's stated strictly positive regime; a negative occurrence such as
\(\mu\to D\) in a constructor argument does not satisfy it.
The declaration generates constructors, data recursion, proposition-valued
induction and constructor computation. These additions change conversion
to admit their \(\iota\)-steps; the beta-only structural proofs above
do not automatically become beta/iota proofs for this extension.
The chapter sets aside the alternative W-type route in its intensional
presentation and refers to a separate extensional treatment. Its
generated constructor recursors are not silently identified with that
other presentation.

For an argument ending in \(\mu\), data recursion supplies a recursive
value family of type \(A_1\to\cdots\to A_k\to C\).
Induction for \(P:\mu\to\mathrm{Prop}\) supplies the family
\(\Pi u_1:A_1\cdots u_k:A_k.P(a_j u_1\cdots u_k)\).
The branch receives every original constructor argument and then these
recursive families in their constructor-argument order. At a constructor,
its recursive family is built by
\(\lambda u_1:A_1\cdots u_k:A_k.
\mathrm{Rec}_\mu\vec f(a_j u_1\cdots u_k)\).
An argument not containing \(\mu\) gets no recursive family. Write
\(\mathcal B_i(C)\) for the type quantifying over all those original
arguments and recursive value families and returning \(C\);
\(\mathcal B_i(P)\) quantifies over the original arguments and
their indexed induction hypotheses and returns
\(P(\mathrm{constr}_i\vec a)\). These metalevel branch-type names
make the generated rule construction explicit; they are not new object
type formers. These
clauses give the actual construction behind the nat, List and Tree rules,
rather than merely naming a generated eliminator
\cite[pp.~1206--1209]{BarendregtGeuvers2001}.
\end{hbdefinition}

\begin{hbdefinition}{def:p12-har-nat-list-tree}{The three fixed inductive types}
The source declarations, each with its complete recursion/induction and
computation rules on the sheet, are
\[
\begin{aligned}
 \mathrm{nat}&:\mathrm{Type},&
 0&:\mathrm{nat},&S&:\mathrm{nat}\to\mathrm{nat},\\
 \mathrm{List}&:\mathrm{Type},&
 \mathrm{Nil}&:\mathrm{List},&
 \mathrm{Cons}&:D\to\mathrm{List}\to\mathrm{List},\\
 \mathrm{Tree}&:\mathrm{Type},&
 \mathrm{Leaf}&:D\to\mathrm{Tree},&
 \mathrm{Join}&:D\to(\mathrm{nat}\to\mathrm{Tree})\to\mathrm{Tree}.
\end{aligned}
\]
Here \(D:\mathrm{Type}\) is fixed for List and Tree. Tree is
constructor-generated and well founded, but every Join has countably many
children; it is not a datatype of infinite non-well-founded branches.
Its data branch receives
\(d:D,t:\mathrm{nat}\to\mathrm{Tree},
r:\mathrm{nat}\to C\). Its induction branch instead receives
\(\Pi n:\mathrm{nat}.P(tn)\). Replacing that family by one
unindexed induction hypothesis would change the source rule.
The recursor notation is schematic at the selected target type/predicate;
no forbidden polymorphic type \(\Pi C:\mathrm{Type}.\cdots\)
is silently assigned to it inside the base calculus
\cite[pp.~1206--1209]{BarendregtGeuvers2001}.
\end{hbdefinition}

\begin{hbexample}{def:p12-har-double-map}{Double, evenness, and the full List map term}
The source terms are
\[
\begin{aligned}
 d&:=\mathrm{Rec}_{\mathrm{nat}}\,0
       (\lambda x:\mathrm{nat}.\lambda y:\mathrm{nat}.S(Sy)),\\
 \mathrm{even}&:=\mathrm{Rec}_{\mathrm{nat}}\,\top
       (\lambda x:\mathrm{nat}.\lambda\alpha:\mathrm{Prop}.\neg\alpha),\\
 \mathrm{map}&:=\lambda f:D\to D.\lambda l:\mathrm{List}.
       \mathrm{Rec}_{\mathrm{List}}\,\mathrm{Nil}
       (\lambda a:D.\lambda k:\mathrm{List}.
          \lambda h:\mathrm{List}.\mathrm{Cons}(fa)h)\,l.
\end{aligned}
\]
The final \(l\) is the actual argument to the recursor. It is needed
for the displayed \((D\to D)\to\mathrm{List}\to\mathrm{List}\)
type; leaving it out under an extra \(\lambda l\) would return a
function instead of the asserted list. Abbreviations expand before
beta/iota computation. Thus \(d0\twoheadrightarrow0\),
\(d(Sx)\twoheadrightarrow S(S(dx))\),
\(\mathrm{even}(0)\twoheadrightarrow\top\), and
\(\mathrm{even}(Sx)\twoheadrightarrow\neg\mathrm{even}(x)\).
With \(\mathrm{true}:\top\), the source induction term is
\[
 \mathrm{Rec}_{\mathrm{nat}}\,\mathrm{true}
 (\lambda x:\mathrm{nat}.\lambda h:\mathrm{even}(dx).
    \lambda z:\neg\mathrm{even}(dx).zh).
\]
It proves \(\Pi x:\mathrm{nat}.\mathrm{even}(dx)\) using
double-negation introduction, not double-negation elimination
\cite[pp.~1207--1208]{BarendregtGeuvers2001}.
\end{hbexample}

\begin{hbtheorem}{thm:p12-har-inductive-checks}{proved=annex:p12.har.inductive-checks}{The actual constructor branches and worked induction term}
For the stated nat/List/Tree recursor surfaces, each displayed constructor
\(\iota\)-contractum has the same prescribed target type as its redex,
for data recursion and for proposition-valued induction. The generated
recursive families have the explicit types in
\cref{def:p12-har-inductive-scheme}. The evenness term has type
\(\Pi x:\mathrm{nat}.\mathrm{even}(dx)\), and
\(\mathrm{map}\,f(\mathrm{Cons}\,a t)
\twoheadrightarrow_{\beta\iota}
\mathrm{Cons}(fa)(\mathrm{map}\,f t)\), with
\(\mathrm{map}\,f\,\mathrm{Nil}\twoheadrightarrow\mathrm{Nil}\).
These are branch typing/computation results for the declared extension,
not a global consistency or beta/iota normalization theorem.
\end{hbtheorem}

\begin{hbexample}{def:p12-har-tree-operations}{The two source Tree tasks made explicit}
For a fixed \(a:D\) and \(n:\mathrm{nat}\), the nth-subtree
operation is the data recursion with branches
\(\lambda d:D.\mathrm{Leaf}\,a\) and
\(\lambda d:D.\lambda t:\mathrm{nat}\to\mathrm{Tree}.
\lambda r:\mathrm{nat}\to\mathrm{Tree}.tn\).
It returns \(tn\) on \(\mathrm{Join}\,d t\) and the stated
default \(\mathrm{Leaf}\,a\) on a leaf. To distinguish a Join
from a single leaf, use target \(\mathrm{Prop}:\mathrm{Type}\),
leaf branch \(\lambda d:D.\bot\), and join branch
\(\lambda d:D.\lambda t:\mathrm{nat}\to\mathrm{Tree}.
\lambda r:\mathrm{nat}\to\mathrm{Prop}.\top\).
The output computes to \(\bot\) or \(\top\) on those
constructors. In the source's phrase ``infinite or just a single leaf,''
infinite means infinitely many nodes supplied by a Join, not an infinite
descending branch in this well-founded inductive type
\cite[p.~1209]{BarendregtGeuvers2001}.
\end{hbexample}

\begin{hbdefinition}{def:p12-har-inductive-propositions}{Inductive conjunction and equality}
For \(\varphi,\psi:\mathrm{Prop}\), the source primitive
inductive conjunction has
\(\mathrm{Pair}:\varphi\to\psi\to(\varphi\land\psi)\).
Its recursor eliminates only to a fixed \(P:\mathrm{Prop}\),
with branch \(\varphi\to\psi\to P\). A family depending
on a conjunction proof would require the absent Prop/Type pair to form
its type. Taking the branches \(\lambda x:\varphi.\lambda y:\psi.x\)
or \(\lambda x:\varphi.\lambda y:\psi.y\) gives projections.
This is a primitive inductive presentation, distinct from the previous
impredicative conjunction encoding.

For \(D:\mathrm{Type}\), source equality is the inductive predicate
\(\mathrm{Eq}:D\to D\to\mathrm{Prop}\), with
\(\mathrm{Refl}:\Pi x:D.\mathrm{Eq}\,x x\).
Its eliminator uses \(P:D\to D\to\mathrm{Prop}\) and
\(f_1:\Pi x:D.Pxx\) to prove
\(\Pi x,y:D.(\mathrm{Eq}\,xy)\to Pxy\).
The source computation display makes its two indices explicit as
\(\mathrm{Rec}_{\mathrm{Eq}}\,x x f_1(\mathrm{Refl}\,x)
\to_\iota f_1x\). Its conclusion predicate depends on the endpoints,
not on their equality proof. This is not the full proof-dependent identity
eliminator of another type theory
\cite[pp.~1209--1211]{BarendregtGeuvers2001}.
\end{hbdefinition}

\begin{hbdefinition}{def:p12-har-dependent-inductives}{The proposed parametric and dependent-constructor extensions}
The source parametric-list scheme is
\(\mathrm{List}:\mathrm{Type}\to\mathrm{Type}\),
\(\mathrm{Nil}:\Pi D:\mathrm{Type}.\mathrm{List}\,D\), and
\(\mathrm{Cons}:\Pi D:\mathrm{Type}.
D\to\mathrm{List}\,D\to\mathrm{List}\,D\).
For each fixed \(D\), the sheet retains its exact recursion/induction
and constructor computations. The constructor quantification requires
the separately proposed Type'/Type product pair; the base system does
not even form the displayed Type-to-Type family kind. The source warns
that its proposed enlargement of lambdaHOL is inconsistent and instead
discusses separately permitting type parameters in a PRED-style setting.
Neither additional product pairs nor a modern universe hierarchy are
silently inserted into the base specification.

In a different extension explicitly adding Prop/Type, the source uses
\(B:\mathrm{Type}\), \(Q:B\to\mathrm{Prop}\), and
\(\mathrm{In}:\Pi z:B.(Qz)\to\mu\), for
\(\mu:\mathrm{Type}\). This yields the strong-sum projections
\[
 \pi_1:=\mathrm{Rec}_\mu(\lambda z:B.\lambda y:Qz.z),
 \qquad
 \pi_2:=\mathrm{Rec}_\mu(\lambda z:B.\lambda y:Qz.y),
\]
where the second recursor uses predicate
\(\lambda x:\mu.Q(\pi_1x)\). Its branch type is checked by
the first projection computation. The predicate domain must be \(B\),
matching the constructor's \(z:B\); a declaration \(Q:A\to\mathrm{Prop}\)
with unrelated \(A\) does not type \(Qz\).
These projections make data depend on proof arguments, so the
disambiguation result for the unchanged base lambdaHOL does not apply
to this extension \cite[pp.~1210--1211]{BarendregtGeuvers2001}.
\end{hbdefinition}

\begin{hbtheorem}{thm:p12-har-inductive-projections}{proved=annex:p12.har.inductive-projections}{The dependent-constructor projections and nondependent predicate elimination}
In the stated Prop/Type extension,
\(\pi_1:\mu\to B\) and
\(\pi_2:\Pi x:\mu.Q(\pi_1x)\), with
\(\pi_1(\mathrm{In}\,b q)\twoheadrightarrow b\) and
\(\pi_2(\mathrm{In}\,b q)\twoheadrightarrow q\).
For the stated primitive conjunction and Eq rules, their constructor
recursor computations have the asserted proposition type; neither one
supplies an eliminator into arbitrary data or one depending on a proof.
\end{hbtheorem}

\begin{proofboundary}{def:p12-har-extension-metatheory}{What the extension arguments do and do not establish}
Constructor-root checks establish the declared computations' local
type preservation. On a supplied well-founded constructor tree, recursion
is constructed by assigning the already obtained child values to each
branch and then applying that branch; induction follows the same
well-founded argument with child proof families. A full formal extension
still needs its exact admissible-signature and formation rules, global
beta/iota subject reduction/confluence/normalization and the stated
universe-consistency or inconsistency arguments. In particular a positive
constructor name does not by itself prove those results, and the partial
group-record formation or warned parametric/sum variants cannot be
discharged by the beta-only base metatheory.
\end{proofboundary}

\subsection{The historical proof-development workflow}

\begin{hbdefinition}{def:p12-har-tactic-workflow}{Source tactics and their term-construction meaning}
The chapter's Figure~7 separates a current context/goal, the
proof-development system receiving tactics, its constructed proof object,
and the proof checker producing a certified statement. Its Coq examples
use historical syntax: \texttt{(x:A)B} denotes \(\Pi x:A.B\),
and \texttt{[x:A]M} denotes \(\lambda x:A.M\). In that input
syntax abstraction and \texttt{->} bind more strongly than application,
so parentheses around \texttt{(P x)} are substantive. This syntax and
the source's version~6.3 theory development are not current Rocq scripts
or a direct implementation of the three-sort beta-only lambdaHOL checker.

For goal \(\Pi x:A.Px\to Qx\), \texttt{Intros} introduces
fresh \(x:A,H:Px\) and leaves goal \(Qx\). A completed term
\(q:Qx\) produces \(\lambda x:A.\lambda H:Px.q\).
For goal \(Qa\) and
\(p:\Pi x:A.Px\to Qx\), \texttt{Apply p} matches the
conclusion with \(Qa\), instantiates \(x=a\), and leaves
the proof hole \(pa\,?\) with goal \(Pa\). Matching alone
does not complete that missing proof.

For \(n:\mathrm{nat}\), \texttt{Elim n} at goal \(Qn\)
creates \(Q0\) and \(\Pi k:\mathrm{nat}.Qk\to Q(Sk)\);
the corresponding term is \(\mathrm{nat\_ind}\,Q f_0 f_S n\).
The source \(\mathrm{nat\_rec}\) separately computes at zero
and successor by its declared constructor reductions. Once every goal is
filled and the term is checked, \texttt{Save fst\_lemma} records its
definition under that name. A transcript or success message without the
completed checking bridge is not itself a proof
\cite[pp.~1211--1214, Fig.~7]{BarendregtGeuvers2001}.
\end{hbdefinition}

\begin{hbdefinition}{def:p12-har-prime-development}{The historical theory-development entry}
The chapter's version~6.3 development uses
\(\mathsf{nat}:\mathsf{Set}\), with source constructor O and successor S,
and inductive conjunction, truth \(\mathsf{TT}\), falsehood
\(\mathsf{FF}\), disjunction and existence in \(\mathsf{Prop}\).
Its \(\mathsf{not}(a)=a\to\mathsf{FF}\) and
\(\mathsf{iff}(a,b)=(a\to b)\land(b\to a)\).
For \(A:\mathsf{Set},P:A\to\mathsf{Prop}\), the existential constructor is
\(\mathsf{ex\_intro}:\Pi x:A.Px\to\mathsf{ex}\,A P\).
The source explicitly declares
\(\mathsf{DN}:\Pi a:\mathsf{Prop}.\neg\neg a\to a\);
that classical assumption remains part of this proof environment.

The full constructor/equality/elimination regime is specified in
\cref{def:p12-har-prime-regime}. Its natural-number constants,
\(\mathsf{Is\_suc}\) cases, inductive leseq/lthan/before definitions, both recursive
arithmetic functions and coefficient-left
div/propdiv/prime/primediv/HPD definitions are given in
\cref{def:p12-har-prime-arithmetic}, together with the additional explicit
subtraction and factorial recursions required for the Euclid continuation.
They retain the historical identifiers and argument order.
The source's abstract \(\mathsf{nat}:\mathsf{Type}\) presentation from
section~3.8, the beta-only lambdaHOL proof-checker, and this Set-based
classical library are different regimes, not aliases for the same rule set.

The chapter itself omits most objects for its listed logical/arithmetic
library and prime-divisor end theorem. Here the corresponding constructions
are supplied in
\cref{thm:p12-har-prime-logical,thm:p12-har-prime-equality,thm:p12-har-prime-order,thm:p12-har-prime-semiring,thm:p12-har-prime-divisibility,thm:p12-har-prime-divisor,thm:p12-har-prime-euclid}.
In particular \(\mathsf{cv\_ind}\) is derived from ordinary \(\mathsf{nat\_ind}\), numerical
definitions are not treated as their arithmetic theorems, and DN is
visible in the prime-divisor proof. The two objects actually printed
by the source remain separately attributed in
\cref{thm:p12-har-development-objects}
\cite[pp.~1215--1220]{BarendregtGeuvers2001}.
\end{hbdefinition}

\subsection{The historical arithmetic library and Euclid development}

\begin{hbdefinition}{def:p12-har-prime-regime}{The exact historical elimination boundary}
The following development reconstructs the chapter's Coq~6.3 example,
not a silently imported modern natural-number library. Its data declaration
is \(\mathsf{nat}:\mathsf{Set}\), with \(O:\mathsf{nat}\) and
\(S:\mathsf{nat}\to\mathsf{nat}\). The licensed operations are
structural recursion and case analysis on this \emph{data} declaration,
and the generated induction principle
\[
 \mathsf{nat\_ind}:\Pi P:\mathsf{nat}\to\mathsf{Prop}.
 P O\to(\Pi n:\mathsf{nat}.P n\to P(S n))\to
 \Pi n:\mathsf{nat}.P n.
\]
The historical logical declarations are likewise retained:
\(\mathsf{and}\ A\ B:\mathsf{Prop}\) has constructor
\(\mathsf{conj}:A\to B\to\mathsf{and}\ A\ B\);
\(\mathsf{TT}:\mathsf{Prop}\) has \(\mathsf{trivial}:\mathsf{TT}\);
\(\mathsf{FF}:\mathsf{Prop}\) has no constructors;
\(\mathsf{or}\ A\ B:\mathsf{Prop}\) has
\(\mathsf{or\_introl}:A\to\mathsf{or}\ A\ B\) and
\(\mathsf{or\_intror}:B\to\mathsf{or}\ A\ B\).
For \(A:\mathsf{Set}\) and \(P:A\to\mathsf{Prop}\),
\(\mathsf{ex}\ A\ P:\mathsf{Prop}\) has
\(\mathsf{ex\_intro}:\Pi x:A.P x\to\mathsf{ex}\ A\ P\).
Here \(\neg a:=a\to\mathsf{FF}\) and
\(a\leftrightarrow b:=(a\to b)\wedge(b\to a)\).
In particular, the displayed source \(\mathsf{Is\_suc}\) case analysis
returns a proposition code, whereas plus, times, predecessor, subtraction
and factorial return data. Induction on \(\mathsf{leseq}\) and elimination
of \(\mathsf{and},\mathsf{or},\mathsf{ex},\mathsf{FF}\) below have
\(\mathsf{Prop}\) targets only. For
\(\mathsf{eq}\ A\ x:\!A\to\mathsf{Prop}\), the constructor is
\(\mathsf{refl\_equal}:\mathsf{eq}\ A\ x\ x\), and its generated
\emph{propositional} eliminator transports \(P x\) along \(x=y\) to
\(P y\), for \(P:A\to\mathsf{Prop}\). No equality reflection,
proof irrelevance, choice, or elimination of an existential proof into
\(\mathsf{Set}\) is used. Finally, retain the explicit source axiom
\[
 \mathsf{DN}:\Pi a:\mathsf{Prop}.((a\to\mathsf{FF})\to\mathsf{FF})\to a.
\]
These are the historical inductive/recursive extensions of the example,
not assertions about the beta-only \(\lambda\mathsf{HOL}\) calculus.
The chapter's historical \((x:A)B\) notation denotes a dependent
product and \([x:A]m\) denotes annotated abstraction; the formulas
below use \(\Pi\) and \(\lambda\) for those same constructions.
Section~3.8's abstract \(\mathsf{nat}:\mathsf{Type}\) and its separately
registered data/proof \(\mathsf{Rec}_{\mathsf{nat}}\) schemas are not
identified with this historical \(\mathsf{Set}\) spelling. The following
proofs are relative to the stated generated principles; a general
consistency theorem for all universe/inductive extensions is a separate
metatheoretic obligation \cite[pp.~1214--1220]{BarendregtGeuvers2001}.
\end{hbdefinition}

\begin{hbdefinition}{def:p12-har-prime-arithmetic}{Source definitions and the needed finite data recursions}
Use the source names \(\mathsf{one}:=S O\), \(\mathsf{two}:=S\mathsf{one}\),
and \(\mathsf{Is\_suc}(O):=\mathsf{FF}\),
\(\mathsf{Is\_suc}(S n):=\mathsf{TT}\). The source equations are
\[
\begin{aligned}
 \mathsf{plus}\ O\ m&=m,& \mathsf{plus}\ (S n)\ m&=S(\mathsf{plus}\ n\ m),\\
 \mathsf{times}\ O\ m&=O,& \mathsf{times}\ (S n)\ m&=
               \mathsf{plus}\ (\mathsf{times}\ n\ m)\ m.
\end{aligned}
\]
Write \(n+m:=\mathsf{plus}\ n\ m\), \(nm:=\mathsf{times}\ n\ m\),
\(n\le m:=\mathsf{leseq}\ n\ m\), and
\(n<m:=\mathsf{leseq}\ (S n)\ m\). The constructors of
\(\mathsf{leseq}\ n\) are \(\mathsf{leseq\_n}:n\le n\) and
\(\mathsf{leseq\_suc}:\Pi m.(n\le m)\to(n\le S m)\).
The instantiated propositional induction schema used for the ordering is
\[
 \Pi n:\mathsf{nat}.\Pi P:\mathsf{nat}\to\mathsf{Prop}.
 P n\to(\Pi m:\mathsf{nat}.(n\le m)\to P m\to P(S m))\to
 \Pi m:\mathsf{nat}.(n\le m)\to P m.
\]
It is generated from these two constructors; it neither changes the
fixed lower bound n nor eliminates into data.
Also \(\mathsf{before}\ n\ P:=\Pi k.(k<n)\to P k\).
The source divisor convention is crucial:
\[
\begin{aligned}
 \mathsf{div}\ d\ n&:=\exists x:\mathsf{nat}.\;xd=n,\\
 \mathsf{propdiv}\ d\ n&:=(\mathsf{one}<d)\wedge(d<n)\wedge\mathsf{div}\ d\ n,\\
 \mathsf{prime}\ n&:=(\mathsf{one}<n)\wedge
                    \neg\exists d.\mathsf{propdiv}\ d\ n,\\
 \mathsf{primediv}\ p\ n&:=\mathsf{prime}\ p\wedge\mathsf{div}\ p\ n,\\
 \mathsf{HPD}\ n&:=\exists p.\mathsf{primediv}\ p\ n.
\end{aligned}
\]
Here \(\exists\) is the source \(\mathsf{ex}:\mathsf{Prop}\), not a
strong sum. Needed data recursions, supplied rather than assumed as
theorems, are
\[
\begin{aligned}
 \mathsf{pred}\ O&:=O,&\mathsf{pred}(S n)&:=n,\\
 O\mathbin{\dot{-}}b&:=O,&(S a)\mathbin{\dot{-}}O&:=S a,&
 (S a)\mathbin{\dot{-}}(S b)&:=a\mathbin{\dot{-}}b,\\
 \mathsf{fac}\ O&:=\mathsf{one},&\mathsf{fac}(S n)&:=(S n)(\mathsf{fac}\ n).
\end{aligned}
\]
Subtraction is structural recursion on the first input, returning a
function of the second; factorial is structural recursion on its input.
All displayed equations are computation equations, not arithmetic
lemmas. More explicitly, the historical \(\mathsf{nat\_rec}\) schema
is instantiated at \(C=\mathsf{nat}\to\mathsf{nat}\) for
\[
\begin{aligned}
 \mathsf{plus}&:=\mathsf{nat\_rec}
  (\lambda m:\mathsf{nat}.m)
  (\lambda n:\mathsf{nat}.\lambda r:\mathsf{nat}\to\mathsf{nat}.
                    \lambda m:\mathsf{nat}.S(r m)),\\
 \mathsf{times}&:=\mathsf{nat\_rec}
  (\lambda m:\mathsf{nat}.O)
  (\lambda n:\mathsf{nat}.\lambda r:\mathsf{nat}\to\mathsf{nat}.
                    \lambda m:\mathsf{nat}.\mathsf{plus}(r m)m),\\
 \mathsf{sub}&:=\mathsf{nat\_rec}
  (\lambda b:\mathsf{nat}.O)
  (\lambda a:\mathsf{nat}.\lambda r:\mathsf{nat}\to\mathsf{nat}.
    \lambda b:\mathsf{nat}.
     \mathsf{nat\_rec}(S a)
       (\lambda c:\mathsf{nat}.\lambda z:\mathsf{nat}.r c)\ b).
\end{aligned}
\]
Here \(a\mathbin{\dot{-}}b:=\mathsf{sub}\ a\ b\), and the inner
sub recursor has \(C=\mathsf{nat}\) and ignores its recursive z.
At that same data target,
\(\mathsf{pred}:=\mathsf{nat\_rec}\ O\
(\lambda n:\mathsf{nat}.\lambda r:\mathsf{nat}.n)\)
and \(\mathsf{fac}:=\mathsf{nat\_rec}\ \mathsf{one}\
(\lambda n:\mathsf{nat}.\lambda r:\mathsf{nat}.\mathsf{times}(S n)r)\).
The base and step of each instance have respectively types C and
\(\mathsf{nat}\to C\to C\); their constructor computation gives
exactly the equations above. This uses several instantiated recursion
schemas, not a new base-calculus constant polymorphic over a top sort.
The source's repeated \((n,m,k:\mathsf{nat})\) binder groups in
\(\mathsf{plus\_ass}\) and \(\mathsf{distr}\) are reconstructed below
with one binder group: no hidden extra parameters are required
\cite[pp.~1219--1220 and p.~1215, note~4]{BarendregtGeuvers2001}.
If the repeated names are instead read as literal binder shadowing,
alpha-rename the outer triple to \(n_0,m_0,k_0\). The resulting
six-argument type has those three arguments unused; its witness is
\(\lambda n_0:\mathsf{nat}.\lambda m_0:\mathsf{nat}.
\lambda k_0:\mathsf{nat}.L\), where L is the fully quantified
single-triple proof of \(\mathsf{plus\_ass}\) or distr constructed below. Thus neither
reading requires an additional arithmetic lemma; no claim about the
historical parser accepting that typography is needed.
\end{hbdefinition}

\begin{hblemma}{thm:p12-har-prime-logical}{proved=annex:p12.har.prime-logical}{The source logical library, with classical dependence separated}
Let \(a,b,c:\mathsf{Prop}\), \(A:\mathsf{Set}\), and
\(P,Q:A\to\mathsf{Prop}\); all these parameters are universally
quantified in the following signatures.
The source constructor/eliminator aliases
\(\mathsf{and\_in},\mathsf{and\_ell},\mathsf{and\_elr},
\mathsf{false\_el},\mathsf{or\_inl},\mathsf{or\_inr},\mathsf{or\_el}\)
have the generated propositional introduction/projection/elimination
types. Also \(\mathsf{all\_el}:\Pi x:A.(\Pi y:A.P y)\to P x\)
and \(\mathsf{ex\_in}:\Pi x:A.P x\to\mathsf{ex}\ A\ P\)
are constructed directly, as is
\(\mathsf{non\_or}:\neg(a\vee b)\to\neg a\wedge\neg b\).
The source \(\mathsf{non\_ex},\mathsf{all\_not},\mathsf{all\_and},
\mathsf{ex\_or}\) have the following types:
\[
\begin{aligned}
 \mathsf{non\_ex}&:\neg(\exists x.P x)\to\Pi x.\neg P x,\\
 \mathsf{all\_not}&:(\Pi x.\neg P x)\to\neg\exists x.P x,\\
 \mathsf{all\_and}&:(\Pi x.P x\wedge Q x)\to
                         (\Pi x.P x)\wedge(\Pi x.Q x),\\
 \mathsf{ex\_or}&:(\exists x.P x\vee Q x)\to
                         (\exists x.P x)\vee(\exists x.Q x).
\end{aligned}
\]
Each bound x has domain A. The source
\(\mathsf{dn\_c}:\neg\neg a\leftrightarrow a\),
\(\mathsf{tnd}:a\vee\neg a\),
\(\mathsf{non\_all}:\neg(\forall x.P x)\to\exists x.\neg P x\),
and \(\mathsf{ex\_c}:(\exists x.P x)\leftrightarrow\neg\forall x.\neg P x\)
are constructed using the displayed \(\mathsf{DN}\). The first four
displayed predicate-logic lemmas and the introduction/elimination aliases
need no classical axiom and require no inhabitant of \(A\).
\end{hblemma}

\begin{hblemma}{thm:p12-har-prime-equality}{proved=annex:p12.har.prime-equality}{Equality transport and natural-number separation}
The source \(\mathsf{sym\_eq},\mathsf{leib},\mathsf{eq\_ind\_r},
\mathsf{f\_equal},\mathsf{no\_conf}\) are derivable from the stated
propositional equality eliminator, at these exact signatures:
\[
\begin{aligned}
 \mathsf{sym\_eq}&:\Pi A:\mathsf{Set}.\Pi x,y:A.(x=y)\to(y=x),\\
 \mathsf{leib}&:\Pi A:\mathsf{Set}.\Pi P:A\to\mathsf{Prop}.
                   \Pi x,y:A.(x=y)\to P x\to P y,\\
 \mathsf{eq\_ind\_r}&:\Pi A:\mathsf{Set}.\Pi x:A.
    \Pi P:A\to\mathsf{Prop}.P x\to\Pi y:A.(y=x)\to P y,\\
 \mathsf{f\_equal}&:\Pi A,B:\mathsf{Set}.\Pi f:A\to B.
                        \Pi x,y:A.(x=y)\to(f x=f y),\\
 \mathsf{no\_conf}&:\Pi n:\mathsf{nat}.\neg(O=S n).
\end{aligned}
\]
Moreover,
\(S n=S m\to n=m\). None uses \(\mathsf{DN}\).
\end{hblemma}

\begin{hblemma}{thm:p12-har-prime-order}{proved=annex:p12.har.prime-order}{The inductive ordering and course-of-values induction}
The source \(\mathsf{leseq\_trans},\mathsf{lthan\_leseq},
\mathsf{non\_lt0},\mathsf{suc\_leseq},\mathsf{lt01},
\mathsf{n0n1lt},\mathsf{cv\_ind}\) are derivable:
\[
\begin{aligned}
 \mathsf{leseq\_trans}&:\Pi x,y,z.(x\le y)\to(y\le z)\to(x\le z),\\
 \mathsf{lthan\_leseq}&:\Pi n,m.(n<m)\to(n\le m),\\
 \mathsf{non\_lt0}&:\Pi n.\neg(n<O),\\
 \mathsf{suc\_leseq}&:\Pi n,m.(S n\le S m)\to(n\le m),\\
 \mathsf{lt01}&:\Pi x.(x=O)\vee(x=\mathsf{one})\vee(\mathsf{one}<x),\\
 \mathsf{n0n1lt}&:\Pi n.\neg(n=O)\to\neg(n=\mathsf{one})\to(\mathsf{one}<n).
\end{aligned}
\]
The omitted natural-number domains are all \(\mathsf{nat}\). Useful derived
forms are zero minimality, trichotomy, antisymmetry, strict
irreflexivity, \(k<S n\leftrightarrow k\le n\), and
\(k\le S n\to(k\le n\vee k=S n)\). In particular, for every
\(P:\mathsf{nat}\to\mathsf{Prop}\),
\[
 (\Pi n.\mathsf{before}\ n\ P\to P n)\to\Pi n.P n
\]
is proved, not supplied as a new induction axiom.
\end{hblemma}

\begin{hblemma}{thm:p12-har-prime-semiring}{proved=annex:p12.har.prime-semiring}{The actual addition and multiplication library}
The source \(\mathsf{plus\_altsuc},\mathsf{plus\_altzero},
\mathsf{plus\_ass},\mathsf{plus\_com},\mathsf{distr},
\mathsf{timesaltzero},\mathsf{timesaltsuc},\mathsf{times\_ass}\)
are proved at the following single-binder types (all variables in nat):
\[
\begin{aligned}
 \mathsf{plus\_altsuc}&:\Pi n,m.n+S m=S(n+m),\\
 \mathsf{plus\_altzero}&:\Pi n.n+O=n,\\
 \mathsf{plus\_ass}&:\Pi n,m,k.n+(m+k)=(n+m)+k,\\
 \mathsf{plus\_com}&:\Pi n,m.n+m=m+n,\\
 \mathsf{distr}&:\Pi n,m,k.(n+m)k=nk+mk,\\
 \mathsf{timesaltzero}&:\Pi n.nO=O,\\
 \mathsf{timesaltsuc}&:\Pi n,m.n(S m)=nm+n,\\
 \mathsf{times\_ass}&:\Pi n,m,k.n(mk)=(nm)k.
\end{aligned}
\]
In addition,
\(nm=mn\), \(\mathsf{one}\,n=n\,\mathsf{one}=n\), both addition
cancellation laws, and
\[
 n\le m\quad\leftrightarrow\quad\exists k.\;n+k=m
\]
are derived. For \(O<d\), multiplication by \(d\) preserves strict
order and reflects weak order. These facts require no \(\mathsf{DN}\).
\end{hblemma}

\begin{hblemma}{thm:p12-har-prime-divisibility}{proved=annex:p12.har.prime-divisibility}{The subtraction dependency actually discharged}
Divisibility is reflexive and transitive, and \(\mathsf{one}\) divides
every natural number. A number \(d>\mathsf{one}\) cannot divide
\(\mathsf{one}\). For every \(d,a,b\), including \(d=O\),
\[
 \mathsf{div}\ d\ a\to\mathsf{div}\ d\ (a+b)\to\mathsf{div}\ d\ b.
\]
For local coefficients \(rd=a\), \(sd=a+b\) and \(O<d\), the
constructed coefficient of \(b\) is \(s\mathbin{\dot{-}}r\).
This is a proof of a \(\mathsf{Prop}\) existential, not a function
extracting a witness from an arbitrary existential proof into
\(\mathsf{Set}\).
\end{hblemma}

\begin{hbtheorem}{thm:p12-har-prime-divisor}{proved=annex:p12.har.prime-divisor}{The source prime-divisor theorem, now with its proof}
The historical declaration has an actual derivation:
\[
 \mathsf{numbers\_gt1\_have\_primediv}:\Pi n:\mathsf{nat}.
             (\mathsf{one}<n)\to\mathsf{HPD}\ n.
\]
Its only classical step is \(\mathsf{tnd}\) on
\(\exists d.\mathsf{propdiv}\ d\ n\), derived from
\(\mathsf{DN}\); the remaining steps use course-of-values induction
and the proved transitivity of divisibility.
\end{hbtheorem}

\begin{hbtheorem}{thm:p12-har-prime-euclid}{proved=annex:p12.har.prime-euclid}{The factorial continuation and Euclid's unboundedness statement}
For every \(N:\mathsf{nat}\),
\[
 \exists p:\mathsf{nat}.\;\mathsf{prime}\ p\wedge N<p.
\]
The proof supplies the factorial bound and the theorem that every
\(\mathsf{one}\le m\le N\) divides \(\mathsf{fac}\ N\), then
applies the source prime-divisor theorem to \(\mathsf{fac}\ N+\mathsf{one}\).
It inherits that theorem's explicit \(\mathsf{DN}\) dependence.
No constructive prime-producing program follows by unlicensed
\(\mathsf{Prop}\)-to-\(\mathsf{Set}\) elimination. This is the first
formula of the source footnote. The footnote's informal variable-length
list display is a different representation; identifying it with a single
coded first-order arithmetic formula requires an explicit coding bridge,
not an assertion that first-order arithmetic cannot express finite lists
\cite[p.~1215, note~4; p.~1220]{BarendregtGeuvers2001}.
There is also a fully typed finite-family formulation within the same
regime. For any \(k:\mathsf{nat}\) and \(a:\mathsf{nat}\to\mathsf{nat}\),
\[
 \exists p.\;\mathsf{prime}\ p\wedge
           \Pi i:\mathsf{nat}.(i<k)\to\neg(p=a i).
\]
The proof bounds the first k values by the actually constructed data
fold \(B O\ a:=O\), \(B(S k)\ a:=B k\ a+a k\), then uses the
unboundedness result. It does not require the listed values to be prime,
so it entails the source's finite-prime-family avoidance assertion.
This dependent-function representation supplies no unproved
identification with an arithmetic sequence-coding library.
\end{hbtheorem}


\begin{hbtheorem}{thm:p12-har-development-objects}{proved=annex:p12.har.development-objects}{The two explicitly supplied source proof objects}
These are the two proof objects actually printed in the chapter, rather
than an assertion that it prints all the library derivations. Their
historical Set/Prop regime is the one of
\cref{def:p12-har-prime-regime}; the additional source-omitted library
proofs are constructed in \cref{thm:p12-har-prime-logical}.
The source \(\mathrm{non\_or}\) term constructs
\(\neg(a\lor b)\to(\neg a\land\neg b)\)
by pairing the functions \(q:a\mapsto p(\mathrm{or\_inl}\,a b q)\)
and \(q:b\mapsto p(\mathrm{or\_inr}\,a b q)\).
For the source existential presentation, its \(\mathrm{ex\_c}\)
term proves
\[
 (\exists x:A.Px)\ \longleftrightarrow\
       \neg\Pi x:A.\neg Px.
\]
The forward implication eliminates the existential into FF; the reverse
constructs double negation of that existential and applies the explicitly
declared DN axiom. Its dependence on DN remains in the proof term.
\end{hbtheorem}

\begin{hbdefinition}{def:p12-har-autarkic}{The source interfaces and their actual reconstructions}
The chapter distinguishes numerical computation, predicate decision, and
symbolic computation. For square root it proposes a represented function
\(F\) with a reusable proof obligation
\(\Pi n.(Fn)^2\le n<(Fn+1)^2\). The concrete computation
\(F\overline{45}\twoheadrightarrow\overline6\), together with
reflexivity and conversion, proves its numeral equation only after that
function's representation is declared; it does not by itself prove the
universal square-root specification.
The actual fixed-target program, universal bound and literal 45 reduction
are supplied in \cref{def:p12-har-autarkic-square-root,thm:p12-har-autarkic-square-root},
under the explicitly historical nat:Set regime.

For primality, \(\mathrm{KPrime}\) represents a characteristic
function with the two obligations
\(\Pi n.\mathrm{Prime}(n)\leftrightarrow
\mathrm{KPrime}(n)=\overline1\) and
\(\Pi n.\mathrm{KPrime}(n)=\overline0\lor
\mathrm{KPrime}(n)=\overline1\). A successful computation is usable
through that checked equivalence. The source notes the expense of its
generic primitive-recursion translation and unary numerals, and separately
mentions the Pocklington/Fermat-based primality work, including the stated
number \(1223334444555554444333221\). Those are attributed examples,
not a large-number certificate supplied by this chapter or newly checked
here. The direct bounded-search reconstruction in
\cref{def:p12-har-kprime-program,thm:p12-har-kprime-correct,thm:p12-har-kprime-61}
supplies the historical prime predicate's full positive/negative/range
specification, 0/1 boundary cases and the printed 61 example. It is not
the generic compiler or Pocklington/Fermat algorithm.

For symbolic normalization, source \(G:\mathrm{term}(\mathrm{Nat})
\to\mathrm{term}(\mathrm{Nat})\) acts on represented syntax and
source \(E\) interprets it. The chapter's desired quoted example is
\[
 G\bigl(\text{`}(x+1)(x+1)\text{'}\bigr)
       \twoheadrightarrow_\iota
       \text{`}x^2+2x+1\text{'}.
\]
For free variables, an explicit valuation \(\nu\) belongs in the
interpreter: \(E_\nu:\mathrm{term}(\mathrm{Nat})\to\mathrm{Nat}\).
The real reusable obligation is
\(\Pi\nu.\Pi t.E_\nu(Gt)=E_\nu(t)\). Instantiating it
at the quoted input and using the two evaluation computations yields
the polynomial identity. The chapter expressly omits its free-variable
bookkeeping; an adequate formal encoding must supply it rather than
identify syntax with values. This source \(E\) is an object-language
interpreter, distinct from the CC proof's metalevel raw interpretation.
The actual typed five-constructor interpreter and finite/open valuation
construction are supplied in
\cref{def:p12-har-autarkic-syntax,def:p12-har-autarkic-interpreter,thm:p12-har-autarkic-interpreter,thm:p12-har-autarkic-substitution}.
The actual sparse polynomial program, normal-form/retraction results and
universal preservation are supplied in
\cref{def:p12-har-poly-normalizer,def:p12-har-poly-quotation,thm:p12-har-poly-normalization,thm:p12-har-poly-preservation,thm:p12-har-poly-retraction}.
Its right-associated x-squared+(2x+1) output and explicit association
bridge to the alternate source parsing are given in
\cref{thm:p12-har-poly-example}; literal AST identity between those
parses is not asserted. Arbitrary fixed carriers require the stated
commutative-semiring laws, not just typed operations.
The same two-level interface supports encoded differentiation,
simplification and continuity arguments, with separate soundness proofs
\cite[pp.~1220--1223]{BarendregtGeuvers2001}.
The following treatments supply the particular F/KPrime/G/interpreter
constructions and, separately, the internal formula compiler and universal
correctness proof of
\cref{thm:p12-har-prcompiler-ok,thm:p12-har-prcompiler-predicate}.
The latter distinguishes actual internal formula computation from external
primitive-recursive prescription templates and from an automatic inverse
on arbitrary propositions. Pocklington/Fermat/large certificates,
differentiation and continuity extensions remain separate. Likewise the proved finite-family avoidance fold in
\cref{thm:p12-har-prime-euclid} is not an arithmetic sequence-coding
bridge. None of these obligations is discharged by the beta-only
normalization theorem or by the historical prime-divisor proof.
\end{hbdefinition}

\subsection{Actual autarkic programs and their reusable proofs}

\begin{hbdefinition}{def:p12-har-autarkic-regime}{The source computation mechanism and this concrete instantiation}
Section~4.3 leaves the program representing square root unspecified,
while requiring its universal adequacy proof; it also leaves the
arithmetic-syntax interpreter and its free-variable bookkeeping to be
constructed. Here those programs are built using the chapter's primitive
recursion mechanism in the explicitly historical arithmetic regime of
\cref{def:p12-har-prime-regime,def:p12-har-prime-arithmetic}.
Thus \(\mathsf{nat}:\mathsf{Set}\), its constructor computation,
generated induction into \(\mathsf{Prop}\), propositional equality,
and the already proved order/semiring library are the assumptions.
The available classical DN axiom is not used in the following programs
or proofs. No proof is eliminated to select data.

The source's abstract \(\mathsf{nat}:\mathsf{Type}\) recursor
schemas describe the same primitive-recursion mechanism at that different
declared sort; they are not being silently substituted for this Set-based
library. Nor are these programs terms of the unchanged beta-only HOL
calculus. Each recursion below has a fixed admitted target and explicitly
typed branches. General universe/inductive extension consistency remains
separate. These are concrete source-based reconstructions, not programs
printed in the chapter or claims of Coq execution
\cite[pp.~1220--1223]{BarendregtGeuvers2001}.
\end{hbdefinition}

\begin{hbdefinition}{def:p12-har-autarkic-eqb}{The actual equality-code and data selector}
Use O and one as the false and true \emph{data codes}. Put
\[
\begin{aligned}
 z&:=\mathsf{nat\_rec}\ \mathsf{one}
       (\lambda j:\mathsf{nat}.\lambda r:\mathsf{nat}.O),\\
 b&:=\lambda n:\mathsf{nat}.\lambda q:\mathsf{nat}\to\mathsf{nat}.
       \lambda m:\mathsf{nat}.
          \mathsf{nat\_rec}\ O
            (\lambda j:\mathsf{nat}.\lambda r:\mathsf{nat}.q j)\ m,\\
 \mathsf{eqb}&:=\mathsf{nat\_rec}\ z\ b,\\
 \mathsf{choose}&:=\lambda c:\mathsf{nat}.\lambda u:\mathsf{nat}.
  \lambda v:\mathsf{nat}.
    \mathsf{nat\_rec}\ v
       (\lambda j:\mathsf{nat}.\lambda r:\mathsf{nat}.u)\ c.
\end{aligned}
\]
The outer eqb recursion has target \(\mathsf{nat}\to\mathsf{nat}\);
the two inner recursions and choose have target nat. Consequently
\(\mathsf{eqb}:\mathsf{nat}\to\mathsf{nat}\to\mathsf{nat}\) and
\(\mathsf{choose}:\mathsf{nat}\to\mathsf{nat}\to\mathsf{nat}\to\mathsf{nat}\).
Their actual computation equations are
\[
\begin{aligned}
 \mathsf{eqb}\ O\ O&\longrightarrow^*\mathsf{one},&
 \mathsf{eqb}\ O\ (S m)&\longrightarrow^*O,\\
 \mathsf{eqb}\ (S n)\ O&\longrightarrow^*O,&
 \mathsf{eqb}\ (S n)\ (S m)&\longrightarrow^*\mathsf{eqb}\ n\ m,\\
 \mathsf{choose}\ O\ u\ v&\longrightarrow^*v,&
 \mathsf{choose}\ (S c)\ u\ v&\longrightarrow^*u.
\end{aligned}
\]
The arrows expand these finite definitions and perform beta/iota
contractions; the equations are not postulated boolean axioms.
\end{hbdefinition}

\begin{hblemma}{thm:p12-har-autarkic-eqb}{proved=annex:p12.har.autarkic-eqb}{The complete equality-code specification}
For all \(n,m:\mathsf{nat}\),
\[
\begin{aligned}
 &\mathsf{eqb}\ n\ m=O\ \vee\ \mathsf{eqb}\ n\ m=\mathsf{one},\\
 &\mathsf{eqb}\ n\ m=\mathsf{one}\ \longleftrightarrow\ n=m,\\
 &\mathsf{eqb}\ n\ m=O\ \longleftrightarrow\ \neg(n=m).
\end{aligned}
\]
This is a constructed decision procedure, not classical excluded middle
used to define a Set-valued program.
\end{hblemma}

\begin{hblemma}{thm:p12-har-autarkic-squares}{proved=annex:p12.har.autarkic-squares}{Square growth and the unique integer square-root bounds}
Write \(Q(r):=\mathsf{times}\ r\ r\). Then
\[
 Q(S r)=Q(r)+S(r+r),\qquad Q(r)<Q(S r).
\]
Squaring preserves weak and strict order on nat. If two values r,s
satisfy \(Q(r)\le n<Q(S r)\) and \(Q(s)\le n<Q(S s)\),
then \(r=s\). These assertions follow from the proved arithmetic
library; no square-root existence assumption is used.
\end{hblemma}

\begin{hbdefinition}{def:p12-har-autarkic-square-root}{The omitted square-root program F, explicitly constructed}
The actual fixed-target recursion is
\[
 F:=\mathsf{nat\_rec}\ O\
  (\lambda n:\mathsf{nat}.\lambda r:\mathsf{nat}.
     \mathsf{choose}\ (\mathsf{eqb}\ (S n)\ (Q(S r)))\ (S r)\ r).
\]
Its step has type \(\mathsf{nat}\to\mathsf{nat}\to\mathsf{nat}\),
so \(F:\mathsf{nat}\to\mathsf{nat}\). Operationally it increments
the current root exactly when the next input reaches the next square.
All branching inputs are computed natural-number data.
For a meta-level natural number k, retain the source's unary numeral
representation, written here
\(\ulcorner k\urcorner_{\mathsf{nat}}:=S^k O\);
this is not a quotation of arithmetic syntax into the term datatype below.
The nat and T subscripts on quotation brackets are editorial sort tags
exposing the two different representation types.
\end{hbdefinition}

\begin{hbtheorem}{thm:p12-har-autarkic-square-root}{proved=annex:p12.har.autarkic-square-root}{The universal F adequacy proof and the actual 45 reduction}
The constructed program satisfies the source's precise proof obligation
\[
 \Pi n:\mathsf{nat}.\quad Q(F n)\le n<Q(S(F n)).
\]
Since \(S(F n)=F n+\mathsf{one}\), this is
\((F n)^2\le n<(F n+1)^2\). Its values are the unique natural-number
integer square roots specified by these bounds. In particular the
explicit definitions admit the finite computation
\[
 F\ulcorner45\urcorner_{\mathsf{nat}}
     \twoheadrightarrow_{\beta\iota\delta}
          \ulcorner6\urcorner_{\mathsf{nat}}.
\]
The annex supplies the invariant's branches and the finite computation
trace. No executed kernel replay is asserted by this mathematical
reduction derivation \cite[p.~1221]{BarendregtGeuvers2001}.
\end{hbtheorem}

\begin{hbdefinition}{def:p12-har-autarkic-syntax}{The actual arithmetic syntax and its five generated branches}
The source \(\mathrm{term}(\mathrm{Nat})\) contains terms over
\(\langle\mathrm{Nat};+,\times,0,1\rangle\); it is not a datatype
of arbitrary typed lambda terms. Make its omitted variable representation
explicit with the positive data declaration \(T:\mathsf{Set}\) and
constructors
\[
\begin{aligned}
 Z,U&:T,&\mathsf{Var}&:\mathsf{nat}\to T,\\
 \mathsf{Add},\mathsf{Mul}&:T\to T\to T.
\end{aligned}
\]
The numeral argument of Var is an index for a free variable name, not its
value. This is the historical Set-based instance of the source's stated
inductive syntax mechanism. For each fixed admitted target C, the generated
fold has seeds \(z,u:C\), variable branch \(v:\mathsf{nat}\to C\),
and two branches \(a,m:T\to T\to C\to C\to C\).
The parenthesized five-argument fold notation means five curried
applications, not a new tuple datatype. Its constructor computations are
\[
\begin{aligned}
 \mathsf{fold}_T^C(z,u,v,a,m)Z&\longrightarrow z,\\
 \mathsf{fold}_T^C(z,u,v,a,m)U&\longrightarrow u,\\
 \mathsf{fold}_T^C(z,u,v,a,m)(\mathsf{Var}\ i)&\longrightarrow v i,\\
 \mathsf{fold}_T^C(z,u,v,a,m)(\mathsf{Add}\ s\ t)&\longrightarrow
    a s t(\mathsf{fold}_T^C(z,u,v,a,m)s)(\mathsf{fold}_T^C(z,u,v,a,m)t),\\
 \mathsf{fold}_T^C(z,u,v,a,m)(\mathsf{Mul}\ s\ t)&\longrightarrow
    m s t(\mathsf{fold}_T^C(z,u,v,a,m)s)(\mathsf{fold}_T^C(z,u,v,a,m)t).
\end{aligned}
\]
Only direct recursive fields occur, so the source strictly positive
branch construction gives precisely the two recursive values/proofs.
For a fixed \(P:T\to\mathsf{Prop}\), the generated induction has
branches \(PZ\), \(PU\), \(\Pi i.P(\mathsf{Var}\ i)\), and
\(\Pi s,t.Ps\to Pt\to P(\mathsf{Add}\ s\ t)\), respectively
\(\Pi s,t.Ps\to Pt\to P(\mathsf{Mul}\ s\ t)\).
These are the particular source-licensed recursion/induction schemas,
not a parametricity assumption or a new top-sort-polymorphic constant
\cite[pp.~1206--1209,1221--1222]{BarendregtGeuvers2001}.
\end{hbdefinition}

\begin{hbdefinition}{def:p12-har-autarkic-interpreter}{The constructed valuation-sensitive interpreter and substitution}
For each fixed \(A:\mathsf{Set}\), take specified values
\(0_A,1_A:A\), operations \(+_A,\times_A:A\to A\to A\),
and a valuation \(\nu:\mathsf{nat}\to A\). Define the actual fold
\[
 E_\nu^A:=\mathsf{fold}_T^A
   (0_A,1_A,\nu,
    \lambda s,t:T.\lambda a,b:A.a+_A b,
    \lambda s,t:T.\lambda a,b:A.a\times_A b).
\]
Every branch has the type just declared; no algebraic laws are needed
to form this interpreter. Its original natural-number instance uses
\(A=\mathsf{nat}\), \(0_A=O\), \(1_A=\mathsf{one}\), and the
already constructed plus/times. Abbreviate that instance by \(E_\nu\).
The arbitrary fixed A version is an explicit interpretation schema,
not a claim that the chapter prints a polymorphic model interpreter.

Quotation of a finite arithmetic expression is defined by its five
clauses: 0 maps to Z, 1 to U, and a name \(x_i\) with meta-level
natural-number index i maps to
\(\mathsf{Var}\ulcorner i\urcorner_{\mathsf{nat}}\);
a sum maps to Add of the two quotations, and a product to Mul of them. Write it
as \(\ulcorner p\urcorner_T\). Under a valuation, \(p[\nu]\)
replaces every \(x_i\) by
\(\nu\ulcorner i\urcorner_{\mathsf{nat}}\) and uses the declared values
and operations. There are no binders in this arithmetic grammar.
Finite open contexts have an actual valuation extension as well. For a
fixed A, define the data selector
\[
 \mathsf{sel}_A:=\lambda c:\mathsf{nat}.\lambda u,v:A.
      \mathsf{nat\_rec}\ v
        (\lambda j:\mathsf{nat}.\lambda r:A.u)\ c.
\]
For distinct meta-level indices \(i_1,\ldots,i_k\) with given
well-typed values \(a_1,\ldots,a_k:A\) in a valid context, take
\[
 \nu_\Gamma(m):=
  \mathsf{sel}_A(\mathsf{eqb}\ m\ulcorner i_1\urcorner_{\mathsf{nat}})a_1
   \bigl(\cdots\mathsf{sel}_A
    (\mathsf{eqb}\ m\ulcorner i_k\urcorner_{\mathsf{nat}})a_k\,0_A\cdots\bigr).
\]
This finite nesting is a fully specified term (at meta-level k=0 it
is the constant \(0_A\) function); the annex proves each assigned-index
computation. It uses neither choice nor a proof-to-data map.
Also define actual syntactic substitution, for \(\sigma:\mathsf{nat}\to T\),
by
\[
 \mathsf{sub}_\sigma:=\mathsf{fold}_T^T
 (Z,U,\sigma,
    \lambda s,t:T.\lambda a,b:T.\mathsf{Add}\ a\ b,
    \lambda s,t:T.\lambda a,b:T.\mathsf{Mul}\ a\ b).
\]
Neither program consults a logical proof to return data.
\end{hbdefinition}

\begin{hbtheorem}{thm:p12-har-autarkic-interpreter}{proved=annex:p12.har.autarkic-interpreter}{Actual interpretation, quotation and open-valuations agreement}
For every finite expression in the declared arithmetic grammar and every
well-typed valuation in a valid context, expansion and constructor reduction
give
\[
 E_\nu^A\ulcorner p\urcorner_T
            \twoheadrightarrow_{\beta\iota\delta}p[\nu].
\]
Here \(p[\nu]\) is the A-valued interpreted expression, not a
purported raw Nat expression of type A. For the Nat instance this is
exactly the source arithmetic quotation/evaluation relationship, with
the missing free-variable valuation supplied. Moreover, if
\(\Pi i:\mathsf{nat}.\nu i=\mu i\) in primitive equality on A,
then \(\Pi t:T.E_\nu^A t=E_\mu^A t\). Finite-context agreement on
the indices actually occurring in a given quotation is sufficient too,
as the annex's constructor-by-constructor proof recipe demonstrates.
No function extensionality or parametricity theorem is assumed.
\end{hbtheorem}

\begin{hbtheorem}{thm:p12-har-autarkic-substitution}{proved=annex:p12.har.autarkic-substitution}{The full open-term substitution equation}
For every fixed A with the declared operations, every
\(\nu:\mathsf{nat}\to A\), \(\sigma:\mathsf{nat}\to T\), and
\(t:T\), put \(\nu_\sigma:=\lambda i:\mathsf{nat}.E_\nu^A(\sigma i)\).
Then
\[
 E_\nu^A(\mathsf{sub}_\sigma t)=E_{\nu_\sigma}^A t.
\]
The equality is proved by the actual five generated induction branches,
not by supposing interpretation automatically commutes with substitution.
\end{hbtheorem}

\begin{hblemma}{thm:p12-har-autarkic-prime-bound}{proved=annex:p12.har.autarkic-prime-bound}{The actual finite coefficient bound needed by primality search}
For \(\mathsf{one}\le d\), every coefficient satisfies
\(x\le xd\). For such a d, consequently,
\[
 \mathsf{div}\ d\ n\to
      \exists x:\mathsf{nat}.(x\le n)\wedge(xd=n).
\]
In particular the source's proper-divisor proposition has the finite
search characterization
\[
\begin{aligned}
 &\exists d:\mathsf{nat}.\mathsf{propdiv}\ d\ n\\
 &\quad\longleftrightarrow\quad
 \exists d,x:\mathsf{nat}.
  (d\le n)\wedge(x\le n)\wedge(\mathsf{one}<d)\wedge(d<n)\wedge(xd=n).
\end{aligned}
\]
Both existentials are in Prop. This bounds the later data-search loops
used by the actual KPrime program in \cref{def:p12-har-kprime-program,thm:p12-har-kprime-correct}.
That program and its characteristic proof are supplied separately;
the loops are not obtained by extracting these Prop witnesses into data.
\end{hblemma}

\begin{proofboundary}{def:p12-har-autarkic-residual}{What these constructions do not supply}
The source square-root program and its full reusable proof are now supplied,
as are the actual arithmetic-syntax interpreter and its open-term laws.
The actual KPrime term and all order/boolean/bounded-fold/characteristic
proofs are supplied in \cref{def:p12-har-kprime-program,thm:p12-har-kprime-correct}.
The actual polynomial program and universal interpretation proof are
supplied in \cref{def:p12-har-poly-normalizer,def:p12-har-poly-quotation,thm:p12-har-poly-preservation}.
Its nontrivial source-pattern output is given with the explicit
association convention in \cref{thm:p12-har-poly-example}; it is not
identity G or a claim of syntax equality across different parses.
The external primitive-recursive representation templates and the distinct
internal formula compiler are supplied in
\cref{thm:p12-har-prcompiler-representation,thm:p12-har-prcompiler-ok,thm:p12-har-prcompiler-predicate}.
They are not consequences of these particular programs or an automatic
inverse from arbitrary propositions to formula data. The Pocklington/Fermat
variant and large-number certificate remain distinct additional constructions.
The global
extension metatheory and arithmetic sequence-coding bridge remain separate.
\end{proofboundary}

\subsection{The actual autarkic primality program}

\begin{hbdefinition}{def:p12-har-kprime-regime}{The exact predicate and computation regime}
The chapter requires a represented KPrime with
\(\mathrm{Prime}(n)\leftrightarrow\mathrm{KPrime}(n)=1\) and a
zero/one range proof. Instantiate that obligation at its explicitly
defined historical predicate \(\mathsf{prime}\) from
\cref{def:p12-har-prime-arithmetic}. The divisor characterization below
justifies the connection with ordinary natural-number primality, rather
than identifying predicates by capitalization alone.

The regime is still \(\mathsf{nat}:\mathsf{Set}\), with the source
primitive data recursion, induction into Prop and propositional equality
of \cref{def:p12-har-prime-regime}. Reuse the actual eqb/choose programs
and their specifications from \cref{def:p12-har-autarkic-eqb,thm:p12-har-autarkic-eqb},
and the proved bounded proper-divisor characterization of
\cref{thm:p12-har-autarkic-prime-bound}. All recursion targets below
are nat or nat-to-nat. DN is not used. This direct bounded-search
reconstruction is neither the cited generic primitive-recursive compiler
nor the separate Pocklington/Fermat algorithm; no modern kernel or
beta-only HOL implementation is substituted
\cite[pp.~1221--1223]{BarendregtGeuvers2001}.
\end{hbdefinition}

\begin{hbdefinition}{def:p12-har-kprime-bits}{The complete data-level boolean operations}
For natural-number data define
\[
\begin{aligned}
 \mathsf{Bit}(a)&:=(a=O)\vee(a=\mathsf{one}),\\
 \mathsf{andb}&:=\lambda a,b:\mathsf{nat}.
                      \mathsf{choose}\ a\ b\ O,\\
 \mathsf{orb}&:=\lambda a,b:\mathsf{nat}.
                      \mathsf{choose}\ a\ \mathsf{one}\ b,\\
 \mathsf{notb}&:=\lambda a:\mathsf{nat}.
                      \mathsf{choose}\ a\ O\ \mathsf{one}.
\end{aligned}
\]
Bit is a proposition about a previously computed datum; it is not the
datatype from which these functions extract their inputs. The operations
have respectively types nat-to-nat-to-nat, nat-to-nat-to-nat and
nat-to-nat, whether or not their inputs satisfy Bit.
\end{hbdefinition}

\begin{hblemma}{thm:p12-har-kprime-bits}{proved=annex:p12.har.kprime-bits}{Both boolean truth values, with the exact range premises}
For Bit a, notb a satisfies Bit. For Bit a and Bit b, andb/orb
also satisfy Bit. At these exact respective premises,
\[
\begin{aligned}
 \mathsf{andb}\ a\ b=\mathsf{one}&\leftrightarrow
          (a=\mathsf{one})\wedge(b=\mathsf{one}),\\
 \mathsf{andb}\ a\ b=O&\leftrightarrow
          (a=O)\vee(b=O),\\
 \mathsf{orb}\ a\ b=\mathsf{one}&\leftrightarrow
          (a=\mathsf{one})\vee(b=\mathsf{one}),\\
 \mathsf{orb}\ a\ b=O&\leftrightarrow(a=O)\wedge(b=O),\\
 \mathsf{notb}\ a=\mathsf{one}&\leftrightarrow a=O,\qquad
 \mathsf{notb}\ a=O\leftrightarrow a=\mathsf{one}.
\end{aligned}
\]
For a Bit datum, \(a=O\leftrightarrow\neg(a=\mathsf{one})\)
and \(a=\mathsf{one}\leftrightarrow\neg(a=O)\), constructively.
\end{hblemma}

\begin{hbdefinition}{def:p12-har-kprime-order}{The actual weak- and strict-order programs}
Put
\[
\begin{aligned}
 l_0&:=\lambda m:\mathsf{nat}.\mathsf{one},\\
 l_S&:=\lambda n:\mathsf{nat}.\lambda q:\mathsf{nat}\to\mathsf{nat}.
  \lambda m:\mathsf{nat}.
    \mathsf{nat\_rec}\ O
      (\lambda j:\mathsf{nat}.\lambda r:\mathsf{nat}.q j)\ m,\\
 \mathsf{leb}&:=\mathsf{nat\_rec}\ l_0\ l_S,\\
 \mathsf{ltb}&:=\lambda n,m:\mathsf{nat}.\mathsf{leb}\ (S n)\ m.
\end{aligned}
\]
The outer recursion target is nat-to-nat, the inner target nat.
Both resulting functions have type nat-to-nat-to-nat. Their computation
clauses are leb O m to one, leb(S n)O to O, and
leb(S n)(S m) to leb n m; ltb expands the source definition of strict order.
\end{hbdefinition}

\begin{hblemma}{thm:p12-har-kprime-order}{proved=annex:p12.har.kprime-order}{The complete weak/strict comparison specifications}
For all n,m:nat the two codes satisfy Bit, and
\[
\begin{aligned}
 \mathsf{leb}\ n\ m=\mathsf{one}&\leftrightarrow n\le m,&
 \mathsf{leb}\ n\ m=O&\leftrightarrow\neg(n\le m),\\
 \mathsf{ltb}\ n\ m=\mathsf{one}&\leftrightarrow n<m,&
 \mathsf{ltb}\ n\ m=O&\leftrightarrow\neg(n<m).
\end{aligned}
\]
These compare the actual source leseq/lthan relations, not an imported
order implementation or classical decision oracle.
\end{hblemma}

\begin{hbdefinition}{def:p12-har-kprime-any}{The inclusive bounded existential fold}
For \(p:\mathsf{nat}\to\mathsf{nat}\), define
\[
 \mathsf{Any}:=\lambda N:\mathsf{nat}.\lambda p:\mathsf{nat}\to\mathsf{nat}.
    \mathsf{nat\_rec}\ (p O)
       (\lambda k:\mathsf{nat}.\lambda r:\mathsf{nat}.
           \mathsf{orb}\ r\ (p(S k)))\ N.
\]
It has type nat-to-(nat-to-nat)-to-nat and actually scans every input
from O through N, \emph{including both endpoints}. Its seed and step
have target nat; it consumes the data function p, not a Prop existential.
\end{hbdefinition}

\begin{hbtheorem}{thm:p12-har-kprime-any}{proved=annex:p12.har.kprime-any}{The complete bounded-fold positive and negative proofs}
For each fixed p with \(\Pi i:\mathsf{nat}.\mathsf{Bit}(p i)\),
the result satisfies Bit, and for every N:nat,
\[
\begin{aligned}
 \mathsf{Any}\ N\ p=\mathsf{one}&\leftrightarrow
             \exists i:\mathsf{nat}.(i\le N)\wedge(p i=\mathsf{one}),\\
 \mathsf{Any}\ N\ p=O&\leftrightarrow
             \Pi i:\mathsf{nat}.(i\le N)\to(p i=O),\\
 \mathsf{Any}\ N\ p=O&\leftrightarrow
             \neg\exists i:\mathsf{nat}.(i\le N)\wedge(p i=\mathsf{one}).
\end{aligned}
\]
Only a code-range premise on the independently supplied function is
assumed. The fold's existential and universal semantic specifications
are proved, not included among its premises.
\end{hbtheorem}

\begin{hbdefinition}{def:p12-har-kprime-program}{The complete nested coefficient/divisor search and KPrime term}
The actual composed terms are
\[
\begin{aligned}
 \mathsf{Coeff}&:=\lambda d,n:\mathsf{nat}.
   \mathsf{Any}\ n\ (\lambda x:\mathsf{nat}.
                \mathsf{eqb}\ (\mathsf{times}\ x\ d)\ n),\\
 \mathsf{Proper}&:=\lambda d,n:\mathsf{nat}.
   \mathsf{andb}\ (\mathsf{ltb}\ \mathsf{one}\ d)
    (\mathsf{andb}\ (\mathsf{ltb}\ d\ n)(\mathsf{Coeff}\ d\ n)),\\
 \mathsf{HasProper}&:=\lambda n:\mathsf{nat}.
        \mathsf{Any}\ n\ (\lambda d:\mathsf{nat}.\mathsf{Proper}\ d\ n),\\
 \mathsf{KPrime}&:=\lambda n:\mathsf{nat}.
   \mathsf{andb}\ (\mathsf{ltb}\ \mathsf{one}\ n)
             (\mathsf{notb}(\mathsf{HasProper}\ n)).
\end{aligned}
\]
Coeff and Proper have type nat-to-nat-to-nat; HasProper and KPrime
have type nat-to-nat. The coefficient is on the \emph{left} of times,
matching the source div d n definition. The two inclusive scans use the
same bound n, justified by the proved bounded-divisor characterization;
strict properness is still computed, not assumed from the loop bounds.
No data witness is extracted from the characterization's Prop existentials.
\end{hbdefinition}

\begin{hblemma}{thm:p12-har-kprime-proper}{proved=annex:p12.har.kprime-proper}{Every nested search code and both proper-divisor directions}
All Coeff/Proper/HasProper outputs satisfy Bit. Coeff d n=one exactly
when \(\exists x.(x\le n)\wedge(xd=n)\). Proper d n=one exactly
when \(\mathsf{propdiv}\ d\ n\); its O output is equivalent to
the negation of that proposition. HasProper n=one exactly when
\(\exists d.\mathsf{propdiv}\ d\ n\); its O output is equivalent
to the negation of that existential. In the reverse Proper direction,
the coefficient bound is used only after the premise one<d has established
one<=d; no invalid cancellation/bound for a zero divisor is used.
\end{hblemma}

\begin{hbtheorem}{thm:p12-har-kprime-correct}{proved=annex:p12.har.kprime-correct}{The reusable source primality characteristic, range and negative decision proof}
The actual program proves all of
\[
\begin{aligned}
 &\Pi n:\mathsf{nat}.\mathsf{Bit}(\mathsf{KPrime}\ n),\\
 &\Pi n:\mathsf{nat}.\quad
       \mathsf{prime}\ n\leftrightarrow\mathsf{KPrime}\ n=\mathsf{one},\\
 &\Pi n:\mathsf{nat}.\quad
       \neg\mathsf{prime}\ n\leftrightarrow\mathsf{KPrime}\ n=O.
\end{aligned}
\]
In particular prime is decidable constructively. Actual constructor
computations give \(\mathsf{KPrime}\ O\twoheadrightarrow O\) and
\(\mathsf{KPrime}\ \mathsf{one}\twoheadrightarrow O\);
both zero and one are refuted as primes. These statements include the
source's two obligations and do not depend on DN
\cite[p.~1221]{BarendregtGeuvers2001}.
As a worked negative search, the actual coefficient/divisor hits at
\(x=d=\overline2,n=\overline4\) give
\(\mathsf{KPrime}\,\overline4\twoheadrightarrow O\) and
\(\neg\mathsf{prime}\,\overline4\), where the bars denote unary
constructor numerals, not primitive decimal constants.
\end{hbtheorem}

\begin{hblemma}{thm:p12-har-kprime-divisor-characterization}{proved=annex:p12.har.kprime-divisor-characterization}{The ordinary divisor characterization proved at the source definitions}
The source predicate prime n is equivalent to
\[
 (\mathsf{one}<n)\wedge
       \Pi d:\mathsf{nat}.\mathsf{div}\ d\ n\to
                       ((d=\mathsf{one})\vee(d=n)).
\]
The forward proof supplies the zero-divisor and positive-coefficient
cases rather than assuming every divisor is automatically bounded by n.
\end{hblemma}

\begin{hblemma}{thm:p12-har-kprime-small-factor}{proved=annex:p12.har.kprime-small-factor}{The actual small-factor and no-multiple certificate lemmas}
For all natural data, \(a\le b\) implies \(ac\le bc\).
If \(qd<n<(S q)d\), then \(\neg\mathsf{div}\ d\ n\).
If one<n and \(n<(S b)(S b)\), then
\[
 (\exists d.\mathsf{propdiv}\ d\ n)\to
        \exists e.(e\le b)\wedge\mathsf{propdiv}\ e\ n.
\]
The factor may be the original divisor or its coefficient; both branches
are constructed inside Prop, without assuming a factorization algorithm.
\end{hblemma}

\begin{hbtheorem}{thm:p12-har-kprime-61}{proved=annex:p12.har.kprime-61}{The printed Prime(61) example, with proof and literal computation}
Let \(\overline{k}:=S^k O\) denote the actual source unary numeral in this
local statement. Then \(\mathsf{prime}\,\overline{61}\) is proved,
and the explicit program admits a finite beta/iota/delta reduction
\[
 \mathsf{KPrime}\,\overline{61}\twoheadrightarrow\mathsf{one}.
\]
The annex supplies the small-factor certificate, all six literal product
intervals, and the inclusive-loop reduction argument. This is not an
executed Coq computation. The printed
1223334444555554444333221 example has a distinct Pocklington/Fermat
attribution; no value of KPrime at that input or certificate for it is
claimed here \cite[pp.~1220,1223]{BarendregtGeuvers2001}.
\end{hbtheorem}

\begin{proofboundary}{def:p12-har-kprime-residual}{Exact variant and remaining source obligations}
The direct bounded-search primality component now has its actual term,
both range/characteristic directions, negative decision and the printed 61
example. The distinct generic computation construction and its universal
correctness proof are supplied in
\cref{thm:p12-har-prcompiler-ok,thm:p12-har-prcompiler-predicate}; this
does not identify its algorithm with direct KPrime or supply arbitrary
proposition reification. The Pocklington/Fermat algorithm and large
certificates remain separate. G's actual polynomial transformer and
universal open-valuation proof are supplied in
\cref{def:p12-har-poly-normalizer,def:p12-har-poly-quotation,thm:p12-har-poly-preservation,thm:p12-har-poly-retraction}.
Differentiation/further simplification/continuity extensions, arithmetic coding
and global extension metatheory remain separate obligations.
\end{proofboundary}

\subsection{The constructed symbolic polynomial normalizer}

\begin{hbdefinition}{def:p12-har-poly-regime}{Exact grammar, algebraic hypotheses and normal-form convention}
Use exactly the five-constructor arithmetic syntax T and E from
\cref{def:p12-har-autarkic-syntax,def:p12-har-autarkic-interpreter}:
Z, U, Var i, Add s t and Mul s t. The following program implements
the omitted source G using positive finite data and their generated
primitive recursion in the explicitly historical Set-based regime.
It is not a beta-only HOL term, an imported polynomial library or the
generic compiler cited elsewhere by the chapter.

For interpretation fix A:Set, typed zero/unit/addition/multiplication,
and primitive book equality on A. Assume the actual commutative-semiring
laws: associative and commutative addition with zero identity;
associative and commutative multiplication with unit identity;
both distributive laws; and zero annihilation on both sides.
These are explicit propositional equation families. They are supplied
by the constructed arithmetic library when A=nat. An arbitrary carrier
with uninterpreted operations does not have these laws.
Every well-typed valuation nu:nat-to-A in a valid context is admitted.
Neither DN, choice, functional extensionality nor parametricity is used.

Normal data are sparse polynomials: a monomial is a nondecreasing finite
word of variable indices, repetitions recording powers; coefficients
are positive naturals; rows have distinct words in strictly descending
lexicographic order, with the empty word smaller than every nonempty
word. The quotation convention is a right-associated sum with no
trailing zero, natural literals made from U/Add, coefficient-one factors
omitted, and no trailing unit in nonempty monomial products. This
explicit convention may differ from implicit source parenthesization;
the example below gives both the actual output and its association bridge
\cite[pp.~1221--1222]{BarendregtGeuvers2001}.
\end{hbdefinition}

\begin{hbdefinition}{def:p12-har-poly-data}{The actual finite data types and recursion support}
Declare the strictly positive data
\[
\begin{aligned}
 W&:\mathsf{Set},& []&:W,&i::w&:W\quad(i:\mathsf{nat},w:W),\\
 R&:\mathsf{Set},&\mathsf{row}&:\mathsf{nat}\to W\to R,\\
 P&:\mathsf{Set},&\epsilon&:P,&r\mathbin{\#}p&:P\quad(r:R,p:P).
\end{aligned}
\]
For a fixed admitted target C, \(\mathrm{Rec}_W\) has seed z:C and branch
nat-to-W-to-C-to-C; \(\mathrm{Rec}_R\) has a single branch nat-to-W-to-C;
\(\mathrm{Rec}_P\) has seed z:C and branch R-to-P-to-C-to-C. Their computations
at constructors are respectively z/f i w(rec w), f c w, and
z/f r p(rec p). Generated Prop induction receives exactly the
recursive-field hypotheses; R has none. These specified instances,
not an assumed list/map library or a top-sort-polymorphic constant,
are all finite-list support used below. A case operation is the same
recursor with its recursive value ignored.

For each fixed C put
\(I_C:=\lambda b:\mathsf{nat}.\lambda u,v:C.
\mathsf{nat\_rec}\ v(\lambda j:\mathsf{nat}.\lambda r:C.u)b\).
Thus \(I_C\) O u v contracts to v and \(I_C\) one u v to u.
Use the already constructed eqb, leb, ltb, andb and orb data codes.
Write \(\mathsf{empty}_W\)/\(\mathsf{empty}_P\) for their constructor cases returning one at
the empty constructor and O at a cons. Row coefficient/word and
word head/tail projections are their constructor cases; head []=O,
tail []=[] are explicit defaults.

All displayed functional equations below are abbreviations for these
fixed-target constructor recursions; branch types and nested recursion
orders are checked in the annex. No well-founded search is silently
postulated from an ordering relation.
\end{hbdefinition}

\begin{hbdefinition}{def:p12-har-poly-word-code}{Computed word equality, order and product}
The actual equality and strict-lexicographic codes are
\[
\begin{aligned}
 \mathsf{weq}([],[])&=\mathsf{one},&
 \mathsf{weq}([],j::v)&=O,&\mathsf{weq}(i::u,[])&=O,\\
 \mathsf{weq}(i::u,j::v)&=
          \mathsf{andb}(\mathsf{eqb}\ i\ j)(\mathsf{weq}(u,v)),\\
 \mathsf{wlt}([],[])&=O,&
 \mathsf{wlt}([],j::v)&=\mathsf{one},&\mathsf{wlt}(i::u,[])&=O,\\
 \mathsf{wlt}(i::u,j::v)&=
   \mathsf{orb}(\mathsf{ltb}\ i\ j)
       (\mathsf{andb}(\mathsf{eqb}\ i\ j)(\mathsf{wlt}(u,v))).
\end{aligned}
\]
Each is an outer \(\mathrm{Rec}_W\) at target W-to-nat, with a case on the
second word in its seed/step and recursive value q:W-to-nat used
only at v. This is an actual nested primitive recursion.
Let u<lex v be the inductive Prop relation with constructors:
[]<lex j::v; i<j implies i::u<lex j::v; and i=j together with
u<lex v implies i::u<lex j::v. The conclusion depends only on
data endpoints, not on a proof argument.

Monomial multiplication merges words:
\[
 \mathsf{MW}([],v)=v,\qquad \mathsf{MW}(i::u,[])=i::u,
\]
and on two conses it computes
\[
 \mathsf{MW}(i::u,j::v)=
 I_W(\mathsf{leb}\ i\ j)
       (i::\mathsf{MW}(u,j::v))
       (j::\mathsf{MW}(i::u,v)).
\]
To make both recursive calls legitimate, its outer \(\mathrm{Rec}_W\) has target
W-to-W, seed lambda v.v, and cons branch
\[
 \lambda i:\mathsf{nat}.\lambda u:W.\lambda q:W\to W.
 \lambda v:W.\mathrm{Rec}_W(i::u)
  (\lambda j:\mathsf{nat}.\lambda v':W.\lambda z:W.
     I_W(\mathsf{leb}\ i\ j)(i::q(j::v'))(j::z))v.
\]
The inner recursive value z is the required same-first-word/shorter-
second-word call. Repeated variable indices are retained, not deduplicated.
\end{hbdefinition}

\begin{hblemma}{thm:p12-har-poly-word-order}{proved=annex:p12.har.poly-word-order}{All word-code cases and the strict total ordering}
Weq and wlt satisfy Bit; weq u v=one iff u=v in primitive equality
on W, and wlt u v=one iff u<lex v. Their O codes are equivalent
to the respective negations. Lex order is irreflexive, transitive,
asymmetric and trichotomous. No word-equality oracle is assumed.
\end{hblemma}

\begin{hbdefinition}{def:p12-har-poly-invariants}{The concrete sparse-data invariants}
Use the inductive element-occurrence relation on W (head and tail
constructors) and row-occurrence relation on P, again with head/tail
constructors. \(\mathrm{Sorted}_W\) has the empty and singleton cases, and
\(\mathrm{Sorted}_W\)(i::j::w) precisely when i<=j and \(\mathrm{Sorted}_W\)(j::w).
Good epsilon holds. Good(row c w \# p) precisely requires c not=O,
\(\mathrm{Sorted}_W\) w, Good p and, for every row c' v occurring in p, v<lex w.
These are Prop predicates with the displayed generated proof branches;
they are not proof fields consumed by the program.
\end{hbdefinition}

\begin{hbdefinition}{def:p12-har-poly-semantics}{Actual coefficient, word and polynomial interpretations}
In the fixed commutative semiring and valuation put
\[
\begin{aligned}
 \chi&:=\mathsf{nat\_rec}\ 0_A
                (\lambda k:\mathsf{nat}.\lambda a:A.a+_A1_A),\\
 V_\nu&:=\mathrm{Rec}_W1_A
       (\lambda i:\mathsf{nat}.\lambda w:W.\lambda a:A.\nu i\times_A a),\\
 B_\nu(\mathsf{row}\ c\ w)&:=\chi c\times_A V_\nu w,\\
 L_\nu&:=\mathrm{Rec}_P0_A
       (\lambda r:R.\lambda p:P.\lambda a:A.B_\nu r+_A a).
\end{aligned}
\]
The row interpretation is its \(\mathrm{Rec}_R\) branch; targets are the fixed A.
These are actual functions, not interpretations postulated to satisfy
the normalizer's desired conclusion.
\end{hbdefinition}

\begin{hblemma}{thm:p12-har-poly-coefficients}{proved=annex:p12.har.poly-coefficients}{Every coefficient homomorphism clause}
\(\chi(0)=0_A\), \(\chi(\mathsf{one})=1_A\),
\(\chi(c+d)=\chi c+_A\chi d\), and
\(\chi(cd)=\chi c\times_A\chi d\).
In the Nat instance \(\chi(n)=n\). The proofs use the declared semiring
laws and generated nat induction, not an assumed homomorphism.
\end{hblemma}

\begin{hblemma}{thm:p12-har-poly-word-merge}{proved=annex:p12.har.poly-word-merge}{The actual merge preserves sortedness and every variable occurrence}
MW preserves \(\mathrm{Sorted}_W\) on two sorted inputs. An index occurs in MW u v
iff it occurs in u or v. In any declared commutative-semiring interpretation,
\(V_\nu(\mathsf{MW}(u,v))=V_\nu(u)\times_A V_\nu(v)\),
where V is the word interpretation defined above. The semantic equation
holds even on unsorted words; commutativity is an explicit premise.
\end{hblemma}

\begin{hbdefinition}{def:p12-har-poly-insert}{Actual insertion, sum, scale and product}
Put mk(c,w,p):=\(I_P\)(eqb c O)p(row c w \# p). Define
\[
 \mathsf{Ins}(c,w,p):=
 I_P(\mathsf{eqb}\ c\ O)p
  \bigl(\mathrm{Rec}_P(\mathsf{mk}(c,w,\epsilon))f_{c,w}\ p\bigr),
\]
where the branch on head row d v, tail s and recursive value z:P is
\[
 f_{c,w}(\mathsf{row}\ d\ v,s,z):=
 I_P(\mathsf{weq}(w,v))\mathsf{mk}(c+d,w,s)
  \bigl(I_P(\mathsf{wlt}(v,w))
    (\mathsf{row}\ c\ w\mathbin{\#}(\mathsf{row}\ d\ v\mathbin{\#}s))
    (\mathsf{row}\ d\ v\mathbin{\#}z)\bigr).
\]
The branch's row pattern is the single \(\mathrm{Rec}_R\)/case branch, not a
destructive mutation. Zero input coefficients return p; an equal key
combines coefficients, a larger new key is prepended, and a smaller
one recurses through the tail. The mk zero test is present also at
the combination branch.
\[
\begin{aligned}
 \mathsf{Sum}(p,q)&:=\mathrm{Rec}_P q
  (\lambda r:R.\lambda s:P.\lambda z:P.
        \mathrm{case}_R(\lambda c:\mathsf{nat}.\lambda w:W.
                            \mathsf{Ins}(c,w,z))r)\ p,\\
 \mathsf{Scale}(c,w,p)&:=\mathrm{Rec}_P\epsilon
  (\lambda r:R.\lambda s:P.\lambda z:P.
   \mathrm{case}_R(\lambda d:\mathsf{nat}.\lambda v:W.
              \mathsf{Ins}(cd,\mathsf{MW}(w,v),z))r)\ p,\\
 \mathsf{Prod}(p,q)&:=\mathrm{Rec}_P\epsilon
  (\lambda r:R.\lambda s:P.\lambda z:P.
   \mathrm{case}_R(\lambda c:\mathsf{nat}.\lambda w:W.
              \mathsf{Sum}(\mathsf{Scale}(c,w,q),z))r)\ p.
\end{aligned}
\]
Every recursive target here is P; all source syntax is expanded into
the displayed ordinary data-recursion applications.
\end{hbdefinition}





\begin{hblemma}{thm:p12-har-poly-insertion}{proved=annex:p12.har.poly-insertion}{Every insertion branch, key order and value}
For Good p and \(\mathrm{Sorted}_W\) w, Ins(c,w,p) is Good, including c=O.
Every output key is either w or an input key. For arbitrary raw p,w,c,
\[
 L_\nu(\mathsf{Ins}(c,w,p))=
           (\chi c\times_A V_\nu w)+_A L_\nu p.
\]
Inserting a positive row whose key is larger than all input keys gives
that row followed by the unchanged input, in primitive book equality.
\end{hblemma}

\begin{hblemma}{thm:p12-har-poly-operations}{proved=annex:p12.har.poly-operations}{The complete computed sum/product and all zero cases}
Sum and Prod preserve Good on Good inputs; Scale(c,w,p) preserves
Good when w is sorted and p is Good. Their actual interpreted values are
\[
\begin{aligned}
 L_\nu(\mathsf{Sum}(p,q))&=L_\nu p+_A L_\nu q,\\
 L_\nu(\mathsf{Scale}(c,w,p))&=(\chi c\times_A V_\nu w)\times_A L_\nu p,\\
 L_\nu(\mathsf{Prod}(p,q))&=L_\nu p\times_A L_\nu q.
\end{aligned}
\]
The semantic equations hold on raw data. Sum(epsilon,q)=q;
Prod(epsilon,q) and Prod(p,epsilon) return epsilon; Scale(0,w,p)
returns epsilon. Zero coefficients are dropped rather than retained
as false nonzero rows.
\end{hblemma}

\begin{hbdefinition}{def:p12-har-poly-normalizer}{All five actual normalization branches}
Write single(c,w):=mk(c,w,epsilon). The actual data normalization is
\[
 N:=\mathsf{fold}_T^P\bigl(
 \epsilon,\mathsf{single}(\mathsf{one},[]),
 \lambda i:\mathsf{nat}.\mathsf{single}(\mathsf{one},i::[]),
 \lambda s,t:T.\lambda p,q:P.\mathsf{Sum}(p,q),
 \lambda s,t:T.\lambda p,q:P.\mathsf{Prod}(p,q)\bigr).
\]
Thus N:T-to-P is an actual fold, not the assertion that some normal
polynomial exists.
\end{hbdefinition}

\begin{hbtheorem}{thm:p12-har-poly-normalization}{proved=annex:p12.har.poly-normalization}{The complete five-branch normal-form and universal value proof}
For every t:T, Good(Nt), and for every valid open valuation in the
declared commutative semiring,
\(L_\nu(Nt)=E_\nu^A t\).
The program is valuation-independent; the preservation theorem is not.
\end{hbtheorem}



\begin{hbdefinition}{def:p12-har-poly-quotation}{Actual quotation-back and G}
The four concrete quotation recursions are
\[
\begin{aligned}
 \mathsf{lit}&:=\mathsf{nat\_rec}\ Z
    (\lambda k:\mathsf{nat}.\lambda t:T.
        I_T(\mathsf{eqb}\ k\ O)U(\mathsf{Add}\ t\ U)),\\
 Q_W&:=\mathrm{Rec}_W U
    (\lambda i:\mathsf{nat}.\lambda w:W.\lambda t:T.
       I_T(\mathsf{empty}_W w)(\mathsf{Var}\ i)
                    (\mathsf{Mul}(\mathsf{Var}\ i)t)),\\
 Q_R(\mathsf{row}\ c\ w)&:=
   I_T(\mathsf{empty}_W w)(\mathsf{lit}\ c)
      (I_T(\mathsf{eqb}\ c\ \mathsf{one})(Q_Ww)
                          (\mathsf{Mul}(\mathsf{lit}\ c)(Q_Ww))),\\
 Q_P&:=\mathrm{Rec}_P Z
   (\lambda r:R.\lambda p:P.\lambda t:T.
        I_T(\mathsf{empty}_P p)(Q_Rr)(\mathsf{Add}(Q_Rr)t)),\\
 G&:=\lambda t:T.Q_P(Nt).
\end{aligned}
\]
\(Q_R\) is the displayed single \(\mathrm{Rec}_R\) branch. All four targets are T.
The source's \(x^2\) and \(2x\) are printing sugar for Mul X X and
Mul(Add U U)X; no additional constructor is inserted into T.
\end{hbdefinition}

\begin{hblemma}{thm:p12-har-poly-quotation}{proved=annex:p12.har.poly-quotation}{All quotation branches and the actual value bridge}
For every raw datum,
\(E_\nu^A(\mathsf{lit}\ c)=\chi c\),
\(E_\nu^A(Q_W w)=V_\nu w\),
\(E_\nu^A(Q_R r)=B_\nu r\), and
\(E_\nu^A(Q_P p)=L_\nu p\).
Each zero, empty, singleton, coefficient-one and other-coefficient case
is constructed in the annex.
\end{hblemma}

\begin{hbtheorem}{thm:p12-har-poly-preservation}{proved=annex:p12.har.poly-preservation}{G's full universal open-valuation preservation}
For every fixed declared commutative semiring, every valid-context
\(\nu:\mathsf{nat}\to A\), and every t:T,
\[
 E_\nu^A(Gt)=E_\nu^A t.
\]
In particular this supplies the source Nat obligation at every open
natural-number valuation. No equality of valuation functions or closure
of the input expression is assumed.
\end{hbtheorem}

\begin{hbtheorem}{thm:p12-har-poly-retraction}{proved=annex:p12.har.poly-retraction}{The actual normal-form retraction and idempotence}
For Good p, \(N(Q_Pp)=p\) in primitive equality on P.
Consequently G(Gt)=Gt in primitive equality on T, and every output is
the declared quotation of Good sparse data. These are book equalities,
not a claim of definitional idempotence on opaque open indices or of
semantic completeness for arbitrary semiring models.
\end{hbtheorem}

\begin{hbtheorem}{thm:p12-har-poly-example}{proved=annex:p12.har.poly-example}{The source square-expansion example and its precise output convention}
Put X:=Var O. The actual computation is
\[
\begin{aligned}
 G(\mathsf{Mul}(\mathsf{Add}\ X\ U)(\mathsf{Add}\ X\ U))
 \twoheadrightarrow_{\beta\iota\delta}
 \mathsf{Add}(\mathsf{Mul}\ X\ X)
  (\mathsf{Add}(\mathsf{Mul}(\mathsf{Add}\ U\ U)X)U).
\end{aligned}
\]
This is \(x^2\)+(2x+1) under the declared right-associated quotation.
The left-associated source reading (\(x^2\)+2x)+1 has the same interpretation
by the explicit addition-associativity law. Literal AST identity with
that alternate parsing is not asserted. This is a nontrivial symbolic
output of the constructed G, not identity G or only a numerical example
\cite[pp.~1221--1222]{BarendregtGeuvers2001}.
\end{hbtheorem}

\begin{proofboundary}{def:p12-har-poly-residual}{What the actual normalizer does not identify}
The program is a sparse formal natural-coefficient polynomial normalizer
for the declared five-constructor grammar. Preservation for a non-Nat
interpretation requires its actual commutative-semiring laws; canonical
data do not imply semantic completeness for a model with extra equations.
Quotation's association convention is explicit rather than an assertion
of a parser-dependent source AST identity. The distinct generic computation
construction is given in
\cref{thm:p12-har-prcompiler-representation,thm:p12-har-prcompiler-ok,thm:p12-har-prcompiler-predicate};
polynomial normalization alone does not supply it or arbitrary proposition
reification. Pocklington/Fermat and large-number certificates, differentiation and
continuity extensions, arithmetic coding and global extension metatheory
remain distinct source obligations. No kernel replay is claimed here.
\end{proofboundary}

\subsection{The generic primitive-recursive computation construction}

\begin{hbdefinition}{def:p12-har-prcompiler-source}{The actual internal program and the external syntax boundary}
The chapter's generic computation promise cites Oostdijk--Geuvers.
The author-held preprint, dated 14 April 2000, gives an \emph{internal}
datatype \(\mathsf{form}:\mathsf{Set}\), an interpretation into
\(\mathsf{Prop}\), an interpreter into \(\mathsf{bool}:\mathsf{Set}\),
and one universal correctness proof. This is stronger than a metalevel
assertion that primitive-recursive functions have representing terms.
Conversely, its atomic arguments are already Coq nat terms: the paper
does not define an internal datatype of primitive-recursive function
prescriptions. Its discussion expressly leaves automatic inverse
translation from an arbitrary proposition to form outside the object
language, and distinguishes the older external Lego generator.
The dated preprint and the eventual 2002 journal article are distinct
manifestations; the chapter's provisional 2001 citation is not a claim
of byte or pagination identity
\cite{OostdijkGeuvers2002ProofByComputation}.
The chapter promise is located at pp.~1222--1223
\cite{BarendregtGeuvers2001}; the construction source is the author preprint
\cite[\S\S4--5, pp.~17--25]{OostdijkGeuvers2000ProofByComputationPreprint}.

Use the historical nat:Set, its data recursion and Prop induction,
and the equality/order library of
\cref{def:p12-har-prime-regime,thm:p12-har-prime-equality,thm:p12-har-prime-order}.
For a fixed admitted target C, write
\(\mathsf{nat\_rec}_C:C\to(\mathsf{nat}\to C\to C)\to\mathsf{nat}\to C\),
with zero and successor computation. The source extension also admits
bool and the strictly positive form declaration below, including its
higher-order recursive fields. These are the selected inductive
extensions of the paper's \(\lambda\mathsf{PRED}\omega\), not
nat:Type, the beta-only HOL calculus, or a theorem of consistency for
every inductive extension. No DN or Prop-to-Set witness extraction is used.
The paper's base sorts are
\(\{\mathsf{Prop},\mathsf{Set},\mathsf{Type},\mathsf{Type}^{s}\}\),
with axioms \(\mathsf{Prop}:\mathsf{Type}\) and
\(\mathsf{Set}:\mathsf{Type}^{s}\). Its product sort pairs are
\((\mathsf{Set},\mathsf{Set}),(\mathsf{Set},\mathsf{Type}),
(\mathsf{Type},\mathsf{Type}),(\mathsf{Prop},\mathsf{Prop}),
(\mathsf{Set},\mathsf{Prop}),(\mathsf{Type},\mathsf{Prop})\),
with result the second sort. Thus the proposition-code target Prop
is in Type, distinct from a proof target in Prop. Constructor computation
belongs to the separately selected inductive extension, not the base
beta conversion \cite[Definition~1, p.~6; \S3]{OostdijkGeuvers2000ProofByComputationPreprint}.
\end{hbdefinition}

\begin{hbdefinition}{def:p12-har-prcompiler-prescriptions}{Explicit arities and all external primitive-recursive templates}
To state function representation precisely, use the following
\emph{editorial metalevel} finite ranked syntax. It is not attributed as
an internal datatype printed in the primary paper. At arity k allow
\(Z_k\); allow \(\mathsf{Succ}\) at arity one; allow
\(P_i^k\) for \(k\ge1\) and \(1\le i\le k\). If f has arity m
and each \(g_i\) arity k, then \(\mathsf{Comp}_k(f;g_1,\ldots,g_m)\)
has arity k. If g has arity k and h arity \(k+2\), then
\(\mathsf{Rec}_k(g,h)\) has arity \(k+1\), with the recursion
argument \emph{last}. Both k and m may be zero. There is no projection
of arity zero. Put \(T_0:=\mathsf{nat}\) and
\(T_{k+1}:=\mathsf{nat}\to T_k\); an empty argument list means
no application, not an assumed unit datatype.

The compiler \(\mathcal C\) emits these actual annotated templates:
\[
\begin{aligned}
 \mathcal C Z_k&:=\lambda x_1:\mathsf{nat}\ldots\lambda x_k:\mathsf{nat}.O,\\
 \mathcal C\mathsf{Succ}&:=\lambda x:\mathsf{nat}.Sx,\\
 \mathcal C P_i^k&:=\lambda x_1:\mathsf{nat}\ldots\lambda x_k:\mathsf{nat}.x_i,\\
 \mathcal C\mathsf{Comp}_k(f;\vec g)&:=
   \lambda\vec x:\mathsf{nat}.(\mathcal C f)
           ((\mathcal C g_1)\vec x)\cdots((\mathcal C g_m)\vec x),\\
 \mathcal C\mathsf{Rec}_k(g,h)&:=\lambda\vec x:\mathsf{nat}.
   \mathsf{nat\_rec}_{\mathsf{nat}}
      ((\mathcal C g)\vec x)
      (\lambda n:\mathsf{nat}.\lambda r:\mathsf{nat}.
                         (\mathcal C h)\vec x\ n\ r).
\end{aligned}
\]
Here \(\lambda\vec x\) abbreviates exactly k annotated abstractions.
At \(k=0\), \(\mathcal C Z_0=O\); at \(m=0\), the composition
body is just \(\mathcal C f:T_0\). The last template itself returns a
function of the recursion argument, so no argument is omitted.
\end{hbdefinition}

\begin{hblemma}{thm:p12-har-prcompiler-typing}{proved=annex:p12.har.prcompiler-typing}{Typing every ranked prescription}
For every finite prescription e of arity k,
\(\mathcal C e:T_k\) in the historical data regime. The template
construction is effective structural recursion on the finite prescription;
it requires no oracle for whether an arbitrary term is primitive recursive.
Its only recursion target is nat, with the indicated parameter context.
\end{hblemma}

\begin{hbdefinition}{def:p12-har-prcompiler-denotation}{The independently specified numerical meaning}
On ordinary metalevel natural numbers define \(\langle e\rangle\)
by zero, successor and projection, by ordinary function composition, and
by the equations
\[
 \langle\mathsf{Rec}_k(g,h)\rangle(\vec a,0)=\langle g\rangle(\vec a),
 \qquad
 \langle\mathsf{Rec}_k(g,h)\rangle(\vec a,n+1)
 =\langle h\rangle(\vec a,n,
          \langle\mathsf{Rec}_k(g,h)\rangle(\vec a,n)).
\]
This numerical recursion specifies a total function by ordinary induction
on n, independently of the emitted syntax. Write \(\overline a:=S^aO\).
\end{hbdefinition}

\begin{hbtheorem}{thm:p12-har-prcompiler-representation}{proved=annex:p12.har.prcompiler-representation}{Uniform numerical representation, including zero arity}
For every prescription e of arity k and every numerical tuple \(\vec a\),
\[
 (\mathcal C e)\overline{a_1}\cdots\overline{a_k}
       \longrightarrow_{\beta\delta\iota}^{*}
                         \overline{\langle e\rangle(\vec a)}.
\]
Consequently every standard-natural-number valuation of its free
parameters gives the independently specified numerical result. This is
a universal representation theorem, not a claim that an arbitrary opaque
open variable contracts to a numeral, or that every Coq nat term belongs
to the finite prescription language.
\end{hbtheorem}

\begin{hbdefinition}{def:p12-har-prcompiler-booleans}{The primary paper's actual boolean datatype and programs}
Retain \(\mathsf{bool}:\mathsf{Set}\) with constructors true and false,
not the earlier natural-number codes O and one. Its generated data
selector is \(\mathsf{if}_C:\mathsf{bool}\to C\to C\to C\),
namely \(\lambda b:\mathsf{bool}.\lambda u:C.\lambda v:C.
\mathsf{Rec}_{\mathsf{bool}}^C\ u\ v\ b\),
with the two constructor computations, and its Prop induction has both
literal cases. Define \(\mathsf{istrue}\ b\) by case analysis into
Prop, with true mapped to TT and false to FF; these are the historical
aliases for the paper's True and False. Explicitly,
\(\mathsf{istrue}:=\lambda b:\mathsf{bool}.
\mathsf{Rec}_{\mathsf{bool}}^{\mathsf{Prop}}\ \mathsf{TT}\ \mathsf{FF}\ b\).
Here Prop:Type is an admitted fixed target of this data recursor
\cite[Example~5, p.~14]{OostdijkGeuvers2000ProofByComputationPreprint}.
Define
\[
\begin{aligned}
 \mathsf{b\_not}\ a&:=\mathsf{if}_{\mathsf{bool}}\ a\ \mathsf{false}\ \mathsf{true},&
 \mathsf{b\_and}\ a\ b&:=\mathsf{if}_{\mathsf{bool}}\ a\ b\ \mathsf{false},\\
 \mathsf{b\_or}\ a\ b&:=\mathsf{if}_{\mathsf{bool}}\ a\ \mathsf{true}\ b,&
 \mathsf{b\_imp}\ a\ b&:=\mathsf{if}_{\mathsf{bool}}\ a\ b\ \mathsf{true},\\
 \mathsf{code}\ b&:=\mathsf{if}_{\mathsf{nat}}\ b\ \mathsf{one}\ O.
\end{aligned}
\]
For the paper's atomic comparisons, the following are complete fixed-target
terms, not only desired equations:
\[
\begin{aligned}
 q_0&:=\lambda m:\mathsf{nat}.\mathsf{true},\\
 e_0&:=\mathsf{nat\_rec}_{\mathsf{bool}}\ \mathsf{true}
             (\lambda j:\mathsf{nat}.\lambda r:\mathsf{bool}.\mathsf{false}),\\
 d&:=\lambda n:\mathsf{nat}.\lambda q:\mathsf{nat}\to\mathsf{bool}.
       \lambda m:\mathsf{nat}.
        \mathsf{nat\_rec}_{\mathsf{bool}}\ \mathsf{false}
           (\lambda j:\mathsf{nat}.\lambda r:\mathsf{bool}.q j)\ m,\\
 \mathsf{b\_le}&:=\mathsf{nat\_rec}_{\mathsf{nat}\to\mathsf{bool}}\ q_0\ d,
 &\mathsf{b\_eq}&:=\mathsf{nat\_rec}_{\mathsf{nat}\to\mathsf{bool}}\ e_0\ d,\\
 \mathsf{b\_lt}\ n\ m&:=\mathsf{b\_le}\ (S n)\ m,
 &\mathsf{b\_ge}\ n\ m&:=\mathsf{b\_le}\ m\ n,\\
 \mathsf{b\_gt}\ n\ m&:=\mathsf{b\_lt}\ m\ n.
\end{aligned}
\]
The equations are exactly zero/zero and zero/successor true for b\_le,
successor/zero false, and successor/successor the previous comparison;
b\_eq differs only at zero/successor, which is false
\cite[Definitions~12--14,17, pp.~21--22]{OostdijkGeuvers2000ProofByComputationPreprint}.
\end{hbdefinition}

\begin{hblemma}{thm:p12-har-prcompiler-boolean}{proved=annex:p12.har.prcompiler-boolean}{Boolean typing, logical operations and the numerical-code bridge}
The displayed programs have their indicated data types. For all a,b:bool,
istrue preserves not, and, or and implication in both directions:
\(\mathsf{istrue}(\mathsf{b\_imp}\ a\ b)
 \leftrightarrow(\mathsf{istrue}\ a\to\mathsf{istrue}\ b)\),
and analogously for the other three connectives. Also
\[
 (\mathsf{code}\ b=O)\vee(\mathsf{code}\ b=\mathsf{one}),\quad
 \mathsf{istrue}\ b\leftrightarrow\mathsf{code}\ b=\mathsf{one},\quad
 \neg\mathsf{istrue}\ b\leftrightarrow\mathsf{code}\ b=O.
\]
The last bridge is an actual computed conversion between two data
representations, not their silent identification.
\end{hblemma}

\begin{hblemma}{thm:p12-har-prcompiler-atoms}{proved=annex:p12.har.prcompiler-atoms}{All five atomic adequacy statements}
For every n,m:nat, istrue of b\_eq, b\_le, b\_lt, b\_ge and b\_gt
is respectively equivalent to \(n=m,n\le m,n<m,m\le n,m<n\).
The equalities and orders are the actual historical library relations.
\end{hblemma}

\begin{hbdefinition}{def:p12-har-prcompiler-bounds}{Both exclusive bounded-quantifier programs}
Let \(C:=(\mathsf{nat}\to\mathsf{bool})\to\mathsf{bool}\). Define
\[
\begin{aligned}
 \mathsf{b\_all}&:=\mathsf{nat\_rec}_C
    (\lambda f:\mathsf{nat}\to\mathsf{bool}.\mathsf{true})
    (\lambda m:\mathsf{nat}.\lambda r:C.
     \lambda f:\mathsf{nat}\to\mathsf{bool}.
                         \mathsf{b\_and}\ (f m)\ (r f)),\\
 \mathsf{b\_ex}&:=\mathsf{nat\_rec}_C
    (\lambda f:\mathsf{nat}\to\mathsf{bool}.\mathsf{false})
    (\lambda m:\mathsf{nat}.\lambda r:C.
     \lambda f:\mathsf{nat}\to\mathsf{bool}.
                         \mathsf{b\_or}\ (f m)\ (r f)).
\end{aligned}
\]
Thus at O the results are true/false, and at S m the endpoint f m is
combined with the old prefix. These scan \(0\le x<b\), not \(x\le b\),
and their evaluation order is the source's descending conjunction or
disjunction \cite[Definitions~15--16, p.~22]{OostdijkGeuvers2000ProofByComputationPreprint}.
\end{hbdefinition}

\begin{hbtheorem}{thm:p12-har-prcompiler-bounds}{proved=annex:p12.har.prcompiler-bounds}{Both bounded quantifier equivalences, with zero and endpoint cases}
For every b:nat and every f:nat-to-bool,
\[
\begin{aligned}
 \mathsf{istrue}(\mathsf{b\_all}\ b\ f)&\leftrightarrow
                 \Pi x:\mathsf{nat}.(x<b)\to\mathsf{istrue}(f x),\\
 \mathsf{istrue}(\mathsf{b\_ex}\ b\ f)&\leftrightarrow
                 \exists x:\mathsf{nat}.(x<b)\wedge\mathsf{istrue}(f x).
\end{aligned}
\]
Both reverse implications are constructive; the existential proof is
used only to prove a proposition, never to construct the boolean datum.
\end{hbtheorem}

\begin{hbdefinition}{def:p12-har-prcompiler-form}{The exact eleven-constructor internal formula datatype}
The source declaration has five constructors
\(\mathsf{f\_lt},\mathsf{f\_le},\mathsf{f\_eq},\mathsf{f\_ge},\mathsf{f\_gt}:
\mathsf{nat}\to\mathsf{nat}\to\mathsf{form}\),
one \(\mathsf{f\_not}:\mathsf{form}\to\mathsf{form}\),
three \(\mathsf{f\_and},\mathsf{f\_or},\mathsf{f\_imp}:
\mathsf{form}\to\mathsf{form}\to\mathsf{form}\), and
\[
 \mathsf{f\_all},\mathsf{f\_ex}:
        \mathsf{nat}\to(\mathsf{nat}\to\mathsf{form})\to\mathsf{form}.
\]
In particular the last argument is a typed family, not an untyped string
or an assumed finite list of formulas
\cite[Definition~7 and generated induction, pp.~18--19]{OostdijkGeuvers2000ProofByComputationPreprint}.
For a fixed target C, its generated fold takes five branches
\(\mathsf{nat}\to\mathsf{nat}\to C\), one branch
\(\mathsf{form}\to C\to C\), three branches
\(\mathsf{form}\to C\to\mathsf{form}\to C\to C\), and two branches
\(\mathsf{nat}\to(\mathsf{nat}\to\mathsf{form})\to
(\mathsf{nat}\to C)\to C\).
At a quantifier constructor the supplied recursive result is exactly
\(\lambda x:\mathsf{nat}.\mathsf{fold}_C(h x)\).
Here \(\mathsf{fold}_C\) means the source-generated primitive
\(\mathsf{Rec}_{\mathsf{form}}^C\) applied to these eleven branches,
in constructor order; it is not an additional postulated evaluator.
The generated Prop induction similarly supplies
\(\Pi x:\mathsf{nat}.P(h x)\) at a family constructor; no unproved
parametricity principle is substituted for this induction.
\end{hbdefinition}

\begin{hbdefinition}{def:p12-har-prcompiler-interpreters}{Every branch of the two internal interpreters}
Define \(\mathcal P:\mathsf{form}\to\mathsf{Prop}\) and
\(\mathsf{comp}:\mathsf{form}\to\mathsf{bool}\) by the preceding
fixed-target folds, at Prop and bool respectively. Their actual branches are
\[
\begin{array}{c|c|c}
 \text{constructor}&\mathcal P&\mathsf{comp}\\\hline
 \mathsf{f\_lt}\ n\ m&n<m&\mathsf{b\_lt}\ n\ m\\
 \mathsf{f\_le}\ n\ m&n\le m&\mathsf{b\_le}\ n\ m\\
 \mathsf{f\_eq}\ n\ m&n=m&\mathsf{b\_eq}\ n\ m\\
 \mathsf{f\_ge}\ n\ m&m\le n&\mathsf{b\_ge}\ n\ m\\
 \mathsf{f\_gt}\ n\ m&m<n&\mathsf{b\_gt}\ n\ m\\
 \mathsf{f\_not}\ p&\neg\mathcal P p&\mathsf{b\_not}(\mathsf{comp}\ p)\\
 \mathsf{f\_and}\ p\ q&\mathcal P p\wedge\mathcal P q&
                   \mathsf{b\_and}(\mathsf{comp}\ p)(\mathsf{comp}\ q)\\
 \mathsf{f\_or}\ p\ q&\mathcal P p\vee\mathcal P q&
                   \mathsf{b\_or}(\mathsf{comp}\ p)(\mathsf{comp}\ q)\\
 \mathsf{f\_imp}\ p\ q&\mathcal P p\to\mathcal P q&
                   \mathsf{b\_imp}(\mathsf{comp}\ p)(\mathsf{comp}\ q)\\
 \mathsf{f\_all}\ b\ h&\Pi x.(x<b)\to\mathcal P(h x)&
                   \mathsf{b\_all}\ b\ (\lambda x.\mathsf{comp}(h x))\\
 \mathsf{f\_ex}\ b\ h&\exists x.(x<b)\wedge\mathcal P(h x)&
                   \mathsf{b\_ex}\ b\ (\lambda x.\mathsf{comp}(h x))
\end{array}
\]
This table gives code templates: for example, the all branches are
\(\lambda b:\mathsf{nat}.\lambda h:\mathsf{nat}\to\mathsf{form}.
\lambda r:\mathsf{nat}\to\mathsf{Prop}.\Pi x.(x<b)\to r x\)
and the corresponding bool branch
\(\lambda b.\lambda h.\lambda r.\mathsf{b\_all}\ b\ r\),
with the fold interface fixing every omitted annotation. Each other row
is the same fully typed abstraction over its constructor fields and
recursive results; a recursive result is used, not a second hidden call
to an unspecified evaluator. In particular the complete comp term is
\[
\begin{aligned}
 \mathsf{comp}:={}&\mathsf{Rec}_{\mathsf{form}}^{\mathsf{bool}}
  \ \mathsf{b\_lt}\ \mathsf{b\_le}\ \mathsf{b\_eq}\ \mathsf{b\_ge}\ \mathsf{b\_gt}\\
 & (\lambda p:\mathsf{form}.\lambda u:\mathsf{bool}.\mathsf{b\_not}\ u)\\
 & (\lambda p:\mathsf{form}.\lambda u:\mathsf{bool}.
     \lambda q:\mathsf{form}.\lambda v:\mathsf{bool}.\mathsf{b\_and}\ u\ v)\\
 & (\lambda p:\mathsf{form}.\lambda u:\mathsf{bool}.
     \lambda q:\mathsf{form}.\lambda v:\mathsf{bool}.\mathsf{b\_or}\ u\ v)\\
 & (\lambda p:\mathsf{form}.\lambda u:\mathsf{bool}.
     \lambda q:\mathsf{form}.\lambda v:\mathsf{bool}.\mathsf{b\_imp}\ u\ v)\\
 & (\lambda b:\mathsf{nat}.\lambda h:\mathsf{nat}\to\mathsf{form}.
     \lambda r:\mathsf{nat}\to\mathsf{bool}.\mathsf{b\_all}\ b\ r)\\
 & (\lambda b:\mathsf{nat}.\lambda h:\mathsf{nat}\to\mathsf{form}.
     \lambda r:\mathsf{nat}\to\mathsf{bool}.\mathsf{b\_ex}\ b\ r).
\end{aligned}
\]
The Prop-target term uses the eleven analogous, fully annotated
branches specified in its typing proof. This primitive-recursion form
reconstructs the source's structural Fixpoint definitions, rather than
claiming a new axiom for comp. The primary paper writes these interpreters
as \(\llbracket-\rrbracket\) and
\(\llparenthesis-\rrparenthesis\), respectively
\cite[Definitions~10--11, pp.~20--21]{OostdijkGeuvers2000ProofByComputationPreprint}.
\end{hbdefinition}

\begin{hblemma}{thm:p12-har-prcompiler-form-typing}{proved=annex:p12.har.prcompiler-form-typing}{All interpreter branches and generated induction are well typed}
The two fold terms have exactly the stated types; every displayed
equation is their constructor computation. The correctness motive
\(P(p):=(\mathcal P p\leftrightarrow\mathsf{istrue}(\mathsf{comp}\ p))\)
is admitted by the generated form induction, including its two
\(\Pi x.P(h x)\) premises.
\end{hblemma}

\begin{hbtheorem}{thm:p12-har-prcompiler-ok}{proved=annex:p12.har.prcompiler-ok}{The single internal universal adequacy proof}
There is a proof term, constructed by form induction,
\[
 \mathsf{ok}:\Pi p:\mathsf{form}.
            \mathcal P p\leftrightarrow\mathsf{istrue}(\mathsf{comp}\ p).
\]
It holds in every valid open parameter context, for arbitrary typed nat
fields and arbitrary typed quantifier families. It does not require
closed inputs, DN, induction on a formula string, or proof irrelevance.
It is the actual internal reusable correctness component of the source
promise, not a separately assumed theorem for each input formula.
\end{hbtheorem}

\begin{hbtheorem}{thm:p12-har-prcompiler-extraction}{proved=annex:p12.har.prcompiler-extraction}{Predicate instances, numerical characteristic codes and closed proof objects}
Put \(F_0:=\mathsf{form}\) and \(F_{k+1}:=\mathsf{nat}\to F_k\).
For any well-typed family \(p:F_k\),
including \(k=0\) with no arguments, define the actual program
\[
 K_p:=\lambda\vec x:\mathsf{nat}.
           \mathsf{code}(\mathsf{comp}(p\vec x)).
\]
It has type \(T_k\), zero/one range, and for every open tuple
\(\mathcal P(p\vec x)\leftrightarrow K_p\vec x=\mathsf{one}\),
as well as \(\neg\mathcal P(p\vec x)\leftrightarrow K_p\vec x=O\).
If a particular closed \(\mathsf{comp}\ p\) contracts to true, then
\((\mathsf{ok}\ p)_{\leftarrow}\ \mathsf{trivial}\) is a proof of
\(\mathcal P p\) by conversion. No claim that every arbitrary closed
axiomatic nat term computes is needed for this implication.
\end{hbtheorem}

\begin{hbtheorem}{thm:p12-har-prcompiler-predicate}{proved=annex:p12.har.prcompiler-predicate}{Every prescribed primitive-recursive predicate has a characteristic program}
Let e be any ranked primitive-recursive prescription of arity k and
define its metalevel numerical predicate
\(R_e(\vec a):\!\iff\langle e\rangle(\vec a)=1\).
Construct, without a supplied per-input correctness theorem,
\[
 p_e:=\lambda\vec x:\mathsf{nat}.
          \mathsf{f\_eq}\ ((\mathcal C e)\vec x)\ \mathsf{one},
 \qquad K_e:=K_{p_e}.
\]
For every numerical tuple, the finite program computation yields
\(K_e\overline{\vec a}\longrightarrow^*\mathsf{one}\) if
\(R_e(\vec a)\), and yields O otherwise. For arbitrary typed open
tuples its internally proved adequacy is
\[
 ((\mathcal C e)\vec x=\mathsf{one})
                 \leftrightarrow K_e\vec x=\mathsf{one},
 \quad
 \neg((\mathcal C e)\vec x=\mathsf{one})
                 \leftrightarrow K_e\vec x=O.
\]
In particular any predicate specified by a finite prescription for
its zero/one characteristic function is included. The convention
\(R_e=\{\vec a:\langle e\rangle\vec a=1\}\) also defines a
predicate when e's values are not confined to zero/one; the resulting
\emph{new} \(K_e\) still has zero/one range. Neither membership testing
for an arbitrary host function nor an internal inverse on Prop is
asserted by this uniform external template plus internal comp program.
\end{hbtheorem}

\begin{hbexample}{def:p12-har-prcompiler-source-prime}{The primary paper's actual division and Prime formula inputs}
At the source unary arithmetic library, Example~9 supplies the two
\emph{formula-valued} programs
\[
\begin{aligned}
 \mathsf{f\_Divides}&:=\lambda n,m:\mathsf{nat}.
  \mathsf{f\_ex}\ (S m)
       (\lambda k:\mathsf{nat}.\mathsf{f\_eq}\ (\mathsf{times}\ k\ n)\ m),\\
 \mathsf{f\_Prime}&:=\lambda n:\mathsf{nat}.
  \mathsf{f\_and}\ (\mathsf{f\_gt}\ n\ \mathsf{one})\\
 &\hspace{35mm}(\mathsf{f\_all}\ n\ (\lambda d:\mathsf{nat}.
    \mathsf{f\_imp}\ (\mathsf{f\_Divides}\ d\ n)
                         (\mathsf{f\_eq}\ d\ \mathsf{one}))).
\end{aligned}
\]
Their types are nat-to-nat-to-form and nat-to-form, respectively.
Expanding \(\mathcal P\) gives, on these same terms,
\(\exists k.(k<S m)\wedge(kn=m)\) and
\((n>\mathsf{one})\wedge\Pi d.(d<n)\to
 ((\exists k.(k<S n)\wedge(kd=n))\to d=\mathsf{one})\).
These are the exact bounded division and Prime inputs of this
implementation, not a new axiom naming ordinary primality. Applying
\cref{thm:p12-har-prcompiler-extraction} supplies their actual
characteristic programs and open universal adequacy proofs without
a second per-input proof search. The bounded coefficient convention,
strict d<n fold and formula's conjunction order are retained
\cite[Examples~6,9, pp.~17,19]{OostdijkGeuvers2000ProofByComputationPreprint}.
No definitional identity with the chapter's separate unbounded div/
proper-divisor prime declaration or the earlier direct KPrime program
is inferred merely from the name. The primary paper reports its
Prime 61 computation and unary/nonoptimized costs historically; this
entry does not claim a newly executed kernel computation or a finite
61 trace for this different generic search
\cite[pp.~24--25]{OostdijkGeuvers2000ProofByComputationPreprint}.
\end{hbexample}

\begin{proofboundary}{def:p12-har-prcompiler-boundary}{What the generic construction supplies and what it does not}
Finite primitive-recursive function prescriptions now have an actual
external syntax compiler and a universal numerical representation proof.
The source's internal formula program and its universal adequacy proof
are separately supplied, with all eleven constructor branches. One can
put compiled primitive-recursive terms in the atomic nat fields and
abstract over free data parameters; this does not require a function-AST
datatype or a choice principle. Arbitrary Prop-to-form reification is not
claimed: it is expressly an unimplemented external interface in the
primary source. The paper's further extension by an atomic R requires
an actual boolean function and a proof
\(\Pi\vec x.R\vec x\leftrightarrow\mathsf{istrue}(f_R\vec x)\);
adding that constructor alone cannot manufacture this proof.
The direct KPrime program already supplies a different concrete
characteristic construction; no identity of its search algorithm with
the generic compiler or the printed Prime 61 implementation is asserted.
Pocklington/Fermat and large certificates, differentiation/continuity,
arithmetic coding, and global inductive-extension/C001 metatheory retain
their separate obligations.
\end{proofboundary}


\subsection{The chapter's dated assistant comparison}

The following retains the 2001 source's comparison, not current products
or a recommendation. The authors explicitly describe the ratings as
intuitive and not to be taken too seriously. Here \(\mathrm{n.s.}\)
means non-standard proof objects/scripts. The historical PP codes are
\(R_1\): a separately chosen reduction notion; \(R_2\): structurally
well-founded recursion; \(R_3\): arbitrary provable equality;
\(R_4\): \(\varepsilon_0\)-recursion. The surrounding text adds
normalization/proof conditions to the chosen reduction mechanism;
the bare code \(R_1\) does not license every unverified rewrite
\cite[pp.~1224--1228]{BarendregtGeuvers2001}.

\begin{longtable}{@{}lcccccccc@{}}
\caption{The source's historical assistant table and ratings.}
\label{tab:p12-har-assistants}\\
\toprule
Assistant & Objects & Rating & PP & Logic & Dep. & Ind. & Auto. & Users\\
\midrule
\endfirsthead
\caption[]{The source's historical assistant table and ratings (continued)}\\
\toprule
Assistant & Objects & Rating & PP & Logic & Dep. & Ind. & Auto. & Users\\
\midrule
\endhead
\midrule\multicolumn{9}{r}{\emph{continued on next page}}\\
\endfoot
\endlastfoot
Agda & yes & +++ & \(\beta\delta\iota R_1\) & int./pred. & yes & yes & none\textsuperscript{8} & -\\
Coq & yes & ++ & \(\beta\delta\iota R_2\) & int. & yes & yes & + & ++\\
Lego & yes & ++ & \(\beta\delta\iota R_1\) & int. & yes & yes & + & +\\
Nuprl & n.s. & ++ & \(\beta\delta\iota R_3\) & int. & yes & yes & + & ++\\
HOL & n.s. & ++ & none & cl. & no & yes & ++ & ++\\
Isabelle & n.s. & ++ & none & t.b.s. & no & no & ++ & ++\\
Mizar & none & + & none & cl. & yes & no & + & ++\\
ACL2 & none & + & \(R_4\) & pra & no & yes\textsuperscript{9} & +++ & +++\\
PVS & none & - & none & cl. & no & no & +++ & +++\\
\bottomrule
\end{longtable}

The source abbreviations are intuitionistic (int.), predicative (pred.),
classical (cl.), object logic chosen by the user (t.b.s.), and its
primitive-recursive arithmetic description (pra). Source note~8 records
Agda's limited higher-order-unification automation; note~9 records ACL2's
coding of induction in its basic inductive type. These columns classify
the source's historical presentations, not contemporary kernel features.

Coq/Lego/Agda are the source's stored-object group. Nuprl's original
extensional judgment is distinguished from its expanded evidence that
records equality proofs. HOL/Isabelle scripts are discussed as replayable
non-standard evidence, with computation handled by explicit rewriting,
not silently by the source's conversion principle. Isabelle's meta-logic
and its internally declared object logics are separate. Mizar, ACL2 and
PVS are the source's no-stored-object group; its PVS discussion includes
predicate-subtype obligations. No row implies current portable-object
support or absence merely from a product name.

Finally the chapter contrasts its pragmatic verification approach with a
small-checker fundamental approach, while keeping logical strength and
computational strength as separate parameters. Its proposed ``universal''
assistant is foundationally broad, not a complete algorithm for all
mathematical truth. Its anticipated research/refereeing/teaching and
program-extraction benefits remain historical projections. Proof-dependent
quotient/subset/partial-function representations require their actual
typing and checking regime, and classical DN in the prime development
does not automatically yield a constructive witness-extraction algorithm
\cite[pp.~1228--1230]{BarendregtGeuvers2001}.

\showcalculus{p12-har-types}
